\documentclass[
  11pt,                % Schriftgröße
  a4paper,             % Papierformat
  headsepline,         % Trennlinie unter der Kopfzeile
  footsepline,         % Trennlinie über der Fußzeile
  twoside,             % Zweiseitig
  openany,             % Neue Kapitel auf beliebigen Seiten
  parskip=full,        % Absatzabstand statt Einzug
  numbers=noenddot,    % Kapitelnummern ohne abschließenden Punkt
  listof=totoc,        % Alle Verzeichnisse in Inhaltsverzeichniss
  headings=optiontoheadandtoc % Ermöglicht head=... und toc=... bei \section
]{scrbook}             % KOMA-Script-Klasse (empfohlen für Thesen)

% ========================================
% PAKETE & EINSTELLUNGEN
% ========================================
\usepackage{settings} % Eigenes Paket mit weiteren Einstellungen

% Spracheinstellungen (Deutsch)
\usepackage[ngerman]{babel}
\usepackage[babel,german=quotes]{csquotes}

% Literaturverwaltung: BibLaTeX mit Biber
\usepackage[
  backend=biber,
  style=ieee,            % Format des Literaturverzeichnisses bleibt IEEE
  citestyle=numeric-comp, % <-- fasst Zitate zu Bereichen zusammen
  sorting=none,
  sortcites=true,
  maxcitenames=2,
  maxbibnames=99,
  isbn=false,
  url=false,
  doi=false,
  eprint=false
]{biblatex}
\addbibresource{bibliography.bib}

% ========================
% DOKUMENTSTART
% ========================
\begin{document}

% \listoftodos % Verzeichnis anzeigen

% ========================
% VORTEIL (römische Seitenzahlen, einfacher Zeilenabstand)
% ========================
\frontmatter          % römische Seitenzahlen, unnummerierte Kapitel

% 00-titlepage.tex
\thispagestyle{empty}
\pdfbookmark[0]{Titelblatt}{title}

\begin{titlepage}

  % Logos oben
  \vspace*{-2.5cm}
  \begin{center}
    \centerlogo[4.2cm]{png/hka_logo}%
    \ifdefstring{\CompanyName}{NoCompanyInvolved}{}{
      \hspace{0.0cm}%
      \centerlogo[4.2cm]{example-image-a}
      % \centerlogo[3cm]{png/ict_logo}%
    }%
  \end{center}

  % Hauptblock zentriert
  \begin{center}
    \vspace{0.1cm}
    \LARGE \textbf{\University}\\
    \vspace{0.3cm}
    \Large -- \Faculty --

    \vspace{1cm}

    \huge \textbf{\Title}\\
    % \vspace{0.4cm}
    \ifdefstring{\Subtitle}{NoSubtitle}{}{
      \LARGE \Subtitle
    }

    \vspace{1cm}

    \Large Zur Erlangung des akademischen Grades \\
    \vspace{0.3cm}
    \LARGE \textbf{\Degree} \\
    (\DegreeShort) \\
    \vspace{0.3cm}
    \Large vorgelegte \\
    \vspace{0.3cm}
    \LARGE \textbf{\DocName} \\
    \vspace{0.3cm}
    \Large von \\
    \vspace{0.3cm}
    \Large \textbf{\Name} (\Id)\\
    \Large geb. am \DateOfBirth \\
    in \PlaceOfBirth \\
  \end{center}

  % Ein einziges vfill, um die Tabelle nach unten zu drücken
  \vfill

  % Tabelle unten (zentriert oder flushleft, wie gewünscht)
  \begin{center} % oder \noindent\begin{tabular}{...} ohne center
    \begin{tabular}{lll}
      Tag der Abgabe      & : & \DeliveryDate \\
      Erstprüfer          & : & \Supervisor \\
      Zweitprüfer         & : & \CoSupervisor\\
      \ifdefstring{\CompanyName}{NoCompanyInvolved}{}{Betreuer \CompanyName & : & \CompSupervisor\\}
    \end{tabular}
  \end{center}

\end{titlepage}
\cleardoubleoddemptypage
% 00-declaration.tex
\addchap{Eigenständigkeitserklärung}

Ich versichere hiermit wahrheitsgemäß, dass ich die vorliegende \DocName{} mit dem Titel
\begin{quote}
  \textbf{\Title} \\
  \textit{\Subtitle}
\end{quote}
selbstständig verfasst und keine anderen als die angegebenen Quellen und Hilfsmittel benutzt habe. Alle Stellen, die wörtlich oder sinngemäß aus veröffentlichten oder nicht veröffentlichten Schriften entnommen sind, habe ich als solche kenntlich gemacht.

Die Arbeit ist in gleicher oder ähnlicher Form noch bei keiner anderen Prüfungsbehörde eingereicht und auch nicht veröffentlicht worden.

\vspace{3cm}

\noindent
\begin{tabular}{@{}p{10cm}@{}}
  \hrulefill \\
  \Location,\hfill den \DeliveryDate, \hfill \Name
\end{tabular}

\clearpage
% !TeX root = ../main.tex
% 00-preface.tex — Kurzfassung / Abstract

\addchap{Abstract}

\vspace{0.5cm}
\noindent\textbf{Keywords:} 



\addchap{Kurzfassung}

\vspace{0.5cm}
\noindent\textbf{Schlagwörter:} 


% End of File

% !TeX root = main.tex
% nomenclature.tex

\nomenclature{$F_\mathrm{x}$}{Längskraft \nomunit{\unit{\newton}}}
\nomenclature{$F_\mathrm{y}$}{Querkraft \nomunit{\unit{\newton}}}
\nomenclature{$F_\mathrm{z}$}{Radlast \nomunit{\unit{\newton}}}

\tableofcontents
\listoffigures
\listoftables
% acronyms.tex
\addchap{Abkürzungsverzeichnis}

\begin{acronym}[IBVS]
	\acro{IBVS}{Image Based Visual Servoing}
	\acro{AI}{Artificial Intelligence}

  \acro{svm}[SVM]{support vector machine}
  \acrodefplural{svm}[SVMs]{support vector machines}
\end{acronym}
\printnomenclature

% ========================
% HAUPTTEIL (arabische Seitenzahlen, 1,5-zeilig)
% ========================
\mainmatter           % arabische Seitenzahlen, nummerierte Kapitel

% !TeX root = ../main.tex
% 02-introduction.tex
\chapter{Einleitung}
\label{chap:introduction}


\section{Motivation und gesellschaftlicher Kontext}
\label{sec:motivation}

\lipsum[1-3]


\section{Problemstellung}
\label{sec:problem}

\lipsum[3-9]


\section{Zielsetzung und Forschungsfragen}
\label{sec:goal_research}
\lipsum[10]

Hieraus leitet sich die folgende übergeordnete Forschungsfrage ab:

\begin{quote}
    \textit{\lipsum[1][11]}
\end{quote}

Zur strukturierten Beantwortung werden vier \acp{TF} definiert, deren thematische Einordnung sowie der jeweilige Bezug zu den \aclp{AP}n~(\acs{AP}) und den bearbeitenden Kapiteln in Tabelle~\ref{tab:forschungsfragen} zusammengefasst sind.

% !TeX root = ../main.tex

\begin{table}[htbp]
    \centering
    \scriptsize
    \caption{Teilforschungsfragen, Arbeitspakete und Einordnung in die Kapitelstruktur}
    \label{tab:forschungsfragen}
    \begin{threeparttable}
        \renewcommand{\arraystretch}{1.6}
        \begin{tabularx}{\textwidth}{c X c c}
            \toprule
            \textbf{Nr.} & \textbf{Forschungsfrage} & \textbf{AP} & \textbf{Kap.} \\
            \midrule

            \textbf{\acs{TF}\,1} & \lipsum[1][12] & 2,\;3 & 3,\;4 \\

            \textbf{\acs{TF}\,2} & \lipsum[1][13] & 1,\;4 & 3,\;5 \\

            \textbf{\acs{TF}\,3} & \lipsum[1][14] & 5,\;6 & 3,\;6 \\

            \textbf{\acs{TF}\,4} & \lipsum[1][15] & 5,\;6 & 7 \\

            \bottomrule
        \end{tabularx}
    \end{threeparttable}
\end{table}

Die abschließende Synthese aller Teilantworten sowie die Beantwortung der übergeordneten Forschungsfrage erfolgt in Kapitel~\ref{chap:conclusion_outlook}.


% End of File
% !TeX root = ../main.tex
% 02-related-work.tex
\chapter{Stand der Technik und Grundlagen}
\label{chap:fundamentals}


\section[head={Sek. 1}]{Sektion die Erste}
\label{sec:section1}

\lipsum[1]


\subsection{Sub-Sektion}
\label{subsec:subsection11}

  \lipsum[2]


\subsection{Sub-Sektion}
\label{subsec:subsection12}

  \lipsum[3]


\subsubsection{Sub-Sub-Sektion}
\label{subsubsec:subsubsection121}

  \lipsum[4]


\subsubsection{Sub-Sub-Sektion}
\label{subsubsec:subsubsection122}

  \lipsum[5]


\section[head={Sek. 2}]{Sektion die Zweite}
\label{sec:section1}

\lipsum[6]



% End of File
% !TeX root = ../main.tex
% 03-methodology.tex
\chapter{Methodik}
\label{chap:methodology}

\lipsum[1]

\begin{figure}[htbp]
    \centering

    %% ── Zeile 1: Plots (vertikal mittig zueinander) ──────────────────
    \begin{minipage}[c]{0.49\textwidth}
      \centering
      \noindent% \documentclass[tikz, border=10pt]{standalone}
% \usepackage[T1]{fontenc}
% \usepackage[utf8]{inputenc}
% \usepackage{tikz}
% \usetikzlibrary{arrows.meta}

\newcommand{\streckeScale}{0.49}

%% ─── Stile ──────────────────────────────────────────────────────────────────
\tikzset{
  road/.style={
    line width=9pt, color=gray!20,
    line cap=round, line join=round
  },
  cline/.style={
    line width=3pt, line cap=round, line join=round, opacity=0.80
  },
  dirarrw/.style={
    -{Stealth[length=4pt, width=2.5pt]},
    line width=0.75pt, color=gray!20
  },
  startmark/.style={
    circle, fill=clrBorder, minimum size=6pt, inner sep=0pt
  },
  endmark/.style={
    circle, fill=white, draw=clrBorder, line width=1.2pt,
    minimum size=6pt, inner sep=0pt
  },
}

% \begin{document}
\noindent\begin{tikzpicture}[scale=\streckeScale]

%%
%% Streckengeometrie (Vogelansicht) – alle Bögen R = 1, 90°-Sweeps:
%%
%%  (0,0) →→→ (4,0)                     Seg 1–3  Gerade
%%  Seg  4  R-Kurve: (4,0)  → (5,−1)  ↓
%%  Seg  5  L-Kurve: (5,−1) → (6,−2)  →
%%  Seg  6  gerade:  (6,−2) → (7,−2)  →
%%  Seg  7  L-Kurve: (7,−2) → (8,−1)  ↑
%%  Seg  8  gerade:  (8,−1) → (8, 0.5)  ↑
%%  Seg  9  L-Kurve: (8, 1) → (7, 1)  ←
%%  Seg 10  R-Kurve: (7, 2) → (6, 2)  ↑
%%  Seg 11  gerade:  (6, 3) → (6, 3)  ↑
%%  Seg 12  R-Kurve: (6,4.5)→ (7, 4)  →
%%

%% ── Straßenoberfläche (grauer Hintergrundstreifen) ───────────────────────
\draw[road]
  (0,0) -- (4,0)
  arc[start angle= 90, end angle=  0, radius=1]   % Seg  4
  arc[start angle=180, end angle=270, radius=1]   % Seg  5
  -- (7,-2)                                        % Seg  6
  arc[start angle=270, end angle=360, radius=1]   % Seg  7
  -- (8, 0.5)                                        % Seg  8
  arc[start angle=  0, end angle= 90, radius=1]   % Seg  9
  arc[start angle=270, end angle=180, radius=1]   % Seg 10
  -- (6, 4)                                      % Seg 11
  arc[start angle=180, end angle= 90, radius=1];  % Seg 12

%% ── Manöverabschnitte (farbkodiert) ─────────────────────────────────────

%  1 · Beschleunigung (gerade)
\draw[cline, color=clrAccel]
  (0,0) -- (1.5,0);

%  2 · Konstante Gerade
\draw[cline, color=clrGerade]
  (1.5,0) -- (3.0,0);

%  3 · Bremsung (gerade)
\draw[cline, color=clrBrake]
  (3.0,0) -- (4.0,0);

%  4 · Konstante Kurve rechts
%      Start (4,0) @ 90°, Mitte (4,−1), Ende (5,−1) ↓
\draw[cline, color=clrCurveR]
  (4,0) arc[start angle=90, end angle=0, radius=1];

%  5 · Konstante Kurve links
%      Start (5,−1) @ 180°, Mitte (6,−1), Ende (6,−2) →
\draw[cline, color=clrCurveL]
  (5,-1) arc[start angle=180, end angle=270, radius=1];

%  6 · Bremsung (gerade)
\draw[cline, color=clrBrake]
  (6,-2) -- (7,-2);

%  7 · Beschleunigte Kurve links
%      Start (7,−2) @ 270°, Mitte (7,−1), Ende (8,−1) ↑
\draw[cline, color=clrCurveAccelL]
  (7,-2) arc[start angle=270, end angle=360, radius=1];

%  8 · Konstante Gerade
\draw[cline, color=clrGerade]
  (8,-1) -- (8,0.5);

%  9 · Gebremste Kurve links
%      Start (8,0.5) @ 0°, Mitte (7,1), Ende (7,1.5) ←
\draw[cline, color=clrCurveBrakeL]
  (8,0.5) arc[start angle=0, end angle=90, radius=1];

% 10 · Beschleunigte Kurve rechts
%      Start (7,1.5) @ 270°, Mitte (7,2), Ende (6,2.5) ↑
\draw[cline, color=clrCurveAccelR]
  (7,1.5) arc[start angle=270, end angle=180, radius=1];

% 11 · Konstante Gerade
\draw[cline, color=clrGerade]
  (6,2.5) -- (6,4);

% 12 · Gebremste Kurve rechts
%      Start (6,4) @ 180°, Mitte (7,4.5), Ende (7,5) →
\draw[cline, color=clrCurveBrakeR]
  (6,4) arc[start angle=180, end angle=90, radius=1];

%% ── Segment-Nummern (außerhalb des Fahrstreifens) ─────────────────────────
\node[font=\scriptsize\bfseries, color=clrAccel]      at ( 0.75,  0.55) {1};
\node[font=\scriptsize\bfseries, color=clrGerade]     at ( 2.25,  0.55) {2};
\node[font=\scriptsize\bfseries, color=clrBrake]      at ( 3.50,  0.55) {3};
\node[font=\scriptsize\bfseries, color=clrCurveR]      at ( 5.10,  0.05) {4};  % außen R-Kurve
\node[font=\scriptsize\bfseries, color=clrCurveL]      at ( 4.90, -2.05) {5};  % außen L-Kurve
\node[font=\scriptsize\bfseries, color=clrBrake]      at ( 6.50, -2.55) {6};
\node[font=\scriptsize\bfseries, color=clrCurveAccelL] at ( 8.10, -2.05) {7};  % außen L-Kurve
\node[font=\scriptsize\bfseries, color=clrGerade]     at ( 8.55, -0.25) {8};
\node[font=\scriptsize\bfseries, color=clrCurveBrakeL] at ( 8.10,  1.55) {9};  % außen L-Kurve
\node[font=\scriptsize\bfseries, color=clrCurveAccelR] at ( 5.90,  1.45) {10}; % außen R-Kurve
\node[font=\scriptsize\bfseries, color=clrGerade]     at ( 5.45,  3.25) {11};
\node[font=\scriptsize\bfseries, color=clrCurveBrakeR] at ( 5.90,  5.05) {12}; % außen R-Kurve

%% ── Fahrtrichtungspfeile ─────────────────────────────────────────────────
\draw[dirarrw] (1.75, 0.00) -- (2.75, 0.00);  % Seg  1  →
\draw[dirarrw] (8.00,-0.75) -- (8.00, 0.25);  % Seg  8  ↑
\draw[dirarrw] (6.00, 2.75) -- (6.00, 3.75);  % Seg 11  ↑

%% ── Start / Ende ─────────────────────────────────────────────────────────
\node[startmark] at (0,0) {};
\node[font=\scriptsize\bfseries, color=clrBorder, anchor=east]
  at (-0.25, 0.00) {Start};

\node[endmark] at (7,5) {};
\node[font=\scriptsize\bfseries, color=clrBorder, anchor=west]
  at ( 7.25, 5) {Ende};

%% ── Legende ──────────────────────────────────────────────────────────────

% Zeile 1  (y = −2.85)
% \draw[cline, color=clrAccel]
%   (0.00,-2.85) -- (0.55,-2.85);
% \node[font=\small, anchor=west] at (0.70,-2.85) {Beschleunigung};

% \draw[cline, color=clrGerade]
%   (3.40,-2.85) -- (3.95,-2.85);
% \node[font=\small, anchor=west] at (4.10,-2.85) {Konst.\,Gerade};

% \draw[cline, color=clrBrake]
%   (6.50,-2.85) -- (7.05,-2.85);
% \node[font=\small, anchor=west] at (7.20,-2.85) {Bremsung};

% % Zeile 2  (y = −3.55)
% \draw[cline, color=clrCurve]
%   (0.00,-3.55) -- (0.55,-3.55);
% \node[font=\small, anchor=west] at (0.70,-3.55) {Konst.\,Kurve};

% \draw[cline, color=clrCurveAccel]
%   (3.40,-3.55) -- (3.95,-3.55);
% \node[font=\small, anchor=west] at (4.10,-3.55) {Beschl.\,Kurve};

% \draw[cline, color=clrCurveBrake]
%   (6.50,-3.55) -- (7.05,-3.55);
% \node[font=\small, anchor=west] at (7.20,-3.55) {Gebr.\,Kurve};

\end{tikzpicture}
% \end{document}
    \end{minipage}%
    \hfill
    \begin{minipage}[c]{0.49\textwidth}
      \centering
      %% Creator: Matplotlib, PGF backend
%%
%% To include the figure in your LaTeX document, write
%%   \input{<filename>.pgf}
%%
%% Make sure the required packages are loaded in your preamble
%%   \usepackage{pgf}
%%
%% Also ensure that all the required font packages are loaded; for instance,
%% the lmodern package is sometimes necessary when using math font.
%%   \usepackage{lmodern}
%%
%% Figures using additional raster images can only be included by \input if
%% they are in the same directory as the main LaTeX file. For loading figures
%% from other directories you can use the `import` package
%%   \usepackage{import}
%%
%% and then include the figures with
%%   \import{<path to file>}{<filename>.pgf}
%%
%% Matplotlib used the following preamble
%%   \def\mathdefault#1{#1}
%%   \everymath=\expandafter{\the\everymath\displaystyle}
%%   \IfFileExists{scrextend.sty}{
%%     \usepackage[fontsize=10.000000pt]{scrextend}
%%   }{
%%     \renewcommand{\normalsize}{\fontsize{10.000000}{12.000000}\selectfont}
%%     \normalsize
%%   }
%%   \usepackage{amsmath}
%%   \usepackage{siunitx}
%%   \makeatletter\@ifpackageloaded{underscore}{}{\usepackage[strings]{underscore}}\makeatother
%%
\begingroup%
\makeatletter%
\begin{pgfpicture}%
\pgfpathrectangle{\pgfpointorigin}{\pgfqpoint{3.000000in}{2.000000in}}%
\pgfusepath{use as bounding box, clip}%
\begin{pgfscope}%
\pgfsetbuttcap%
\pgfsetmiterjoin%
\definecolor{currentfill}{rgb}{1.000000,1.000000,1.000000}%
\pgfsetfillcolor{currentfill}%
\pgfsetlinewidth{0.000000pt}%
\definecolor{currentstroke}{rgb}{1.000000,1.000000,1.000000}%
\pgfsetstrokecolor{currentstroke}%
\pgfsetdash{}{0pt}%
\pgfpathmoveto{\pgfqpoint{0.000000in}{0.000000in}}%
\pgfpathlineto{\pgfqpoint{3.000000in}{0.000000in}}%
\pgfpathlineto{\pgfqpoint{3.000000in}{2.000000in}}%
\pgfpathlineto{\pgfqpoint{0.000000in}{2.000000in}}%
\pgfpathlineto{\pgfqpoint{0.000000in}{0.000000in}}%
\pgfpathclose%
\pgfusepath{fill}%
\end{pgfscope}%
\begin{pgfscope}%
\pgfsetbuttcap%
\pgfsetmiterjoin%
\definecolor{currentfill}{rgb}{1.000000,1.000000,1.000000}%
\pgfsetfillcolor{currentfill}%
\pgfsetlinewidth{0.000000pt}%
\definecolor{currentstroke}{rgb}{0.000000,0.000000,0.000000}%
\pgfsetstrokecolor{currentstroke}%
\pgfsetstrokeopacity{0.000000}%
\pgfsetdash{}{0pt}%
\pgfpathmoveto{\pgfqpoint{0.464911in}{0.456793in}}%
\pgfpathlineto{\pgfqpoint{2.442496in}{0.456793in}}%
\pgfpathlineto{\pgfqpoint{2.442496in}{1.910105in}}%
\pgfpathlineto{\pgfqpoint{0.464911in}{1.910105in}}%
\pgfpathlineto{\pgfqpoint{0.464911in}{0.456793in}}%
\pgfpathclose%
\pgfusepath{fill}%
\end{pgfscope}%
\begin{pgfscope}%
\pgfpathrectangle{\pgfqpoint{0.464911in}{0.456793in}}{\pgfqpoint{1.977584in}{1.453311in}}%
\pgfusepath{clip}%
\pgfsetbuttcap%
\pgfsetmiterjoin%
\definecolor{currentfill}{rgb}{0.121569,0.466667,0.705882}%
\pgfsetfillcolor{currentfill}%
\pgfsetfillopacity{0.200000}%
\pgfsetlinewidth{0.000000pt}%
\definecolor{currentstroke}{rgb}{0.121569,0.466667,0.705882}%
\pgfsetstrokecolor{currentstroke}%
\pgfsetstrokeopacity{0.200000}%
\pgfsetdash{}{0pt}%
\pgfpathmoveto{\pgfqpoint{0.520568in}{0.456793in}}%
\pgfpathlineto{\pgfqpoint{0.649379in}{0.456793in}}%
\pgfpathlineto{\pgfqpoint{0.649379in}{1.910105in}}%
\pgfpathlineto{\pgfqpoint{0.520568in}{1.910105in}}%
\pgfpathlineto{\pgfqpoint{0.520568in}{0.456793in}}%
\pgfpathclose%
\pgfusepath{fill}%
\end{pgfscope}%
\begin{pgfscope}%
\pgfpathrectangle{\pgfqpoint{0.464911in}{0.456793in}}{\pgfqpoint{1.977584in}{1.453311in}}%
\pgfusepath{clip}%
\pgfsetbuttcap%
\pgfsetmiterjoin%
\definecolor{currentfill}{rgb}{0.750000,0.750000,0.750000}%
\pgfsetfillcolor{currentfill}%
\pgfsetfillopacity{0.200000}%
\pgfsetlinewidth{0.000000pt}%
\definecolor{currentstroke}{rgb}{0.750000,0.750000,0.750000}%
\pgfsetstrokecolor{currentstroke}%
\pgfsetstrokeopacity{0.200000}%
\pgfsetdash{}{0pt}%
\pgfpathmoveto{\pgfqpoint{0.674881in}{0.456793in}}%
\pgfpathlineto{\pgfqpoint{0.737797in}{0.456793in}}%
\pgfpathlineto{\pgfqpoint{0.737797in}{1.910105in}}%
\pgfpathlineto{\pgfqpoint{0.674881in}{1.910105in}}%
\pgfpathlineto{\pgfqpoint{0.674881in}{0.456793in}}%
\pgfpathclose%
\pgfusepath{fill}%
\end{pgfscope}%
\begin{pgfscope}%
\pgfpathrectangle{\pgfqpoint{0.464911in}{0.456793in}}{\pgfqpoint{1.977584in}{1.453311in}}%
\pgfusepath{clip}%
\pgfsetbuttcap%
\pgfsetmiterjoin%
\definecolor{currentfill}{rgb}{1.000000,0.498039,0.054902}%
\pgfsetfillcolor{currentfill}%
\pgfsetfillopacity{0.200000}%
\pgfsetlinewidth{0.000000pt}%
\definecolor{currentstroke}{rgb}{1.000000,0.498039,0.054902}%
\pgfsetstrokecolor{currentstroke}%
\pgfsetstrokeopacity{0.200000}%
\pgfsetdash{}{0pt}%
\pgfpathmoveto{\pgfqpoint{0.740775in}{0.456793in}}%
\pgfpathlineto{\pgfqpoint{0.767952in}{0.456793in}}%
\pgfpathlineto{\pgfqpoint{0.767952in}{1.910105in}}%
\pgfpathlineto{\pgfqpoint{0.740775in}{1.910105in}}%
\pgfpathlineto{\pgfqpoint{0.740775in}{0.456793in}}%
\pgfpathclose%
\pgfusepath{fill}%
\end{pgfscope}%
\begin{pgfscope}%
\pgfpathrectangle{\pgfqpoint{0.464911in}{0.456793in}}{\pgfqpoint{1.977584in}{1.453311in}}%
\pgfusepath{clip}%
\pgfsetbuttcap%
\pgfsetmiterjoin%
\definecolor{currentfill}{rgb}{0.464314,0.323137,0.592941}%
\pgfsetfillcolor{currentfill}%
\pgfsetfillopacity{0.200000}%
\pgfsetlinewidth{0.000000pt}%
\definecolor{currentstroke}{rgb}{0.464314,0.323137,0.592941}%
\pgfsetstrokecolor{currentstroke}%
\pgfsetstrokeopacity{0.200000}%
\pgfsetdash{}{0pt}%
\pgfpathmoveto{\pgfqpoint{0.768138in}{0.456793in}}%
\pgfpathlineto{\pgfqpoint{0.805367in}{0.456793in}}%
\pgfpathlineto{\pgfqpoint{0.805367in}{1.910105in}}%
\pgfpathlineto{\pgfqpoint{0.768138in}{1.910105in}}%
\pgfpathlineto{\pgfqpoint{0.768138in}{0.456793in}}%
\pgfpathclose%
\pgfusepath{fill}%
\end{pgfscope}%
\begin{pgfscope}%
\pgfpathrectangle{\pgfqpoint{0.464911in}{0.456793in}}{\pgfqpoint{1.977584in}{1.453311in}}%
\pgfusepath{clip}%
\pgfsetbuttcap%
\pgfsetmiterjoin%
\definecolor{currentfill}{rgb}{0.138039,0.501961,0.138039}%
\pgfsetfillcolor{currentfill}%
\pgfsetfillopacity{0.200000}%
\pgfsetlinewidth{0.000000pt}%
\definecolor{currentstroke}{rgb}{0.138039,0.501961,0.138039}%
\pgfsetstrokecolor{currentstroke}%
\pgfsetstrokeopacity{0.200000}%
\pgfsetdash{}{0pt}%
\pgfpathmoveto{\pgfqpoint{0.805553in}{0.456793in}}%
\pgfpathlineto{\pgfqpoint{0.939762in}{0.456793in}}%
\pgfpathlineto{\pgfqpoint{0.939762in}{1.910105in}}%
\pgfpathlineto{\pgfqpoint{0.805553in}{1.910105in}}%
\pgfpathlineto{\pgfqpoint{0.805553in}{0.456793in}}%
\pgfpathclose%
\pgfusepath{fill}%
\end{pgfscope}%
\begin{pgfscope}%
\pgfpathrectangle{\pgfqpoint{0.464911in}{0.456793in}}{\pgfqpoint{1.977584in}{1.453311in}}%
\pgfusepath{clip}%
\pgfsetbuttcap%
\pgfsetmiterjoin%
\definecolor{currentfill}{rgb}{0.207059,0.752941,0.207059}%
\pgfsetfillcolor{currentfill}%
\pgfsetfillopacity{0.200000}%
\pgfsetlinewidth{0.000000pt}%
\definecolor{currentstroke}{rgb}{0.207059,0.752941,0.207059}%
\pgfsetstrokecolor{currentstroke}%
\pgfsetstrokeopacity{0.200000}%
\pgfsetdash{}{0pt}%
\pgfpathmoveto{\pgfqpoint{0.951117in}{0.456793in}}%
\pgfpathlineto{\pgfqpoint{1.108780in}{0.456793in}}%
\pgfpathlineto{\pgfqpoint{1.108780in}{1.910105in}}%
\pgfpathlineto{\pgfqpoint{0.951117in}{1.910105in}}%
\pgfpathlineto{\pgfqpoint{0.951117in}{0.456793in}}%
\pgfpathclose%
\pgfusepath{fill}%
\end{pgfscope}%
\begin{pgfscope}%
\pgfpathrectangle{\pgfqpoint{0.464911in}{0.456793in}}{\pgfqpoint{1.977584in}{1.453311in}}%
\pgfusepath{clip}%
\pgfsetbuttcap%
\pgfsetmiterjoin%
\definecolor{currentfill}{rgb}{0.696471,0.484706,0.889412}%
\pgfsetfillcolor{currentfill}%
\pgfsetfillopacity{0.200000}%
\pgfsetlinewidth{0.000000pt}%
\definecolor{currentstroke}{rgb}{0.696471,0.484706,0.889412}%
\pgfsetstrokecolor{currentstroke}%
\pgfsetstrokeopacity{0.200000}%
\pgfsetdash{}{0pt}%
\pgfpathmoveto{\pgfqpoint{1.108966in}{0.456793in}}%
\pgfpathlineto{\pgfqpoint{1.131117in}{0.456793in}}%
\pgfpathlineto{\pgfqpoint{1.131117in}{1.910105in}}%
\pgfpathlineto{\pgfqpoint{1.108966in}{1.910105in}}%
\pgfpathlineto{\pgfqpoint{1.108966in}{0.456793in}}%
\pgfpathclose%
\pgfusepath{fill}%
\end{pgfscope}%
\begin{pgfscope}%
\pgfpathrectangle{\pgfqpoint{0.464911in}{0.456793in}}{\pgfqpoint{1.977584in}{1.453311in}}%
\pgfusepath{clip}%
\pgfsetbuttcap%
\pgfsetmiterjoin%
\definecolor{currentfill}{rgb}{1.000000,0.498039,0.054902}%
\pgfsetfillcolor{currentfill}%
\pgfsetfillopacity{0.200000}%
\pgfsetlinewidth{0.000000pt}%
\definecolor{currentstroke}{rgb}{1.000000,0.498039,0.054902}%
\pgfsetstrokecolor{currentstroke}%
\pgfsetstrokeopacity{0.200000}%
\pgfsetdash{}{0pt}%
\pgfpathmoveto{\pgfqpoint{1.131304in}{0.456793in}}%
\pgfpathlineto{\pgfqpoint{1.188636in}{0.456793in}}%
\pgfpathlineto{\pgfqpoint{1.188636in}{1.910105in}}%
\pgfpathlineto{\pgfqpoint{1.131304in}{1.910105in}}%
\pgfpathlineto{\pgfqpoint{1.131304in}{0.456793in}}%
\pgfpathclose%
\pgfusepath{fill}%
\end{pgfscope}%
\begin{pgfscope}%
\pgfpathrectangle{\pgfqpoint{0.464911in}{0.456793in}}{\pgfqpoint{1.977584in}{1.453311in}}%
\pgfusepath{clip}%
\pgfsetbuttcap%
\pgfsetmiterjoin%
\definecolor{currentfill}{rgb}{0.121569,0.466667,0.705882}%
\pgfsetfillcolor{currentfill}%
\pgfsetfillopacity{0.200000}%
\pgfsetlinewidth{0.000000pt}%
\definecolor{currentstroke}{rgb}{0.121569,0.466667,0.705882}%
\pgfsetstrokecolor{currentstroke}%
\pgfsetstrokeopacity{0.200000}%
\pgfsetdash{}{0pt}%
\pgfpathmoveto{\pgfqpoint{1.201480in}{0.456793in}}%
\pgfpathlineto{\pgfqpoint{1.233496in}{0.456793in}}%
\pgfpathlineto{\pgfqpoint{1.233496in}{1.910105in}}%
\pgfpathlineto{\pgfqpoint{1.201480in}{1.910105in}}%
\pgfpathlineto{\pgfqpoint{1.201480in}{0.456793in}}%
\pgfpathclose%
\pgfusepath{fill}%
\end{pgfscope}%
\begin{pgfscope}%
\pgfpathrectangle{\pgfqpoint{0.464911in}{0.456793in}}{\pgfqpoint{1.977584in}{1.453311in}}%
\pgfusepath{clip}%
\pgfsetbuttcap%
\pgfsetmiterjoin%
\definecolor{currentfill}{rgb}{1.000000,0.183529,0.188235}%
\pgfsetfillcolor{currentfill}%
\pgfsetfillopacity{0.200000}%
\pgfsetlinewidth{0.000000pt}%
\definecolor{currentstroke}{rgb}{1.000000,0.183529,0.188235}%
\pgfsetstrokecolor{currentstroke}%
\pgfsetstrokeopacity{0.200000}%
\pgfsetdash{}{0pt}%
\pgfpathmoveto{\pgfqpoint{1.233682in}{0.456793in}}%
\pgfpathlineto{\pgfqpoint{1.423921in}{0.456793in}}%
\pgfpathlineto{\pgfqpoint{1.423921in}{1.910105in}}%
\pgfpathlineto{\pgfqpoint{1.233682in}{1.910105in}}%
\pgfpathlineto{\pgfqpoint{1.233682in}{0.456793in}}%
\pgfpathclose%
\pgfusepath{fill}%
\end{pgfscope}%
\begin{pgfscope}%
\pgfpathrectangle{\pgfqpoint{0.464911in}{0.456793in}}{\pgfqpoint{1.977584in}{1.453311in}}%
\pgfusepath{clip}%
\pgfsetbuttcap%
\pgfsetmiterjoin%
\definecolor{currentfill}{rgb}{0.696471,0.484706,0.889412}%
\pgfsetfillcolor{currentfill}%
\pgfsetfillopacity{0.200000}%
\pgfsetlinewidth{0.000000pt}%
\definecolor{currentstroke}{rgb}{0.696471,0.484706,0.889412}%
\pgfsetstrokecolor{currentstroke}%
\pgfsetstrokeopacity{0.200000}%
\pgfsetdash{}{0pt}%
\pgfpathmoveto{\pgfqpoint{1.467106in}{0.456793in}}%
\pgfpathlineto{\pgfqpoint{1.661253in}{0.456793in}}%
\pgfpathlineto{\pgfqpoint{1.661253in}{1.910105in}}%
\pgfpathlineto{\pgfqpoint{1.467106in}{1.910105in}}%
\pgfpathlineto{\pgfqpoint{1.467106in}{0.456793in}}%
\pgfpathclose%
\pgfusepath{fill}%
\end{pgfscope}%
\begin{pgfscope}%
\pgfpathrectangle{\pgfqpoint{0.464911in}{0.456793in}}{\pgfqpoint{1.977584in}{1.453311in}}%
\pgfusepath{clip}%
\pgfsetbuttcap%
\pgfsetmiterjoin%
\definecolor{currentfill}{rgb}{0.750000,0.750000,0.750000}%
\pgfsetfillcolor{currentfill}%
\pgfsetfillopacity{0.200000}%
\pgfsetlinewidth{0.000000pt}%
\definecolor{currentstroke}{rgb}{0.750000,0.750000,0.750000}%
\pgfsetstrokecolor{currentstroke}%
\pgfsetstrokeopacity{0.200000}%
\pgfsetdash{}{0pt}%
\pgfpathmoveto{\pgfqpoint{1.675214in}{0.456793in}}%
\pgfpathlineto{\pgfqpoint{1.697551in}{0.456793in}}%
\pgfpathlineto{\pgfqpoint{1.697551in}{1.910105in}}%
\pgfpathlineto{\pgfqpoint{1.675214in}{1.910105in}}%
\pgfpathlineto{\pgfqpoint{1.675214in}{0.456793in}}%
\pgfpathclose%
\pgfusepath{fill}%
\end{pgfscope}%
\begin{pgfscope}%
\pgfpathrectangle{\pgfqpoint{0.464911in}{0.456793in}}{\pgfqpoint{1.977584in}{1.453311in}}%
\pgfusepath{clip}%
\pgfsetbuttcap%
\pgfsetmiterjoin%
\definecolor{currentfill}{rgb}{0.750000,0.750000,0.750000}%
\pgfsetfillcolor{currentfill}%
\pgfsetfillopacity{0.200000}%
\pgfsetlinewidth{0.000000pt}%
\definecolor{currentstroke}{rgb}{0.750000,0.750000,0.750000}%
\pgfsetstrokecolor{currentstroke}%
\pgfsetstrokeopacity{0.200000}%
\pgfsetdash{}{0pt}%
\pgfpathmoveto{\pgfqpoint{1.737013in}{0.456793in}}%
\pgfpathlineto{\pgfqpoint{1.807748in}{0.456793in}}%
\pgfpathlineto{\pgfqpoint{1.807748in}{1.910105in}}%
\pgfpathlineto{\pgfqpoint{1.737013in}{1.910105in}}%
\pgfpathlineto{\pgfqpoint{1.737013in}{0.456793in}}%
\pgfpathclose%
\pgfusepath{fill}%
\end{pgfscope}%
\begin{pgfscope}%
\pgfpathrectangle{\pgfqpoint{0.464911in}{0.456793in}}{\pgfqpoint{1.977584in}{1.453311in}}%
\pgfusepath{clip}%
\pgfsetbuttcap%
\pgfsetmiterjoin%
\definecolor{currentfill}{rgb}{0.121569,0.466667,0.705882}%
\pgfsetfillcolor{currentfill}%
\pgfsetfillopacity{0.200000}%
\pgfsetlinewidth{0.000000pt}%
\definecolor{currentstroke}{rgb}{0.121569,0.466667,0.705882}%
\pgfsetstrokecolor{currentstroke}%
\pgfsetstrokeopacity{0.200000}%
\pgfsetdash{}{0pt}%
\pgfpathmoveto{\pgfqpoint{1.843487in}{0.456793in}}%
\pgfpathlineto{\pgfqpoint{1.884252in}{0.456793in}}%
\pgfpathlineto{\pgfqpoint{1.884252in}{1.910105in}}%
\pgfpathlineto{\pgfqpoint{1.843487in}{1.910105in}}%
\pgfpathlineto{\pgfqpoint{1.843487in}{0.456793in}}%
\pgfpathclose%
\pgfusepath{fill}%
\end{pgfscope}%
\begin{pgfscope}%
\pgfpathrectangle{\pgfqpoint{0.464911in}{0.456793in}}{\pgfqpoint{1.977584in}{1.453311in}}%
\pgfusepath{clip}%
\pgfsetbuttcap%
\pgfsetmiterjoin%
\definecolor{currentfill}{rgb}{0.671373,0.122353,0.125490}%
\pgfsetfillcolor{currentfill}%
\pgfsetfillopacity{0.200000}%
\pgfsetlinewidth{0.000000pt}%
\definecolor{currentstroke}{rgb}{0.671373,0.122353,0.125490}%
\pgfsetstrokecolor{currentstroke}%
\pgfsetstrokeopacity{0.200000}%
\pgfsetdash{}{0pt}%
\pgfpathmoveto{\pgfqpoint{1.884439in}{0.456793in}}%
\pgfpathlineto{\pgfqpoint{2.071885in}{0.456793in}}%
\pgfpathlineto{\pgfqpoint{2.071885in}{1.910105in}}%
\pgfpathlineto{\pgfqpoint{1.884439in}{1.910105in}}%
\pgfpathlineto{\pgfqpoint{1.884439in}{0.456793in}}%
\pgfpathclose%
\pgfusepath{fill}%
\end{pgfscope}%
\begin{pgfscope}%
\pgfpathrectangle{\pgfqpoint{0.464911in}{0.456793in}}{\pgfqpoint{1.977584in}{1.453311in}}%
\pgfusepath{clip}%
\pgfsetbuttcap%
\pgfsetmiterjoin%
\definecolor{currentfill}{rgb}{0.464314,0.323137,0.592941}%
\pgfsetfillcolor{currentfill}%
\pgfsetfillopacity{0.200000}%
\pgfsetlinewidth{0.000000pt}%
\definecolor{currentstroke}{rgb}{0.464314,0.323137,0.592941}%
\pgfsetstrokecolor{currentstroke}%
\pgfsetstrokeopacity{0.200000}%
\pgfsetdash{}{0pt}%
\pgfpathmoveto{\pgfqpoint{2.117117in}{0.456793in}}%
\pgfpathlineto{\pgfqpoint{2.301958in}{0.456793in}}%
\pgfpathlineto{\pgfqpoint{2.301958in}{1.910105in}}%
\pgfpathlineto{\pgfqpoint{2.117117in}{1.910105in}}%
\pgfpathlineto{\pgfqpoint{2.117117in}{0.456793in}}%
\pgfpathclose%
\pgfusepath{fill}%
\end{pgfscope}%
\begin{pgfscope}%
\pgfpathrectangle{\pgfqpoint{0.464911in}{0.456793in}}{\pgfqpoint{1.977584in}{1.453311in}}%
\pgfusepath{clip}%
\pgfsetbuttcap%
\pgfsetmiterjoin%
\definecolor{currentfill}{rgb}{0.750000,0.750000,0.750000}%
\pgfsetfillcolor{currentfill}%
\pgfsetfillopacity{0.200000}%
\pgfsetlinewidth{0.000000pt}%
\definecolor{currentstroke}{rgb}{0.750000,0.750000,0.750000}%
\pgfsetstrokecolor{currentstroke}%
\pgfsetstrokeopacity{0.200000}%
\pgfsetdash{}{0pt}%
\pgfpathmoveto{\pgfqpoint{2.316291in}{0.456793in}}%
\pgfpathlineto{\pgfqpoint{2.341978in}{0.456793in}}%
\pgfpathlineto{\pgfqpoint{2.341978in}{1.910105in}}%
\pgfpathlineto{\pgfqpoint{2.316291in}{1.910105in}}%
\pgfpathlineto{\pgfqpoint{2.316291in}{0.456793in}}%
\pgfpathclose%
\pgfusepath{fill}%
\end{pgfscope}%
\begin{pgfscope}%
\pgfpathrectangle{\pgfqpoint{0.464911in}{0.456793in}}{\pgfqpoint{1.977584in}{1.453311in}}%
\pgfusepath{clip}%
\pgfsetbuttcap%
\pgfsetmiterjoin%
\definecolor{currentfill}{rgb}{0.750000,0.750000,0.750000}%
\pgfsetfillcolor{currentfill}%
\pgfsetfillopacity{0.200000}%
\pgfsetlinewidth{0.000000pt}%
\definecolor{currentstroke}{rgb}{0.750000,0.750000,0.750000}%
\pgfsetstrokecolor{currentstroke}%
\pgfsetstrokeopacity{0.200000}%
\pgfsetdash{}{0pt}%
\pgfpathmoveto{\pgfqpoint{2.382930in}{0.456793in}}%
\pgfpathlineto{\pgfqpoint{2.442496in}{0.456793in}}%
\pgfpathlineto{\pgfqpoint{2.442496in}{1.910105in}}%
\pgfpathlineto{\pgfqpoint{2.382930in}{1.910105in}}%
\pgfpathlineto{\pgfqpoint{2.382930in}{0.456793in}}%
\pgfpathclose%
\pgfusepath{fill}%
\end{pgfscope}%
\begin{pgfscope}%
\pgfsetbuttcap%
\pgfsetroundjoin%
\definecolor{currentfill}{rgb}{0.000000,0.000000,0.000000}%
\pgfsetfillcolor{currentfill}%
\pgfsetlinewidth{0.803000pt}%
\definecolor{currentstroke}{rgb}{0.000000,0.000000,0.000000}%
\pgfsetstrokecolor{currentstroke}%
\pgfsetdash{}{0pt}%
\pgfsys@defobject{currentmarker}{\pgfqpoint{0.000000in}{-0.048611in}}{\pgfqpoint{0.000000in}{0.000000in}}{%
\pgfpathmoveto{\pgfqpoint{0.000000in}{0.000000in}}%
\pgfpathlineto{\pgfqpoint{0.000000in}{-0.048611in}}%
\pgfusepath{stroke,fill}%
}%
\begin{pgfscope}%
\pgfsys@transformshift{0.464911in}{0.456793in}%
\pgfsys@useobject{currentmarker}{}%
\end{pgfscope}%
\end{pgfscope}%
\begin{pgfscope}%
\definecolor{textcolor}{rgb}{0.000000,0.000000,0.000000}%
\pgfsetstrokecolor{textcolor}%
\pgfsetfillcolor{textcolor}%
\pgftext[x=0.464911in,y=0.359571in,,top]{\color{textcolor}{\rmfamily\fontsize{10.000000}{12.000000}\selectfont\catcode`\^=\active\def^{\ifmmode\sp\else\^{}\fi}\catcode`\%=\active\def%{\%}0}}%
\end{pgfscope}%
\begin{pgfscope}%
\pgfsetbuttcap%
\pgfsetroundjoin%
\definecolor{currentfill}{rgb}{0.000000,0.000000,0.000000}%
\pgfsetfillcolor{currentfill}%
\pgfsetlinewidth{0.803000pt}%
\definecolor{currentstroke}{rgb}{0.000000,0.000000,0.000000}%
\pgfsetstrokecolor{currentstroke}%
\pgfsetdash{}{0pt}%
\pgfsys@defobject{currentmarker}{\pgfqpoint{0.000000in}{-0.048611in}}{\pgfqpoint{0.000000in}{0.000000in}}{%
\pgfpathmoveto{\pgfqpoint{0.000000in}{0.000000in}}%
\pgfpathlineto{\pgfqpoint{0.000000in}{-0.048611in}}%
\pgfusepath{stroke,fill}%
}%
\begin{pgfscope}%
\pgfsys@transformshift{1.395627in}{0.456793in}%
\pgfsys@useobject{currentmarker}{}%
\end{pgfscope}%
\end{pgfscope}%
\begin{pgfscope}%
\definecolor{textcolor}{rgb}{0.000000,0.000000,0.000000}%
\pgfsetstrokecolor{textcolor}%
\pgfsetfillcolor{textcolor}%
\pgftext[x=1.395627in,y=0.359571in,,top]{\color{textcolor}{\rmfamily\fontsize{10.000000}{12.000000}\selectfont\catcode`\^=\active\def^{\ifmmode\sp\else\^{}\fi}\catcode`\%=\active\def%{\%}50}}%
\end{pgfscope}%
\begin{pgfscope}%
\pgfsetbuttcap%
\pgfsetroundjoin%
\definecolor{currentfill}{rgb}{0.000000,0.000000,0.000000}%
\pgfsetfillcolor{currentfill}%
\pgfsetlinewidth{0.803000pt}%
\definecolor{currentstroke}{rgb}{0.000000,0.000000,0.000000}%
\pgfsetstrokecolor{currentstroke}%
\pgfsetdash{}{0pt}%
\pgfsys@defobject{currentmarker}{\pgfqpoint{0.000000in}{-0.048611in}}{\pgfqpoint{0.000000in}{0.000000in}}{%
\pgfpathmoveto{\pgfqpoint{0.000000in}{0.000000in}}%
\pgfpathlineto{\pgfqpoint{0.000000in}{-0.048611in}}%
\pgfusepath{stroke,fill}%
}%
\begin{pgfscope}%
\pgfsys@transformshift{2.326342in}{0.456793in}%
\pgfsys@useobject{currentmarker}{}%
\end{pgfscope}%
\end{pgfscope}%
\begin{pgfscope}%
\definecolor{textcolor}{rgb}{0.000000,0.000000,0.000000}%
\pgfsetstrokecolor{textcolor}%
\pgfsetfillcolor{textcolor}%
\pgftext[x=2.326342in,y=0.359571in,,top]{\color{textcolor}{\rmfamily\fontsize{10.000000}{12.000000}\selectfont\catcode`\^=\active\def^{\ifmmode\sp\else\^{}\fi}\catcode`\%=\active\def%{\%}100}}%
\end{pgfscope}%
\begin{pgfscope}%
\definecolor{textcolor}{rgb}{0.000000,0.000000,0.000000}%
\pgfsetstrokecolor{textcolor}%
\pgfsetfillcolor{textcolor}%
\pgftext[x=1.453703in,y=0.180559in,,top]{\color{textcolor}{\rmfamily\fontsize{10.000000}{12.000000}\selectfont\catcode`\^=\active\def^{\ifmmode\sp\else\^{}\fi}\catcode`\%=\active\def%{\%}$t$ / \unit{\second}}}%
\end{pgfscope}%
\begin{pgfscope}%
\pgfpathrectangle{\pgfqpoint{0.464911in}{0.456793in}}{\pgfqpoint{1.977584in}{1.453311in}}%
\pgfusepath{clip}%
\pgfsetrectcap%
\pgfsetroundjoin%
\pgfsetlinewidth{0.501875pt}%
\definecolor{currentstroke}{rgb}{0.690196,0.690196,0.690196}%
\pgfsetstrokecolor{currentstroke}%
\pgfsetstrokeopacity{0.300000}%
\pgfsetdash{}{0pt}%
\pgfpathmoveto{\pgfqpoint{0.464911in}{0.456793in}}%
\pgfpathlineto{\pgfqpoint{2.442496in}{0.456793in}}%
\pgfusepath{stroke}%
\end{pgfscope}%
\begin{pgfscope}%
\pgfsetbuttcap%
\pgfsetroundjoin%
\definecolor{currentfill}{rgb}{0.000000,0.000000,0.000000}%
\pgfsetfillcolor{currentfill}%
\pgfsetlinewidth{0.803000pt}%
\definecolor{currentstroke}{rgb}{0.000000,0.000000,0.000000}%
\pgfsetstrokecolor{currentstroke}%
\pgfsetdash{}{0pt}%
\pgfsys@defobject{currentmarker}{\pgfqpoint{-0.048611in}{0.000000in}}{\pgfqpoint{-0.000000in}{0.000000in}}{%
\pgfpathmoveto{\pgfqpoint{-0.000000in}{0.000000in}}%
\pgfpathlineto{\pgfqpoint{-0.048611in}{0.000000in}}%
\pgfusepath{stroke,fill}%
}%
\begin{pgfscope}%
\pgfsys@transformshift{0.464911in}{0.456793in}%
\pgfsys@useobject{currentmarker}{}%
\end{pgfscope}%
\end{pgfscope}%
\begin{pgfscope}%
\definecolor{textcolor}{rgb}{0.000000,0.000000,0.000000}%
\pgfsetstrokecolor{textcolor}%
\pgfsetfillcolor{textcolor}%
\pgftext[x=0.251948in, y=0.408568in, left, base]{\color{textcolor}{\rmfamily\fontsize{10.000000}{12.000000}\selectfont\catcode`\^=\active\def^{\ifmmode\sp\else\^{}\fi}\catcode`\%=\active\def%{\%}-6}}%
\end{pgfscope}%
\begin{pgfscope}%
\pgfpathrectangle{\pgfqpoint{0.464911in}{0.456793in}}{\pgfqpoint{1.977584in}{1.453311in}}%
\pgfusepath{clip}%
\pgfsetrectcap%
\pgfsetroundjoin%
\pgfsetlinewidth{0.501875pt}%
\definecolor{currentstroke}{rgb}{0.690196,0.690196,0.690196}%
\pgfsetstrokecolor{currentstroke}%
\pgfsetstrokeopacity{0.300000}%
\pgfsetdash{}{0pt}%
\pgfpathmoveto{\pgfqpoint{0.464911in}{0.699012in}}%
\pgfpathlineto{\pgfqpoint{2.442496in}{0.699012in}}%
\pgfusepath{stroke}%
\end{pgfscope}%
\begin{pgfscope}%
\pgfsetbuttcap%
\pgfsetroundjoin%
\definecolor{currentfill}{rgb}{0.000000,0.000000,0.000000}%
\pgfsetfillcolor{currentfill}%
\pgfsetlinewidth{0.803000pt}%
\definecolor{currentstroke}{rgb}{0.000000,0.000000,0.000000}%
\pgfsetstrokecolor{currentstroke}%
\pgfsetdash{}{0pt}%
\pgfsys@defobject{currentmarker}{\pgfqpoint{-0.048611in}{0.000000in}}{\pgfqpoint{-0.000000in}{0.000000in}}{%
\pgfpathmoveto{\pgfqpoint{-0.000000in}{0.000000in}}%
\pgfpathlineto{\pgfqpoint{-0.048611in}{0.000000in}}%
\pgfusepath{stroke,fill}%
}%
\begin{pgfscope}%
\pgfsys@transformshift{0.464911in}{0.699012in}%
\pgfsys@useobject{currentmarker}{}%
\end{pgfscope}%
\end{pgfscope}%
\begin{pgfscope}%
\definecolor{textcolor}{rgb}{0.000000,0.000000,0.000000}%
\pgfsetstrokecolor{textcolor}%
\pgfsetfillcolor{textcolor}%
\pgftext[x=0.251948in, y=0.650787in, left, base]{\color{textcolor}{\rmfamily\fontsize{10.000000}{12.000000}\selectfont\catcode`\^=\active\def^{\ifmmode\sp\else\^{}\fi}\catcode`\%=\active\def%{\%}-4}}%
\end{pgfscope}%
\begin{pgfscope}%
\pgfpathrectangle{\pgfqpoint{0.464911in}{0.456793in}}{\pgfqpoint{1.977584in}{1.453311in}}%
\pgfusepath{clip}%
\pgfsetrectcap%
\pgfsetroundjoin%
\pgfsetlinewidth{0.501875pt}%
\definecolor{currentstroke}{rgb}{0.690196,0.690196,0.690196}%
\pgfsetstrokecolor{currentstroke}%
\pgfsetstrokeopacity{0.300000}%
\pgfsetdash{}{0pt}%
\pgfpathmoveto{\pgfqpoint{0.464911in}{0.941230in}}%
\pgfpathlineto{\pgfqpoint{2.442496in}{0.941230in}}%
\pgfusepath{stroke}%
\end{pgfscope}%
\begin{pgfscope}%
\pgfsetbuttcap%
\pgfsetroundjoin%
\definecolor{currentfill}{rgb}{0.000000,0.000000,0.000000}%
\pgfsetfillcolor{currentfill}%
\pgfsetlinewidth{0.803000pt}%
\definecolor{currentstroke}{rgb}{0.000000,0.000000,0.000000}%
\pgfsetstrokecolor{currentstroke}%
\pgfsetdash{}{0pt}%
\pgfsys@defobject{currentmarker}{\pgfqpoint{-0.048611in}{0.000000in}}{\pgfqpoint{-0.000000in}{0.000000in}}{%
\pgfpathmoveto{\pgfqpoint{-0.000000in}{0.000000in}}%
\pgfpathlineto{\pgfqpoint{-0.048611in}{0.000000in}}%
\pgfusepath{stroke,fill}%
}%
\begin{pgfscope}%
\pgfsys@transformshift{0.464911in}{0.941230in}%
\pgfsys@useobject{currentmarker}{}%
\end{pgfscope}%
\end{pgfscope}%
\begin{pgfscope}%
\definecolor{textcolor}{rgb}{0.000000,0.000000,0.000000}%
\pgfsetstrokecolor{textcolor}%
\pgfsetfillcolor{textcolor}%
\pgftext[x=0.251948in, y=0.893005in, left, base]{\color{textcolor}{\rmfamily\fontsize{10.000000}{12.000000}\selectfont\catcode`\^=\active\def^{\ifmmode\sp\else\^{}\fi}\catcode`\%=\active\def%{\%}-2}}%
\end{pgfscope}%
\begin{pgfscope}%
\pgfpathrectangle{\pgfqpoint{0.464911in}{0.456793in}}{\pgfqpoint{1.977584in}{1.453311in}}%
\pgfusepath{clip}%
\pgfsetrectcap%
\pgfsetroundjoin%
\pgfsetlinewidth{0.501875pt}%
\definecolor{currentstroke}{rgb}{0.690196,0.690196,0.690196}%
\pgfsetstrokecolor{currentstroke}%
\pgfsetstrokeopacity{0.300000}%
\pgfsetdash{}{0pt}%
\pgfpathmoveto{\pgfqpoint{0.464911in}{1.183449in}}%
\pgfpathlineto{\pgfqpoint{2.442496in}{1.183449in}}%
\pgfusepath{stroke}%
\end{pgfscope}%
\begin{pgfscope}%
\pgfsetbuttcap%
\pgfsetroundjoin%
\definecolor{currentfill}{rgb}{0.000000,0.000000,0.000000}%
\pgfsetfillcolor{currentfill}%
\pgfsetlinewidth{0.803000pt}%
\definecolor{currentstroke}{rgb}{0.000000,0.000000,0.000000}%
\pgfsetstrokecolor{currentstroke}%
\pgfsetdash{}{0pt}%
\pgfsys@defobject{currentmarker}{\pgfqpoint{-0.048611in}{0.000000in}}{\pgfqpoint{-0.000000in}{0.000000in}}{%
\pgfpathmoveto{\pgfqpoint{-0.000000in}{0.000000in}}%
\pgfpathlineto{\pgfqpoint{-0.048611in}{0.000000in}}%
\pgfusepath{stroke,fill}%
}%
\begin{pgfscope}%
\pgfsys@transformshift{0.464911in}{1.183449in}%
\pgfsys@useobject{currentmarker}{}%
\end{pgfscope}%
\end{pgfscope}%
\begin{pgfscope}%
\definecolor{textcolor}{rgb}{0.000000,0.000000,0.000000}%
\pgfsetstrokecolor{textcolor}%
\pgfsetfillcolor{textcolor}%
\pgftext[x=0.298244in, y=1.135224in, left, base]{\color{textcolor}{\rmfamily\fontsize{10.000000}{12.000000}\selectfont\catcode`\^=\active\def^{\ifmmode\sp\else\^{}\fi}\catcode`\%=\active\def%{\%}0}}%
\end{pgfscope}%
\begin{pgfscope}%
\pgfpathrectangle{\pgfqpoint{0.464911in}{0.456793in}}{\pgfqpoint{1.977584in}{1.453311in}}%
\pgfusepath{clip}%
\pgfsetrectcap%
\pgfsetroundjoin%
\pgfsetlinewidth{0.501875pt}%
\definecolor{currentstroke}{rgb}{0.690196,0.690196,0.690196}%
\pgfsetstrokecolor{currentstroke}%
\pgfsetstrokeopacity{0.300000}%
\pgfsetdash{}{0pt}%
\pgfpathmoveto{\pgfqpoint{0.464911in}{1.425668in}}%
\pgfpathlineto{\pgfqpoint{2.442496in}{1.425668in}}%
\pgfusepath{stroke}%
\end{pgfscope}%
\begin{pgfscope}%
\pgfsetbuttcap%
\pgfsetroundjoin%
\definecolor{currentfill}{rgb}{0.000000,0.000000,0.000000}%
\pgfsetfillcolor{currentfill}%
\pgfsetlinewidth{0.803000pt}%
\definecolor{currentstroke}{rgb}{0.000000,0.000000,0.000000}%
\pgfsetstrokecolor{currentstroke}%
\pgfsetdash{}{0pt}%
\pgfsys@defobject{currentmarker}{\pgfqpoint{-0.048611in}{0.000000in}}{\pgfqpoint{-0.000000in}{0.000000in}}{%
\pgfpathmoveto{\pgfqpoint{-0.000000in}{0.000000in}}%
\pgfpathlineto{\pgfqpoint{-0.048611in}{0.000000in}}%
\pgfusepath{stroke,fill}%
}%
\begin{pgfscope}%
\pgfsys@transformshift{0.464911in}{1.425668in}%
\pgfsys@useobject{currentmarker}{}%
\end{pgfscope}%
\end{pgfscope}%
\begin{pgfscope}%
\definecolor{textcolor}{rgb}{0.000000,0.000000,0.000000}%
\pgfsetstrokecolor{textcolor}%
\pgfsetfillcolor{textcolor}%
\pgftext[x=0.298244in, y=1.377442in, left, base]{\color{textcolor}{\rmfamily\fontsize{10.000000}{12.000000}\selectfont\catcode`\^=\active\def^{\ifmmode\sp\else\^{}\fi}\catcode`\%=\active\def%{\%}2}}%
\end{pgfscope}%
\begin{pgfscope}%
\pgfpathrectangle{\pgfqpoint{0.464911in}{0.456793in}}{\pgfqpoint{1.977584in}{1.453311in}}%
\pgfusepath{clip}%
\pgfsetrectcap%
\pgfsetroundjoin%
\pgfsetlinewidth{0.501875pt}%
\definecolor{currentstroke}{rgb}{0.690196,0.690196,0.690196}%
\pgfsetstrokecolor{currentstroke}%
\pgfsetstrokeopacity{0.300000}%
\pgfsetdash{}{0pt}%
\pgfpathmoveto{\pgfqpoint{0.464911in}{1.667886in}}%
\pgfpathlineto{\pgfqpoint{2.442496in}{1.667886in}}%
\pgfusepath{stroke}%
\end{pgfscope}%
\begin{pgfscope}%
\pgfsetbuttcap%
\pgfsetroundjoin%
\definecolor{currentfill}{rgb}{0.000000,0.000000,0.000000}%
\pgfsetfillcolor{currentfill}%
\pgfsetlinewidth{0.803000pt}%
\definecolor{currentstroke}{rgb}{0.000000,0.000000,0.000000}%
\pgfsetstrokecolor{currentstroke}%
\pgfsetdash{}{0pt}%
\pgfsys@defobject{currentmarker}{\pgfqpoint{-0.048611in}{0.000000in}}{\pgfqpoint{-0.000000in}{0.000000in}}{%
\pgfpathmoveto{\pgfqpoint{-0.000000in}{0.000000in}}%
\pgfpathlineto{\pgfqpoint{-0.048611in}{0.000000in}}%
\pgfusepath{stroke,fill}%
}%
\begin{pgfscope}%
\pgfsys@transformshift{0.464911in}{1.667886in}%
\pgfsys@useobject{currentmarker}{}%
\end{pgfscope}%
\end{pgfscope}%
\begin{pgfscope}%
\definecolor{textcolor}{rgb}{0.000000,0.000000,0.000000}%
\pgfsetstrokecolor{textcolor}%
\pgfsetfillcolor{textcolor}%
\pgftext[x=0.298244in, y=1.619661in, left, base]{\color{textcolor}{\rmfamily\fontsize{10.000000}{12.000000}\selectfont\catcode`\^=\active\def^{\ifmmode\sp\else\^{}\fi}\catcode`\%=\active\def%{\%}4}}%
\end{pgfscope}%
\begin{pgfscope}%
\pgfpathrectangle{\pgfqpoint{0.464911in}{0.456793in}}{\pgfqpoint{1.977584in}{1.453311in}}%
\pgfusepath{clip}%
\pgfsetrectcap%
\pgfsetroundjoin%
\pgfsetlinewidth{0.501875pt}%
\definecolor{currentstroke}{rgb}{0.690196,0.690196,0.690196}%
\pgfsetstrokecolor{currentstroke}%
\pgfsetstrokeopacity{0.300000}%
\pgfsetdash{}{0pt}%
\pgfpathmoveto{\pgfqpoint{0.464911in}{1.910105in}}%
\pgfpathlineto{\pgfqpoint{2.442496in}{1.910105in}}%
\pgfusepath{stroke}%
\end{pgfscope}%
\begin{pgfscope}%
\pgfsetbuttcap%
\pgfsetroundjoin%
\definecolor{currentfill}{rgb}{0.000000,0.000000,0.000000}%
\pgfsetfillcolor{currentfill}%
\pgfsetlinewidth{0.803000pt}%
\definecolor{currentstroke}{rgb}{0.000000,0.000000,0.000000}%
\pgfsetstrokecolor{currentstroke}%
\pgfsetdash{}{0pt}%
\pgfsys@defobject{currentmarker}{\pgfqpoint{-0.048611in}{0.000000in}}{\pgfqpoint{-0.000000in}{0.000000in}}{%
\pgfpathmoveto{\pgfqpoint{-0.000000in}{0.000000in}}%
\pgfpathlineto{\pgfqpoint{-0.048611in}{0.000000in}}%
\pgfusepath{stroke,fill}%
}%
\begin{pgfscope}%
\pgfsys@transformshift{0.464911in}{1.910105in}%
\pgfsys@useobject{currentmarker}{}%
\end{pgfscope}%
\end{pgfscope}%
\begin{pgfscope}%
\definecolor{textcolor}{rgb}{0.000000,0.000000,0.000000}%
\pgfsetstrokecolor{textcolor}%
\pgfsetfillcolor{textcolor}%
\pgftext[x=0.298244in, y=1.861879in, left, base]{\color{textcolor}{\rmfamily\fontsize{10.000000}{12.000000}\selectfont\catcode`\^=\active\def^{\ifmmode\sp\else\^{}\fi}\catcode`\%=\active\def%{\%}6}}%
\end{pgfscope}%
\begin{pgfscope}%
\definecolor{textcolor}{rgb}{0.000000,0.000000,0.000000}%
\pgfsetstrokecolor{textcolor}%
\pgfsetfillcolor{textcolor}%
\pgftext[x=0.196393in,y=1.183449in,,bottom,rotate=90.000000]{\color{textcolor}{\rmfamily\fontsize{10.000000}{12.000000}\selectfont\catcode`\^=\active\def^{\ifmmode\sp\else\^{}\fi}\catcode`\%=\active\def%{\%}$a$ / \unit{\metre\per\second\squared}}}%
\end{pgfscope}%
\begin{pgfscope}%
\pgfpathrectangle{\pgfqpoint{0.464911in}{0.456793in}}{\pgfqpoint{1.977584in}{1.453311in}}%
\pgfusepath{clip}%
\pgfsetrectcap%
\pgfsetroundjoin%
\pgfsetlinewidth{0.803000pt}%
\definecolor{currentstroke}{rgb}{0.121569,0.466667,0.705882}%
\pgfsetstrokecolor{currentstroke}%
\pgfsetdash{}{0pt}%
\pgfpathmoveto{\pgfqpoint{0.464911in}{1.183445in}}%
\pgfpathlineto{\pgfqpoint{0.500092in}{1.183224in}}%
\pgfpathlineto{\pgfqpoint{0.501768in}{1.203650in}}%
\pgfpathlineto{\pgfqpoint{0.504560in}{1.257546in}}%
\pgfpathlineto{\pgfqpoint{0.508096in}{1.261969in}}%
\pgfpathlineto{\pgfqpoint{0.512750in}{1.268483in}}%
\pgfpathlineto{\pgfqpoint{0.513495in}{1.269636in}}%
\pgfpathlineto{\pgfqpoint{0.514798in}{1.299361in}}%
\pgfpathlineto{\pgfqpoint{0.516659in}{1.396998in}}%
\pgfpathlineto{\pgfqpoint{0.518520in}{1.628460in}}%
\pgfpathlineto{\pgfqpoint{0.519823in}{1.685607in}}%
\pgfpathlineto{\pgfqpoint{0.528014in}{1.685607in}}%
\pgfpathlineto{\pgfqpoint{0.529689in}{1.630694in}}%
\pgfpathlineto{\pgfqpoint{0.537693in}{1.630694in}}%
\pgfpathlineto{\pgfqpoint{0.540858in}{1.654489in}}%
\pgfpathlineto{\pgfqpoint{0.541416in}{1.655554in}}%
\pgfpathlineto{\pgfqpoint{0.543464in}{1.660967in}}%
\pgfpathlineto{\pgfqpoint{0.545139in}{1.661171in}}%
\pgfpathlineto{\pgfqpoint{0.551468in}{1.662264in}}%
\pgfpathlineto{\pgfqpoint{0.553515in}{1.662901in}}%
\pgfpathlineto{\pgfqpoint{0.554074in}{1.662901in}}%
\pgfpathlineto{\pgfqpoint{0.554818in}{1.637045in}}%
\pgfpathlineto{\pgfqpoint{0.556494in}{1.619147in}}%
\pgfpathlineto{\pgfqpoint{0.558727in}{1.611859in}}%
\pgfpathlineto{\pgfqpoint{0.563567in}{1.612372in}}%
\pgfpathlineto{\pgfqpoint{0.568593in}{1.637489in}}%
\pgfpathlineto{\pgfqpoint{0.570268in}{1.639871in}}%
\pgfpathlineto{\pgfqpoint{0.571013in}{1.640620in}}%
\pgfpathlineto{\pgfqpoint{0.576039in}{1.655998in}}%
\pgfpathlineto{\pgfqpoint{0.578086in}{1.657410in}}%
\pgfpathlineto{\pgfqpoint{0.589069in}{1.656931in}}%
\pgfpathlineto{\pgfqpoint{0.593908in}{1.639143in}}%
\pgfpathlineto{\pgfqpoint{0.595770in}{1.637677in}}%
\pgfpathlineto{\pgfqpoint{0.597631in}{1.637522in}}%
\pgfpathlineto{\pgfqpoint{0.599307in}{1.636422in}}%
\pgfpathlineto{\pgfqpoint{0.600237in}{1.635199in}}%
\pgfpathlineto{\pgfqpoint{0.602285in}{1.607994in}}%
\pgfpathlineto{\pgfqpoint{0.608055in}{1.607012in}}%
\pgfpathlineto{\pgfqpoint{0.610475in}{1.606405in}}%
\pgfpathlineto{\pgfqpoint{0.615873in}{1.605315in}}%
\pgfpathlineto{\pgfqpoint{0.618107in}{1.604829in}}%
\pgfpathlineto{\pgfqpoint{0.620713in}{1.603733in}}%
\pgfpathlineto{\pgfqpoint{0.624064in}{1.599693in}}%
\pgfpathlineto{\pgfqpoint{0.625739in}{1.587702in}}%
\pgfpathlineto{\pgfqpoint{0.629089in}{1.552053in}}%
\pgfpathlineto{\pgfqpoint{0.634301in}{1.498155in}}%
\pgfpathlineto{\pgfqpoint{0.634488in}{1.498155in}}%
\pgfpathlineto{\pgfqpoint{0.635977in}{1.475796in}}%
\pgfpathlineto{\pgfqpoint{0.638210in}{1.435799in}}%
\pgfpathlineto{\pgfqpoint{0.641747in}{1.408048in}}%
\pgfpathlineto{\pgfqpoint{0.643795in}{1.395293in}}%
\pgfpathlineto{\pgfqpoint{0.645284in}{1.388511in}}%
\pgfpathlineto{\pgfqpoint{0.646028in}{1.387463in}}%
\pgfpathlineto{\pgfqpoint{0.647145in}{1.379588in}}%
\pgfpathlineto{\pgfqpoint{0.650496in}{1.317718in}}%
\pgfpathlineto{\pgfqpoint{0.653288in}{1.304755in}}%
\pgfpathlineto{\pgfqpoint{0.655149in}{1.298302in}}%
\pgfpathlineto{\pgfqpoint{0.656639in}{1.276946in}}%
\pgfpathlineto{\pgfqpoint{0.658500in}{1.271626in}}%
\pgfpathlineto{\pgfqpoint{0.662409in}{1.260754in}}%
\pgfpathlineto{\pgfqpoint{0.666132in}{1.254219in}}%
\pgfpathlineto{\pgfqpoint{0.666504in}{1.253568in}}%
\pgfpathlineto{\pgfqpoint{0.670786in}{1.242650in}}%
\pgfpathlineto{\pgfqpoint{0.672089in}{1.242092in}}%
\pgfpathlineto{\pgfqpoint{0.672833in}{1.240988in}}%
\pgfpathlineto{\pgfqpoint{0.674322in}{1.237050in}}%
\pgfpathlineto{\pgfqpoint{0.676928in}{1.220457in}}%
\pgfpathlineto{\pgfqpoint{0.681768in}{1.209135in}}%
\pgfpathlineto{\pgfqpoint{0.683257in}{1.200967in}}%
\pgfpathlineto{\pgfqpoint{0.692006in}{1.195774in}}%
\pgfpathlineto{\pgfqpoint{0.694984in}{1.194700in}}%
\pgfpathlineto{\pgfqpoint{0.705967in}{1.189107in}}%
\pgfpathlineto{\pgfqpoint{0.716391in}{1.183696in}}%
\pgfpathlineto{\pgfqpoint{0.724767in}{1.167176in}}%
\pgfpathlineto{\pgfqpoint{0.730724in}{1.143094in}}%
\pgfpathlineto{\pgfqpoint{0.735005in}{1.123852in}}%
\pgfpathlineto{\pgfqpoint{0.737797in}{1.118367in}}%
\pgfpathlineto{\pgfqpoint{0.745057in}{0.719210in}}%
\pgfpathlineto{\pgfqpoint{0.746918in}{0.705766in}}%
\pgfpathlineto{\pgfqpoint{0.755294in}{0.705766in}}%
\pgfpathlineto{\pgfqpoint{0.761809in}{0.769345in}}%
\pgfpathlineto{\pgfqpoint{0.777259in}{0.939379in}}%
\pgfpathlineto{\pgfqpoint{0.784333in}{1.045909in}}%
\pgfpathlineto{\pgfqpoint{0.789917in}{1.097782in}}%
\pgfpathlineto{\pgfqpoint{0.795129in}{1.127718in}}%
\pgfpathlineto{\pgfqpoint{0.797177in}{1.128106in}}%
\pgfpathlineto{\pgfqpoint{0.804250in}{1.129899in}}%
\pgfpathlineto{\pgfqpoint{0.809648in}{1.146753in}}%
\pgfpathlineto{\pgfqpoint{0.812999in}{1.147848in}}%
\pgfpathlineto{\pgfqpoint{0.815233in}{1.151300in}}%
\pgfpathlineto{\pgfqpoint{0.820817in}{1.159231in}}%
\pgfpathlineto{\pgfqpoint{0.824726in}{1.163384in}}%
\pgfpathlineto{\pgfqpoint{0.827332in}{1.175978in}}%
\pgfpathlineto{\pgfqpoint{0.832172in}{1.186509in}}%
\pgfpathlineto{\pgfqpoint{0.847994in}{1.192918in}}%
\pgfpathlineto{\pgfqpoint{0.857115in}{1.193861in}}%
\pgfpathlineto{\pgfqpoint{0.873123in}{1.192755in}}%
\pgfpathlineto{\pgfqpoint{0.920217in}{1.185903in}}%
\pgfpathlineto{\pgfqpoint{0.937901in}{1.186305in}}%
\pgfpathlineto{\pgfqpoint{0.947394in}{1.183781in}}%
\pgfpathlineto{\pgfqpoint{0.966567in}{1.178143in}}%
\pgfpathlineto{\pgfqpoint{0.979597in}{1.179236in}}%
\pgfpathlineto{\pgfqpoint{1.006960in}{1.184116in}}%
\pgfpathlineto{\pgfqpoint{1.028739in}{1.185307in}}%
\pgfpathlineto{\pgfqpoint{1.071179in}{1.184672in}}%
\pgfpathlineto{\pgfqpoint{1.083837in}{1.183609in}}%
\pgfpathlineto{\pgfqpoint{1.089607in}{1.178030in}}%
\pgfpathlineto{\pgfqpoint{1.093517in}{1.171660in}}%
\pgfpathlineto{\pgfqpoint{1.095936in}{1.162194in}}%
\pgfpathlineto{\pgfqpoint{1.098170in}{1.144933in}}%
\pgfpathlineto{\pgfqpoint{1.100776in}{1.125161in}}%
\pgfpathlineto{\pgfqpoint{1.107477in}{1.122854in}}%
\pgfpathlineto{\pgfqpoint{1.108594in}{1.120284in}}%
\pgfpathlineto{\pgfqpoint{1.109339in}{1.086719in}}%
\pgfpathlineto{\pgfqpoint{1.115668in}{0.681007in}}%
\pgfpathlineto{\pgfqpoint{1.118274in}{0.653460in}}%
\pgfpathlineto{\pgfqpoint{1.121624in}{0.652532in}}%
\pgfpathlineto{\pgfqpoint{1.125161in}{0.648987in}}%
\pgfpathlineto{\pgfqpoint{1.128511in}{0.649941in}}%
\pgfpathlineto{\pgfqpoint{1.130559in}{0.654846in}}%
\pgfpathlineto{\pgfqpoint{1.132793in}{0.658237in}}%
\pgfpathlineto{\pgfqpoint{1.135399in}{0.665385in}}%
\pgfpathlineto{\pgfqpoint{1.137074in}{0.673265in}}%
\pgfpathlineto{\pgfqpoint{1.138563in}{0.674248in}}%
\pgfpathlineto{\pgfqpoint{1.142100in}{0.689534in}}%
\pgfpathlineto{\pgfqpoint{1.146195in}{0.694786in}}%
\pgfpathlineto{\pgfqpoint{1.146753in}{0.694786in}}%
\pgfpathlineto{\pgfqpoint{1.148801in}{0.700578in}}%
\pgfpathlineto{\pgfqpoint{1.151407in}{0.699871in}}%
\pgfpathlineto{\pgfqpoint{1.152710in}{0.699349in}}%
\pgfpathlineto{\pgfqpoint{1.154571in}{0.699862in}}%
\pgfpathlineto{\pgfqpoint{1.159597in}{0.714769in}}%
\pgfpathlineto{\pgfqpoint{1.160714in}{0.713945in}}%
\pgfpathlineto{\pgfqpoint{1.162762in}{0.712194in}}%
\pgfpathlineto{\pgfqpoint{1.164251in}{0.713155in}}%
\pgfpathlineto{\pgfqpoint{1.168532in}{0.724496in}}%
\pgfpathlineto{\pgfqpoint{1.170580in}{0.740333in}}%
\pgfpathlineto{\pgfqpoint{1.172069in}{0.909058in}}%
\pgfpathlineto{\pgfqpoint{1.175606in}{0.912118in}}%
\pgfpathlineto{\pgfqpoint{1.176350in}{0.913165in}}%
\pgfpathlineto{\pgfqpoint{1.179887in}{0.927998in}}%
\pgfpathlineto{\pgfqpoint{1.181190in}{0.930553in}}%
\pgfpathlineto{\pgfqpoint{1.183237in}{0.930553in}}%
\pgfpathlineto{\pgfqpoint{1.186588in}{0.964434in}}%
\pgfpathlineto{\pgfqpoint{1.189194in}{1.041894in}}%
\pgfpathlineto{\pgfqpoint{1.192917in}{1.084144in}}%
\pgfpathlineto{\pgfqpoint{1.194034in}{1.090437in}}%
\pgfpathlineto{\pgfqpoint{1.196454in}{1.228237in}}%
\pgfpathlineto{\pgfqpoint{1.201293in}{1.398358in}}%
\pgfpathlineto{\pgfqpoint{1.206505in}{1.507716in}}%
\pgfpathlineto{\pgfqpoint{1.211345in}{1.568954in}}%
\pgfpathlineto{\pgfqpoint{1.213579in}{1.577812in}}%
\pgfpathlineto{\pgfqpoint{1.214882in}{1.578697in}}%
\pgfpathlineto{\pgfqpoint{1.217674in}{1.581390in}}%
\pgfpathlineto{\pgfqpoint{1.225306in}{1.580369in}}%
\pgfpathlineto{\pgfqpoint{1.226981in}{1.576410in}}%
\pgfpathlineto{\pgfqpoint{1.230890in}{1.519377in}}%
\pgfpathlineto{\pgfqpoint{1.233496in}{1.515023in}}%
\pgfpathlineto{\pgfqpoint{1.234241in}{1.514613in}}%
\pgfpathlineto{\pgfqpoint{1.235544in}{1.505815in}}%
\pgfpathlineto{\pgfqpoint{1.239639in}{1.504729in}}%
\pgfpathlineto{\pgfqpoint{1.246712in}{1.501106in}}%
\pgfpathlineto{\pgfqpoint{1.248015in}{1.502056in}}%
\pgfpathlineto{\pgfqpoint{1.249877in}{1.502772in}}%
\pgfpathlineto{\pgfqpoint{1.252110in}{1.501754in}}%
\pgfpathlineto{\pgfqpoint{1.255275in}{1.494494in}}%
\pgfpathlineto{\pgfqpoint{1.257509in}{1.490891in}}%
\pgfpathlineto{\pgfqpoint{1.258439in}{1.490087in}}%
\pgfpathlineto{\pgfqpoint{1.259556in}{1.485262in}}%
\pgfpathlineto{\pgfqpoint{1.260859in}{1.452835in}}%
\pgfpathlineto{\pgfqpoint{1.263279in}{1.418585in}}%
\pgfpathlineto{\pgfqpoint{1.264954in}{1.414082in}}%
\pgfpathlineto{\pgfqpoint{1.267188in}{1.412468in}}%
\pgfpathlineto{\pgfqpoint{1.268491in}{1.398629in}}%
\pgfpathlineto{\pgfqpoint{1.273889in}{1.396899in}}%
\pgfpathlineto{\pgfqpoint{1.283755in}{1.392345in}}%
\pgfpathlineto{\pgfqpoint{1.288781in}{1.380666in}}%
\pgfpathlineto{\pgfqpoint{1.300694in}{1.345557in}}%
\pgfpathlineto{\pgfqpoint{1.304417in}{1.338112in}}%
\pgfpathlineto{\pgfqpoint{1.312421in}{1.322447in}}%
\pgfpathlineto{\pgfqpoint{1.312979in}{1.321485in}}%
\pgfpathlineto{\pgfqpoint{1.314655in}{1.315790in}}%
\pgfpathlineto{\pgfqpoint{1.317261in}{1.292992in}}%
\pgfpathlineto{\pgfqpoint{1.319867in}{1.286382in}}%
\pgfpathlineto{\pgfqpoint{1.321356in}{1.285321in}}%
\pgfpathlineto{\pgfqpoint{1.321914in}{1.284969in}}%
\pgfpathlineto{\pgfqpoint{1.323217in}{1.273611in}}%
\pgfpathlineto{\pgfqpoint{1.327498in}{1.271587in}}%
\pgfpathlineto{\pgfqpoint{1.333269in}{1.267846in}}%
\pgfpathlineto{\pgfqpoint{1.334572in}{1.266744in}}%
\pgfpathlineto{\pgfqpoint{1.359515in}{1.238817in}}%
\pgfpathlineto{\pgfqpoint{1.363424in}{1.236955in}}%
\pgfpathlineto{\pgfqpoint{1.372545in}{1.237955in}}%
\pgfpathlineto{\pgfqpoint{1.376268in}{1.243514in}}%
\pgfpathlineto{\pgfqpoint{1.380549in}{1.254063in}}%
\pgfpathlineto{\pgfqpoint{1.385947in}{1.273997in}}%
\pgfpathlineto{\pgfqpoint{1.391904in}{1.304701in}}%
\pgfpathlineto{\pgfqpoint{1.398233in}{1.348723in}}%
\pgfpathlineto{\pgfqpoint{1.402514in}{1.375997in}}%
\pgfpathlineto{\pgfqpoint{1.403445in}{1.375997in}}%
\pgfpathlineto{\pgfqpoint{1.413869in}{1.475251in}}%
\pgfpathlineto{\pgfqpoint{1.415358in}{1.496211in}}%
\pgfpathlineto{\pgfqpoint{1.422990in}{1.495887in}}%
\pgfpathlineto{\pgfqpoint{1.425968in}{1.393636in}}%
\pgfpathlineto{\pgfqpoint{1.427457in}{1.359508in}}%
\pgfpathlineto{\pgfqpoint{1.429319in}{1.325698in}}%
\pgfpathlineto{\pgfqpoint{1.430994in}{1.299407in}}%
\pgfpathlineto{\pgfqpoint{1.439184in}{1.156979in}}%
\pgfpathlineto{\pgfqpoint{1.446072in}{1.105582in}}%
\pgfpathlineto{\pgfqpoint{1.447933in}{1.095829in}}%
\pgfpathlineto{\pgfqpoint{1.449422in}{1.020905in}}%
\pgfpathlineto{\pgfqpoint{1.452214in}{0.885027in}}%
\pgfpathlineto{\pgfqpoint{1.455379in}{0.866033in}}%
\pgfpathlineto{\pgfqpoint{1.461708in}{0.866624in}}%
\pgfpathlineto{\pgfqpoint{1.466733in}{0.890606in}}%
\pgfpathlineto{\pgfqpoint{1.485720in}{1.029447in}}%
\pgfpathlineto{\pgfqpoint{1.491490in}{1.054327in}}%
\pgfpathlineto{\pgfqpoint{1.496516in}{1.066999in}}%
\pgfpathlineto{\pgfqpoint{1.500798in}{1.072434in}}%
\pgfpathlineto{\pgfqpoint{1.503031in}{1.072547in}}%
\pgfpathlineto{\pgfqpoint{1.511594in}{1.071522in}}%
\pgfpathlineto{\pgfqpoint{1.516992in}{1.065403in}}%
\pgfpathlineto{\pgfqpoint{1.524252in}{1.053531in}}%
\pgfpathlineto{\pgfqpoint{1.533187in}{1.034294in}}%
\pgfpathlineto{\pgfqpoint{1.543424in}{1.006745in}}%
\pgfpathlineto{\pgfqpoint{1.550498in}{0.988578in}}%
\pgfpathlineto{\pgfqpoint{1.552732in}{0.987538in}}%
\pgfpathlineto{\pgfqpoint{1.563156in}{0.979166in}}%
\pgfpathlineto{\pgfqpoint{1.564831in}{0.956622in}}%
\pgfpathlineto{\pgfqpoint{1.569298in}{0.951949in}}%
\pgfpathlineto{\pgfqpoint{1.575255in}{0.940696in}}%
\pgfpathlineto{\pgfqpoint{1.581956in}{0.940696in}}%
\pgfpathlineto{\pgfqpoint{1.587168in}{0.994692in}}%
\pgfpathlineto{\pgfqpoint{1.590332in}{1.030734in}}%
\pgfpathlineto{\pgfqpoint{1.598709in}{1.030734in}}%
\pgfpathlineto{\pgfqpoint{1.599081in}{1.014729in}}%
\pgfpathlineto{\pgfqpoint{1.600384in}{0.990979in}}%
\pgfpathlineto{\pgfqpoint{1.601501in}{0.990024in}}%
\pgfpathlineto{\pgfqpoint{1.602804in}{0.982467in}}%
\pgfpathlineto{\pgfqpoint{1.605410in}{0.947747in}}%
\pgfpathlineto{\pgfqpoint{1.608202in}{0.920279in}}%
\pgfpathlineto{\pgfqpoint{1.616392in}{0.921081in}}%
\pgfpathlineto{\pgfqpoint{1.617509in}{0.941599in}}%
\pgfpathlineto{\pgfqpoint{1.622721in}{1.034390in}}%
\pgfpathlineto{\pgfqpoint{1.624583in}{1.044358in}}%
\pgfpathlineto{\pgfqpoint{1.625141in}{1.044358in}}%
\pgfpathlineto{\pgfqpoint{1.630912in}{1.101633in}}%
\pgfpathlineto{\pgfqpoint{1.636310in}{1.113368in}}%
\pgfpathlineto{\pgfqpoint{1.640219in}{1.114201in}}%
\pgfpathlineto{\pgfqpoint{1.645989in}{1.114167in}}%
\pgfpathlineto{\pgfqpoint{1.651201in}{1.112108in}}%
\pgfpathlineto{\pgfqpoint{1.651574in}{1.111262in}}%
\pgfpathlineto{\pgfqpoint{1.652318in}{1.082377in}}%
\pgfpathlineto{\pgfqpoint{1.656413in}{0.876705in}}%
\pgfpathlineto{\pgfqpoint{1.664045in}{0.876705in}}%
\pgfpathlineto{\pgfqpoint{1.666279in}{0.958780in}}%
\pgfpathlineto{\pgfqpoint{1.668513in}{1.054616in}}%
\pgfpathlineto{\pgfqpoint{1.670560in}{1.077343in}}%
\pgfpathlineto{\pgfqpoint{1.673538in}{1.093609in}}%
\pgfpathlineto{\pgfqpoint{1.675214in}{1.137538in}}%
\pgfpathlineto{\pgfqpoint{1.678378in}{1.170275in}}%
\pgfpathlineto{\pgfqpoint{1.682659in}{1.200093in}}%
\pgfpathlineto{\pgfqpoint{1.684335in}{1.202995in}}%
\pgfpathlineto{\pgfqpoint{1.692525in}{1.202995in}}%
\pgfpathlineto{\pgfqpoint{1.694200in}{1.184065in}}%
\pgfpathlineto{\pgfqpoint{1.696992in}{1.129385in}}%
\pgfpathlineto{\pgfqpoint{1.698854in}{1.103754in}}%
\pgfpathlineto{\pgfqpoint{1.704624in}{1.101366in}}%
\pgfpathlineto{\pgfqpoint{1.706113in}{1.099650in}}%
\pgfpathlineto{\pgfqpoint{1.707603in}{1.094367in}}%
\pgfpathlineto{\pgfqpoint{1.708719in}{1.025176in}}%
\pgfpathlineto{\pgfqpoint{1.710395in}{0.969143in}}%
\pgfpathlineto{\pgfqpoint{1.718585in}{0.969143in}}%
\pgfpathlineto{\pgfqpoint{1.721936in}{1.060376in}}%
\pgfpathlineto{\pgfqpoint{1.726217in}{1.111149in}}%
\pgfpathlineto{\pgfqpoint{1.734221in}{1.113111in}}%
\pgfpathlineto{\pgfqpoint{1.735152in}{1.116867in}}%
\pgfpathlineto{\pgfqpoint{1.737758in}{1.143790in}}%
\pgfpathlineto{\pgfqpoint{1.743528in}{1.200328in}}%
\pgfpathlineto{\pgfqpoint{1.746879in}{1.217293in}}%
\pgfpathlineto{\pgfqpoint{1.748740in}{1.219739in}}%
\pgfpathlineto{\pgfqpoint{1.757117in}{1.219413in}}%
\pgfpathlineto{\pgfqpoint{1.762515in}{1.196953in}}%
\pgfpathlineto{\pgfqpoint{1.766796in}{1.187855in}}%
\pgfpathlineto{\pgfqpoint{1.770519in}{1.184476in}}%
\pgfpathlineto{\pgfqpoint{1.774056in}{1.184005in}}%
\pgfpathlineto{\pgfqpoint{1.784480in}{1.182642in}}%
\pgfpathlineto{\pgfqpoint{1.787086in}{1.178936in}}%
\pgfpathlineto{\pgfqpoint{1.790622in}{1.169314in}}%
\pgfpathlineto{\pgfqpoint{1.794531in}{1.151958in}}%
\pgfpathlineto{\pgfqpoint{1.801233in}{1.118454in}}%
\pgfpathlineto{\pgfqpoint{1.803466in}{1.117282in}}%
\pgfpathlineto{\pgfqpoint{1.807375in}{1.115894in}}%
\pgfpathlineto{\pgfqpoint{1.808492in}{1.096064in}}%
\pgfpathlineto{\pgfqpoint{1.810726in}{1.061855in}}%
\pgfpathlineto{\pgfqpoint{1.818916in}{1.061855in}}%
\pgfpathlineto{\pgfqpoint{1.824314in}{1.165758in}}%
\pgfpathlineto{\pgfqpoint{1.826362in}{1.185329in}}%
\pgfpathlineto{\pgfqpoint{1.828037in}{1.185631in}}%
\pgfpathlineto{\pgfqpoint{1.834738in}{1.186971in}}%
\pgfpathlineto{\pgfqpoint{1.836227in}{1.193562in}}%
\pgfpathlineto{\pgfqpoint{1.837903in}{1.211133in}}%
\pgfpathlineto{\pgfqpoint{1.840323in}{1.263944in}}%
\pgfpathlineto{\pgfqpoint{1.843115in}{1.383126in}}%
\pgfpathlineto{\pgfqpoint{1.846465in}{1.529444in}}%
\pgfpathlineto{\pgfqpoint{1.854469in}{1.529444in}}%
\pgfpathlineto{\pgfqpoint{1.855959in}{1.492844in}}%
\pgfpathlineto{\pgfqpoint{1.863963in}{1.492844in}}%
\pgfpathlineto{\pgfqpoint{1.870106in}{1.569083in}}%
\pgfpathlineto{\pgfqpoint{1.871409in}{1.571435in}}%
\pgfpathlineto{\pgfqpoint{1.877179in}{1.572498in}}%
\pgfpathlineto{\pgfqpoint{1.879413in}{1.578702in}}%
\pgfpathlineto{\pgfqpoint{1.880716in}{1.578702in}}%
\pgfpathlineto{\pgfqpoint{1.881460in}{1.562825in}}%
\pgfpathlineto{\pgfqpoint{1.883508in}{1.541240in}}%
\pgfpathlineto{\pgfqpoint{1.884997in}{1.538651in}}%
\pgfpathlineto{\pgfqpoint{1.887789in}{1.535218in}}%
\pgfpathlineto{\pgfqpoint{1.890209in}{1.472655in}}%
\pgfpathlineto{\pgfqpoint{1.897841in}{1.473703in}}%
\pgfpathlineto{\pgfqpoint{1.900075in}{1.476726in}}%
\pgfpathlineto{\pgfqpoint{1.902494in}{1.477837in}}%
\pgfpathlineto{\pgfqpoint{1.911243in}{1.476951in}}%
\pgfpathlineto{\pgfqpoint{1.911802in}{1.475291in}}%
\pgfpathlineto{\pgfqpoint{1.916083in}{1.428672in}}%
\pgfpathlineto{\pgfqpoint{1.919433in}{1.426160in}}%
\pgfpathlineto{\pgfqpoint{1.920736in}{1.379230in}}%
\pgfpathlineto{\pgfqpoint{1.940281in}{1.378283in}}%
\pgfpathlineto{\pgfqpoint{1.950705in}{1.360125in}}%
\pgfpathlineto{\pgfqpoint{1.966528in}{1.319282in}}%
\pgfpathlineto{\pgfqpoint{1.968575in}{1.318522in}}%
\pgfpathlineto{\pgfqpoint{1.970064in}{1.311266in}}%
\pgfpathlineto{\pgfqpoint{1.974718in}{1.283398in}}%
\pgfpathlineto{\pgfqpoint{1.977510in}{1.279120in}}%
\pgfpathlineto{\pgfqpoint{1.978813in}{1.270175in}}%
\pgfpathlineto{\pgfqpoint{1.984956in}{1.265590in}}%
\pgfpathlineto{\pgfqpoint{1.988306in}{1.263205in}}%
\pgfpathlineto{\pgfqpoint{1.989796in}{1.262225in}}%
\pgfpathlineto{\pgfqpoint{2.010830in}{1.236370in}}%
\pgfpathlineto{\pgfqpoint{2.013436in}{1.235627in}}%
\pgfpathlineto{\pgfqpoint{2.022184in}{1.236727in}}%
\pgfpathlineto{\pgfqpoint{2.025907in}{1.242941in}}%
\pgfpathlineto{\pgfqpoint{2.030375in}{1.254996in}}%
\pgfpathlineto{\pgfqpoint{2.035959in}{1.277261in}}%
\pgfpathlineto{\pgfqpoint{2.042288in}{1.312225in}}%
\pgfpathlineto{\pgfqpoint{2.048617in}{1.359476in}}%
\pgfpathlineto{\pgfqpoint{2.051037in}{1.370834in}}%
\pgfpathlineto{\pgfqpoint{2.051967in}{1.370834in}}%
\pgfpathlineto{\pgfqpoint{2.065183in}{1.500177in}}%
\pgfpathlineto{\pgfqpoint{2.066486in}{1.508944in}}%
\pgfpathlineto{\pgfqpoint{2.071698in}{1.508944in}}%
\pgfpathlineto{\pgfqpoint{2.075049in}{1.434589in}}%
\pgfpathlineto{\pgfqpoint{2.086404in}{1.207167in}}%
\pgfpathlineto{\pgfqpoint{2.086776in}{1.203711in}}%
\pgfpathlineto{\pgfqpoint{2.090313in}{1.141853in}}%
\pgfpathlineto{\pgfqpoint{2.094222in}{1.103675in}}%
\pgfpathlineto{\pgfqpoint{2.096642in}{1.096621in}}%
\pgfpathlineto{\pgfqpoint{2.098689in}{0.958793in}}%
\pgfpathlineto{\pgfqpoint{2.100923in}{0.879217in}}%
\pgfpathlineto{\pgfqpoint{2.103715in}{0.866717in}}%
\pgfpathlineto{\pgfqpoint{2.110789in}{0.867624in}}%
\pgfpathlineto{\pgfqpoint{2.115814in}{0.889893in}}%
\pgfpathlineto{\pgfqpoint{2.138152in}{1.045695in}}%
\pgfpathlineto{\pgfqpoint{2.143550in}{1.065668in}}%
\pgfpathlineto{\pgfqpoint{2.148203in}{1.075437in}}%
\pgfpathlineto{\pgfqpoint{2.150995in}{1.077593in}}%
\pgfpathlineto{\pgfqpoint{2.160303in}{1.076651in}}%
\pgfpathlineto{\pgfqpoint{2.165328in}{1.069744in}}%
\pgfpathlineto{\pgfqpoint{2.173891in}{1.053656in}}%
\pgfpathlineto{\pgfqpoint{2.186921in}{1.024290in}}%
\pgfpathlineto{\pgfqpoint{2.210375in}{0.963195in}}%
\pgfpathlineto{\pgfqpoint{2.213353in}{0.943833in}}%
\pgfpathlineto{\pgfqpoint{2.215215in}{0.942742in}}%
\pgfpathlineto{\pgfqpoint{2.219124in}{0.936038in}}%
\pgfpathlineto{\pgfqpoint{2.220427in}{0.932355in}}%
\pgfpathlineto{\pgfqpoint{2.220799in}{0.933780in}}%
\pgfpathlineto{\pgfqpoint{2.223591in}{0.939430in}}%
\pgfpathlineto{\pgfqpoint{2.228617in}{0.939973in}}%
\pgfpathlineto{\pgfqpoint{2.229548in}{0.946063in}}%
\pgfpathlineto{\pgfqpoint{2.235132in}{1.013658in}}%
\pgfpathlineto{\pgfqpoint{2.237552in}{1.032314in}}%
\pgfpathlineto{\pgfqpoint{2.245742in}{1.032314in}}%
\pgfpathlineto{\pgfqpoint{2.248162in}{0.988316in}}%
\pgfpathlineto{\pgfqpoint{2.249093in}{0.986970in}}%
\pgfpathlineto{\pgfqpoint{2.250582in}{0.975689in}}%
\pgfpathlineto{\pgfqpoint{2.255980in}{0.926785in}}%
\pgfpathlineto{\pgfqpoint{2.263240in}{0.926785in}}%
\pgfpathlineto{\pgfqpoint{2.264729in}{0.943420in}}%
\pgfpathlineto{\pgfqpoint{2.270313in}{1.037791in}}%
\pgfpathlineto{\pgfqpoint{2.272175in}{1.045952in}}%
\pgfpathlineto{\pgfqpoint{2.272547in}{1.045952in}}%
\pgfpathlineto{\pgfqpoint{2.279620in}{1.104235in}}%
\pgfpathlineto{\pgfqpoint{2.287066in}{1.104093in}}%
\pgfpathlineto{\pgfqpoint{2.288183in}{1.086842in}}%
\pgfpathlineto{\pgfqpoint{2.290975in}{1.049102in}}%
\pgfpathlineto{\pgfqpoint{2.291161in}{1.049102in}}%
\pgfpathlineto{\pgfqpoint{2.296746in}{0.873430in}}%
\pgfpathlineto{\pgfqpoint{2.304377in}{0.873968in}}%
\pgfpathlineto{\pgfqpoint{2.306239in}{0.910052in}}%
\pgfpathlineto{\pgfqpoint{2.308845in}{1.042108in}}%
\pgfpathlineto{\pgfqpoint{2.310892in}{1.073804in}}%
\pgfpathlineto{\pgfqpoint{2.314615in}{1.107510in}}%
\pgfpathlineto{\pgfqpoint{2.316291in}{1.136063in}}%
\pgfpathlineto{\pgfqpoint{2.319827in}{1.169583in}}%
\pgfpathlineto{\pgfqpoint{2.324109in}{1.194647in}}%
\pgfpathlineto{\pgfqpoint{2.327273in}{1.203932in}}%
\pgfpathlineto{\pgfqpoint{2.336022in}{1.203795in}}%
\pgfpathlineto{\pgfqpoint{2.338255in}{1.186052in}}%
\pgfpathlineto{\pgfqpoint{2.341048in}{1.145676in}}%
\pgfpathlineto{\pgfqpoint{2.343654in}{1.103487in}}%
\pgfpathlineto{\pgfqpoint{2.349982in}{1.100113in}}%
\pgfpathlineto{\pgfqpoint{2.351658in}{1.095075in}}%
\pgfpathlineto{\pgfqpoint{2.352588in}{1.056189in}}%
\pgfpathlineto{\pgfqpoint{2.355008in}{0.963613in}}%
\pgfpathlineto{\pgfqpoint{2.363199in}{0.963613in}}%
\pgfpathlineto{\pgfqpoint{2.366177in}{1.047039in}}%
\pgfpathlineto{\pgfqpoint{2.371575in}{1.109453in}}%
\pgfpathlineto{\pgfqpoint{2.373064in}{1.110834in}}%
\pgfpathlineto{\pgfqpoint{2.379021in}{1.112335in}}%
\pgfpathlineto{\pgfqpoint{2.380510in}{1.114402in}}%
\pgfpathlineto{\pgfqpoint{2.382185in}{1.128135in}}%
\pgfpathlineto{\pgfqpoint{2.390934in}{1.211146in}}%
\pgfpathlineto{\pgfqpoint{2.393540in}{1.219633in}}%
\pgfpathlineto{\pgfqpoint{2.403033in}{1.218759in}}%
\pgfpathlineto{\pgfqpoint{2.408431in}{1.196235in}}%
\pgfpathlineto{\pgfqpoint{2.412713in}{1.186935in}}%
\pgfpathlineto{\pgfqpoint{2.416808in}{1.182782in}}%
\pgfpathlineto{\pgfqpoint{2.420903in}{1.181497in}}%
\pgfpathlineto{\pgfqpoint{2.441937in}{1.182604in}}%
\pgfpathlineto{\pgfqpoint{2.442496in}{1.182628in}}%
\pgfpathlineto{\pgfqpoint{2.442496in}{1.182628in}}%
\pgfusepath{stroke}%
\end{pgfscope}%
\begin{pgfscope}%
\pgfpathrectangle{\pgfqpoint{0.464911in}{0.456793in}}{\pgfqpoint{1.977584in}{1.453311in}}%
\pgfusepath{clip}%
\pgfsetrectcap%
\pgfsetroundjoin%
\pgfsetlinewidth{0.803000pt}%
\definecolor{currentstroke}{rgb}{1.000000,0.498039,0.054902}%
\pgfsetstrokecolor{currentstroke}%
\pgfsetdash{}{0pt}%
\pgfpathmoveto{\pgfqpoint{0.464911in}{1.183560in}}%
\pgfpathlineto{\pgfqpoint{0.516659in}{1.183778in}}%
\pgfpathlineto{\pgfqpoint{0.521313in}{1.184402in}}%
\pgfpathlineto{\pgfqpoint{0.524105in}{1.183584in}}%
\pgfpathlineto{\pgfqpoint{0.525594in}{1.182968in}}%
\pgfpathlineto{\pgfqpoint{0.551654in}{1.181863in}}%
\pgfpathlineto{\pgfqpoint{0.559100in}{1.179662in}}%
\pgfpathlineto{\pgfqpoint{0.560961in}{1.179908in}}%
\pgfpathlineto{\pgfqpoint{0.589441in}{1.184261in}}%
\pgfpathlineto{\pgfqpoint{0.603029in}{1.184342in}}%
\pgfpathlineto{\pgfqpoint{0.605635in}{1.182645in}}%
\pgfpathlineto{\pgfqpoint{0.609544in}{1.183728in}}%
\pgfpathlineto{\pgfqpoint{0.659431in}{1.184790in}}%
\pgfpathlineto{\pgfqpoint{0.663526in}{1.185659in}}%
\pgfpathlineto{\pgfqpoint{0.703174in}{1.184564in}}%
\pgfpathlineto{\pgfqpoint{0.721975in}{1.181826in}}%
\pgfpathlineto{\pgfqpoint{0.727001in}{1.182920in}}%
\pgfpathlineto{\pgfqpoint{0.730537in}{1.185301in}}%
\pgfpathlineto{\pgfqpoint{0.732585in}{1.188236in}}%
\pgfpathlineto{\pgfqpoint{0.740961in}{1.190373in}}%
\pgfpathlineto{\pgfqpoint{0.742451in}{1.190742in}}%
\pgfpathlineto{\pgfqpoint{0.751385in}{1.189749in}}%
\pgfpathlineto{\pgfqpoint{0.754178in}{1.184365in}}%
\pgfpathlineto{\pgfqpoint{0.757900in}{1.171713in}}%
\pgfpathlineto{\pgfqpoint{0.762926in}{1.145611in}}%
\pgfpathlineto{\pgfqpoint{0.770000in}{1.095605in}}%
\pgfpathlineto{\pgfqpoint{0.780052in}{1.007064in}}%
\pgfpathlineto{\pgfqpoint{0.804809in}{0.777045in}}%
\pgfpathlineto{\pgfqpoint{0.810393in}{0.749929in}}%
\pgfpathlineto{\pgfqpoint{0.815419in}{0.735356in}}%
\pgfpathlineto{\pgfqpoint{0.818769in}{0.730503in}}%
\pgfpathlineto{\pgfqpoint{0.821561in}{0.727439in}}%
\pgfpathlineto{\pgfqpoint{0.822864in}{0.726817in}}%
\pgfpathlineto{\pgfqpoint{0.823981in}{0.725713in}}%
\pgfpathlineto{\pgfqpoint{0.825657in}{0.725019in}}%
\pgfpathlineto{\pgfqpoint{0.829938in}{0.723467in}}%
\pgfpathlineto{\pgfqpoint{0.832172in}{0.722087in}}%
\pgfpathlineto{\pgfqpoint{0.834219in}{0.718198in}}%
\pgfpathlineto{\pgfqpoint{0.837197in}{0.717451in}}%
\pgfpathlineto{\pgfqpoint{0.871634in}{0.715835in}}%
\pgfpathlineto{\pgfqpoint{0.885967in}{0.715962in}}%
\pgfpathlineto{\pgfqpoint{0.904395in}{0.714512in}}%
\pgfpathlineto{\pgfqpoint{0.907560in}{0.716192in}}%
\pgfpathlineto{\pgfqpoint{0.910352in}{0.720162in}}%
\pgfpathlineto{\pgfqpoint{0.913330in}{0.728666in}}%
\pgfpathlineto{\pgfqpoint{0.916867in}{0.746186in}}%
\pgfpathlineto{\pgfqpoint{0.921148in}{0.780208in}}%
\pgfpathlineto{\pgfqpoint{0.926360in}{0.841492in}}%
\pgfpathlineto{\pgfqpoint{0.933061in}{0.949774in}}%
\pgfpathlineto{\pgfqpoint{0.942368in}{1.140527in}}%
\pgfpathlineto{\pgfqpoint{0.957818in}{1.456744in}}%
\pgfpathlineto{\pgfqpoint{0.965636in}{1.569373in}}%
\pgfpathlineto{\pgfqpoint{0.970848in}{1.616747in}}%
\pgfpathlineto{\pgfqpoint{0.975129in}{1.637234in}}%
\pgfpathlineto{\pgfqpoint{0.978852in}{1.645303in}}%
\pgfpathlineto{\pgfqpoint{0.982203in}{1.647994in}}%
\pgfpathlineto{\pgfqpoint{0.995978in}{1.646892in}}%
\pgfpathlineto{\pgfqpoint{1.039163in}{1.636611in}}%
\pgfpathlineto{\pgfqpoint{1.054054in}{1.637710in}}%
\pgfpathlineto{\pgfqpoint{1.063361in}{1.637908in}}%
\pgfpathlineto{\pgfqpoint{1.066898in}{1.636838in}}%
\pgfpathlineto{\pgfqpoint{1.070249in}{1.633132in}}%
\pgfpathlineto{\pgfqpoint{1.074158in}{1.625030in}}%
\pgfpathlineto{\pgfqpoint{1.078625in}{1.610169in}}%
\pgfpathlineto{\pgfqpoint{1.084023in}{1.583706in}}%
\pgfpathlineto{\pgfqpoint{1.090724in}{1.538719in}}%
\pgfpathlineto{\pgfqpoint{1.100218in}{1.458124in}}%
\pgfpathlineto{\pgfqpoint{1.120135in}{1.289092in}}%
\pgfpathlineto{\pgfqpoint{1.125905in}{1.252839in}}%
\pgfpathlineto{\pgfqpoint{1.131304in}{1.231859in}}%
\pgfpathlineto{\pgfqpoint{1.132420in}{1.228406in}}%
\pgfpathlineto{\pgfqpoint{1.136143in}{1.218109in}}%
\pgfpathlineto{\pgfqpoint{1.140983in}{1.211982in}}%
\pgfpathlineto{\pgfqpoint{1.141728in}{1.210690in}}%
\pgfpathlineto{\pgfqpoint{1.143403in}{1.206219in}}%
\pgfpathlineto{\pgfqpoint{1.147684in}{1.203103in}}%
\pgfpathlineto{\pgfqpoint{1.150104in}{1.199868in}}%
\pgfpathlineto{\pgfqpoint{1.157364in}{1.195995in}}%
\pgfpathlineto{\pgfqpoint{1.160528in}{1.193863in}}%
\pgfpathlineto{\pgfqpoint{1.167043in}{1.192766in}}%
\pgfpathlineto{\pgfqpoint{1.175047in}{1.189665in}}%
\pgfpathlineto{\pgfqpoint{1.178398in}{1.187334in}}%
\pgfpathlineto{\pgfqpoint{1.181190in}{1.187246in}}%
\pgfpathlineto{\pgfqpoint{1.184168in}{1.187381in}}%
\pgfpathlineto{\pgfqpoint{1.195895in}{1.188406in}}%
\pgfpathlineto{\pgfqpoint{1.198687in}{1.189708in}}%
\pgfpathlineto{\pgfqpoint{1.202038in}{1.191436in}}%
\pgfpathlineto{\pgfqpoint{1.208367in}{1.198557in}}%
\pgfpathlineto{\pgfqpoint{1.214510in}{1.208293in}}%
\pgfpathlineto{\pgfqpoint{1.224003in}{1.228310in}}%
\pgfpathlineto{\pgfqpoint{1.229401in}{1.246615in}}%
\pgfpathlineto{\pgfqpoint{1.233682in}{1.256288in}}%
\pgfpathlineto{\pgfqpoint{1.234613in}{1.267281in}}%
\pgfpathlineto{\pgfqpoint{1.235916in}{1.274000in}}%
\pgfpathlineto{\pgfqpoint{1.238336in}{1.281028in}}%
\pgfpathlineto{\pgfqpoint{1.244665in}{1.306458in}}%
\pgfpathlineto{\pgfqpoint{1.249318in}{1.324051in}}%
\pgfpathlineto{\pgfqpoint{1.262907in}{1.382476in}}%
\pgfpathlineto{\pgfqpoint{1.267188in}{1.395569in}}%
\pgfpathlineto{\pgfqpoint{1.269236in}{1.412057in}}%
\pgfpathlineto{\pgfqpoint{1.271283in}{1.418224in}}%
\pgfpathlineto{\pgfqpoint{1.276867in}{1.439607in}}%
\pgfpathlineto{\pgfqpoint{1.283382in}{1.458037in}}%
\pgfpathlineto{\pgfqpoint{1.292690in}{1.488658in}}%
\pgfpathlineto{\pgfqpoint{1.305906in}{1.526348in}}%
\pgfpathlineto{\pgfqpoint{1.317819in}{1.553811in}}%
\pgfpathlineto{\pgfqpoint{1.321914in}{1.560520in}}%
\pgfpathlineto{\pgfqpoint{1.323589in}{1.567589in}}%
\pgfpathlineto{\pgfqpoint{1.325823in}{1.570663in}}%
\pgfpathlineto{\pgfqpoint{1.332338in}{1.582789in}}%
\pgfpathlineto{\pgfqpoint{1.336806in}{1.586997in}}%
\pgfpathlineto{\pgfqpoint{1.348160in}{1.601713in}}%
\pgfpathlineto{\pgfqpoint{1.360073in}{1.614007in}}%
\pgfpathlineto{\pgfqpoint{1.362307in}{1.614429in}}%
\pgfpathlineto{\pgfqpoint{1.370870in}{1.613372in}}%
\pgfpathlineto{\pgfqpoint{1.374406in}{1.606828in}}%
\pgfpathlineto{\pgfqpoint{1.378688in}{1.593812in}}%
\pgfpathlineto{\pgfqpoint{1.383900in}{1.570318in}}%
\pgfpathlineto{\pgfqpoint{1.390415in}{1.529509in}}%
\pgfpathlineto{\pgfqpoint{1.398605in}{1.462126in}}%
\pgfpathlineto{\pgfqpoint{1.413124in}{1.321125in}}%
\pgfpathlineto{\pgfqpoint{1.415358in}{1.299105in}}%
\pgfpathlineto{\pgfqpoint{1.417592in}{1.289467in}}%
\pgfpathlineto{\pgfqpoint{1.420756in}{1.255974in}}%
\pgfpathlineto{\pgfqpoint{1.422804in}{1.240097in}}%
\pgfpathlineto{\pgfqpoint{1.428388in}{1.203917in}}%
\pgfpathlineto{\pgfqpoint{1.434345in}{1.176673in}}%
\pgfpathlineto{\pgfqpoint{1.436764in}{1.175561in}}%
\pgfpathlineto{\pgfqpoint{1.441604in}{1.165068in}}%
\pgfpathlineto{\pgfqpoint{1.445513in}{1.164977in}}%
\pgfpathlineto{\pgfqpoint{1.451656in}{1.165661in}}%
\pgfpathlineto{\pgfqpoint{1.453145in}{1.172222in}}%
\pgfpathlineto{\pgfqpoint{1.455006in}{1.175624in}}%
\pgfpathlineto{\pgfqpoint{1.457426in}{1.189953in}}%
\pgfpathlineto{\pgfqpoint{1.467478in}{1.259724in}}%
\pgfpathlineto{\pgfqpoint{1.480694in}{1.380100in}}%
\pgfpathlineto{\pgfqpoint{1.492607in}{1.481267in}}%
\pgfpathlineto{\pgfqpoint{1.500425in}{1.530711in}}%
\pgfpathlineto{\pgfqpoint{1.506010in}{1.553482in}}%
\pgfpathlineto{\pgfqpoint{1.510105in}{1.562050in}}%
\pgfpathlineto{\pgfqpoint{1.520156in}{1.561061in}}%
\pgfpathlineto{\pgfqpoint{1.525368in}{1.552866in}}%
\pgfpathlineto{\pgfqpoint{1.532814in}{1.536503in}}%
\pgfpathlineto{\pgfqpoint{1.542494in}{1.509410in}}%
\pgfpathlineto{\pgfqpoint{1.554965in}{1.467156in}}%
\pgfpathlineto{\pgfqpoint{1.590705in}{1.343076in}}%
\pgfpathlineto{\pgfqpoint{1.596289in}{1.334102in}}%
\pgfpathlineto{\pgfqpoint{1.601315in}{1.325567in}}%
\pgfpathlineto{\pgfqpoint{1.608761in}{1.308143in}}%
\pgfpathlineto{\pgfqpoint{1.609691in}{1.306166in}}%
\pgfpathlineto{\pgfqpoint{1.611925in}{1.295432in}}%
\pgfpathlineto{\pgfqpoint{1.613042in}{1.293849in}}%
\pgfpathlineto{\pgfqpoint{1.617509in}{1.280318in}}%
\pgfpathlineto{\pgfqpoint{1.621418in}{1.273043in}}%
\pgfpathlineto{\pgfqpoint{1.628678in}{1.265073in}}%
\pgfpathlineto{\pgfqpoint{1.634262in}{1.262902in}}%
\pgfpathlineto{\pgfqpoint{1.637613in}{1.260661in}}%
\pgfpathlineto{\pgfqpoint{1.645431in}{1.257492in}}%
\pgfpathlineto{\pgfqpoint{1.649340in}{1.256414in}}%
\pgfpathlineto{\pgfqpoint{1.652877in}{1.253823in}}%
\pgfpathlineto{\pgfqpoint{1.656227in}{1.246184in}}%
\pgfpathlineto{\pgfqpoint{1.663487in}{1.228613in}}%
\pgfpathlineto{\pgfqpoint{1.664045in}{1.226500in}}%
\pgfpathlineto{\pgfqpoint{1.665907in}{1.221720in}}%
\pgfpathlineto{\pgfqpoint{1.672608in}{1.220155in}}%
\pgfpathlineto{\pgfqpoint{1.675028in}{1.218064in}}%
\pgfpathlineto{\pgfqpoint{1.682473in}{1.217805in}}%
\pgfpathlineto{\pgfqpoint{1.698668in}{1.216716in}}%
\pgfpathlineto{\pgfqpoint{1.707789in}{1.211278in}}%
\pgfpathlineto{\pgfqpoint{1.710953in}{1.205888in}}%
\pgfpathlineto{\pgfqpoint{1.715234in}{1.195655in}}%
\pgfpathlineto{\pgfqpoint{1.716537in}{1.193256in}}%
\pgfpathlineto{\pgfqpoint{1.742225in}{1.188990in}}%
\pgfpathlineto{\pgfqpoint{1.773870in}{1.187886in}}%
\pgfpathlineto{\pgfqpoint{1.849630in}{1.180575in}}%
\pgfpathlineto{\pgfqpoint{1.854656in}{1.178610in}}%
\pgfpathlineto{\pgfqpoint{1.857076in}{1.172985in}}%
\pgfpathlineto{\pgfqpoint{1.867127in}{1.161842in}}%
\pgfpathlineto{\pgfqpoint{1.873642in}{1.147521in}}%
\pgfpathlineto{\pgfqpoint{1.877551in}{1.133653in}}%
\pgfpathlineto{\pgfqpoint{1.878854in}{1.130165in}}%
\pgfpathlineto{\pgfqpoint{1.879971in}{1.128817in}}%
\pgfpathlineto{\pgfqpoint{1.881088in}{1.126170in}}%
\pgfpathlineto{\pgfqpoint{1.887045in}{1.103811in}}%
\pgfpathlineto{\pgfqpoint{1.887789in}{1.100314in}}%
\pgfpathlineto{\pgfqpoint{1.895049in}{1.064144in}}%
\pgfpathlineto{\pgfqpoint{1.896538in}{1.060796in}}%
\pgfpathlineto{\pgfqpoint{1.899888in}{1.045359in}}%
\pgfpathlineto{\pgfqpoint{1.910871in}{0.996323in}}%
\pgfpathlineto{\pgfqpoint{1.912360in}{0.993722in}}%
\pgfpathlineto{\pgfqpoint{1.923715in}{0.938832in}}%
\pgfpathlineto{\pgfqpoint{1.925762in}{0.932726in}}%
\pgfpathlineto{\pgfqpoint{1.929113in}{0.929234in}}%
\pgfpathlineto{\pgfqpoint{1.933580in}{0.915651in}}%
\pgfpathlineto{\pgfqpoint{1.958896in}{0.840702in}}%
\pgfpathlineto{\pgfqpoint{1.970064in}{0.813902in}}%
\pgfpathlineto{\pgfqpoint{1.971926in}{0.809996in}}%
\pgfpathlineto{\pgfqpoint{1.975649in}{0.802320in}}%
\pgfpathlineto{\pgfqpoint{1.977882in}{0.798544in}}%
\pgfpathlineto{\pgfqpoint{1.983839in}{0.783019in}}%
\pgfpathlineto{\pgfqpoint{1.993891in}{0.771262in}}%
\pgfpathlineto{\pgfqpoint{2.006176in}{0.755302in}}%
\pgfpathlineto{\pgfqpoint{2.011016in}{0.751494in}}%
\pgfpathlineto{\pgfqpoint{2.020137in}{0.752512in}}%
\pgfpathlineto{\pgfqpoint{2.023674in}{0.759287in}}%
\pgfpathlineto{\pgfqpoint{2.027955in}{0.772647in}}%
\pgfpathlineto{\pgfqpoint{2.033167in}{0.796662in}}%
\pgfpathlineto{\pgfqpoint{2.039682in}{0.838062in}}%
\pgfpathlineto{\pgfqpoint{2.047686in}{0.904649in}}%
\pgfpathlineto{\pgfqpoint{2.059227in}{1.020085in}}%
\pgfpathlineto{\pgfqpoint{2.060902in}{1.034968in}}%
\pgfpathlineto{\pgfqpoint{2.063508in}{1.043722in}}%
\pgfpathlineto{\pgfqpoint{2.069092in}{1.108791in}}%
\pgfpathlineto{\pgfqpoint{2.070209in}{1.110000in}}%
\pgfpathlineto{\pgfqpoint{2.071698in}{1.131910in}}%
\pgfpathlineto{\pgfqpoint{2.073560in}{1.140861in}}%
\pgfpathlineto{\pgfqpoint{2.074863in}{1.143122in}}%
\pgfpathlineto{\pgfqpoint{2.075980in}{1.149835in}}%
\pgfpathlineto{\pgfqpoint{2.080075in}{1.178767in}}%
\pgfpathlineto{\pgfqpoint{2.084170in}{1.190159in}}%
\pgfpathlineto{\pgfqpoint{2.085287in}{1.191469in}}%
\pgfpathlineto{\pgfqpoint{2.088265in}{1.201551in}}%
\pgfpathlineto{\pgfqpoint{2.095711in}{1.203279in}}%
\pgfpathlineto{\pgfqpoint{2.098689in}{1.202232in}}%
\pgfpathlineto{\pgfqpoint{2.102226in}{1.198700in}}%
\pgfpathlineto{\pgfqpoint{2.106135in}{1.183233in}}%
\pgfpathlineto{\pgfqpoint{2.121026in}{1.076080in}}%
\pgfpathlineto{\pgfqpoint{2.147831in}{0.846592in}}%
\pgfpathlineto{\pgfqpoint{2.153974in}{0.816067in}}%
\pgfpathlineto{\pgfqpoint{2.158441in}{0.803391in}}%
\pgfpathlineto{\pgfqpoint{2.160489in}{0.801894in}}%
\pgfpathlineto{\pgfqpoint{2.169424in}{0.802902in}}%
\pgfpathlineto{\pgfqpoint{2.174636in}{0.811676in}}%
\pgfpathlineto{\pgfqpoint{2.182640in}{0.829963in}}%
\pgfpathlineto{\pgfqpoint{2.193436in}{0.860632in}}%
\pgfpathlineto{\pgfqpoint{2.207583in}{0.908264in}}%
\pgfpathlineto{\pgfqpoint{2.220241in}{0.957229in}}%
\pgfpathlineto{\pgfqpoint{2.222288in}{0.964850in}}%
\pgfpathlineto{\pgfqpoint{2.223964in}{0.969678in}}%
\pgfpathlineto{\pgfqpoint{2.227314in}{0.985872in}}%
\pgfpathlineto{\pgfqpoint{2.237180in}{1.015150in}}%
\pgfpathlineto{\pgfqpoint{2.242950in}{1.024601in}}%
\pgfpathlineto{\pgfqpoint{2.246487in}{1.031281in}}%
\pgfpathlineto{\pgfqpoint{2.248534in}{1.033751in}}%
\pgfpathlineto{\pgfqpoint{2.250396in}{1.037057in}}%
\pgfpathlineto{\pgfqpoint{2.255422in}{1.052893in}}%
\pgfpathlineto{\pgfqpoint{2.256911in}{1.056972in}}%
\pgfpathlineto{\pgfqpoint{2.258772in}{1.060618in}}%
\pgfpathlineto{\pgfqpoint{2.261564in}{1.073279in}}%
\pgfpathlineto{\pgfqpoint{2.267521in}{1.086828in}}%
\pgfpathlineto{\pgfqpoint{2.273105in}{1.095142in}}%
\pgfpathlineto{\pgfqpoint{2.276642in}{1.097087in}}%
\pgfpathlineto{\pgfqpoint{2.288555in}{1.106514in}}%
\pgfpathlineto{\pgfqpoint{2.291534in}{1.108331in}}%
\pgfpathlineto{\pgfqpoint{2.294326in}{1.114099in}}%
\pgfpathlineto{\pgfqpoint{2.306239in}{1.142819in}}%
\pgfpathlineto{\pgfqpoint{2.318710in}{1.149923in}}%
\pgfpathlineto{\pgfqpoint{2.331182in}{1.148830in}}%
\pgfpathlineto{\pgfqpoint{2.337139in}{1.148429in}}%
\pgfpathlineto{\pgfqpoint{2.344026in}{1.149493in}}%
\pgfpathlineto{\pgfqpoint{2.352775in}{1.153492in}}%
\pgfpathlineto{\pgfqpoint{2.355008in}{1.157301in}}%
\pgfpathlineto{\pgfqpoint{2.358731in}{1.167858in}}%
\pgfpathlineto{\pgfqpoint{2.360034in}{1.170242in}}%
\pgfpathlineto{\pgfqpoint{2.363012in}{1.172324in}}%
\pgfpathlineto{\pgfqpoint{2.366363in}{1.173844in}}%
\pgfpathlineto{\pgfqpoint{2.376229in}{1.175236in}}%
\pgfpathlineto{\pgfqpoint{2.384791in}{1.175725in}}%
\pgfpathlineto{\pgfqpoint{2.422020in}{1.176818in}}%
\pgfpathlineto{\pgfqpoint{2.442496in}{1.179677in}}%
\pgfpathlineto{\pgfqpoint{2.442496in}{1.179677in}}%
\pgfusepath{stroke}%
\end{pgfscope}%
\begin{pgfscope}%
\pgfsetrectcap%
\pgfsetmiterjoin%
\pgfsetlinewidth{0.803000pt}%
\definecolor{currentstroke}{rgb}{0.000000,0.000000,0.000000}%
\pgfsetstrokecolor{currentstroke}%
\pgfsetdash{}{0pt}%
\pgfpathmoveto{\pgfqpoint{0.464911in}{0.456793in}}%
\pgfpathlineto{\pgfqpoint{0.464911in}{1.910105in}}%
\pgfusepath{stroke}%
\end{pgfscope}%
\begin{pgfscope}%
\pgfsetrectcap%
\pgfsetmiterjoin%
\pgfsetlinewidth{0.803000pt}%
\definecolor{currentstroke}{rgb}{0.000000,0.000000,0.000000}%
\pgfsetstrokecolor{currentstroke}%
\pgfsetdash{}{0pt}%
\pgfpathmoveto{\pgfqpoint{0.464911in}{0.456793in}}%
\pgfpathlineto{\pgfqpoint{2.442496in}{0.456793in}}%
\pgfusepath{stroke}%
\end{pgfscope}%
\begin{pgfscope}%
\pgfsetbuttcap%
\pgfsetroundjoin%
\definecolor{currentfill}{rgb}{0.000000,0.000000,0.000000}%
\pgfsetfillcolor{currentfill}%
\pgfsetlinewidth{0.803000pt}%
\definecolor{currentstroke}{rgb}{0.000000,0.000000,0.000000}%
\pgfsetstrokecolor{currentstroke}%
\pgfsetdash{}{0pt}%
\pgfsys@defobject{currentmarker}{\pgfqpoint{0.000000in}{0.000000in}}{\pgfqpoint{0.048611in}{0.000000in}}{%
\pgfpathmoveto{\pgfqpoint{0.000000in}{0.000000in}}%
\pgfpathlineto{\pgfqpoint{0.048611in}{0.000000in}}%
\pgfusepath{stroke,fill}%
}%
\begin{pgfscope}%
\pgfsys@transformshift{2.442496in}{0.456793in}%
\pgfsys@useobject{currentmarker}{}%
\end{pgfscope}%
\end{pgfscope}%
\begin{pgfscope}%
\definecolor{textcolor}{rgb}{0.000000,0.000000,0.000000}%
\pgfsetstrokecolor{textcolor}%
\pgfsetfillcolor{textcolor}%
\pgftext[x=2.539718in, y=0.408568in, left, base]{\color{textcolor}{\rmfamily\fontsize{10.000000}{12.000000}\selectfont\catcode`\^=\active\def^{\ifmmode\sp\else\^{}\fi}\catcode`\%=\active\def%{\%}0}}%
\end{pgfscope}%
\begin{pgfscope}%
\pgfsetbuttcap%
\pgfsetroundjoin%
\definecolor{currentfill}{rgb}{0.000000,0.000000,0.000000}%
\pgfsetfillcolor{currentfill}%
\pgfsetlinewidth{0.803000pt}%
\definecolor{currentstroke}{rgb}{0.000000,0.000000,0.000000}%
\pgfsetstrokecolor{currentstroke}%
\pgfsetdash{}{0pt}%
\pgfsys@defobject{currentmarker}{\pgfqpoint{0.000000in}{0.000000in}}{\pgfqpoint{0.048611in}{0.000000in}}{%
\pgfpathmoveto{\pgfqpoint{0.000000in}{0.000000in}}%
\pgfpathlineto{\pgfqpoint{0.048611in}{0.000000in}}%
\pgfusepath{stroke,fill}%
}%
\begin{pgfscope}%
\pgfsys@transformshift{2.442496in}{0.699012in}%
\pgfsys@useobject{currentmarker}{}%
\end{pgfscope}%
\end{pgfscope}%
\begin{pgfscope}%
\definecolor{textcolor}{rgb}{0.000000,0.000000,0.000000}%
\pgfsetstrokecolor{textcolor}%
\pgfsetfillcolor{textcolor}%
\pgftext[x=2.539718in, y=0.650787in, left, base]{\color{textcolor}{\rmfamily\fontsize{10.000000}{12.000000}\selectfont\catcode`\^=\active\def^{\ifmmode\sp\else\^{}\fi}\catcode`\%=\active\def%{\%}20}}%
\end{pgfscope}%
\begin{pgfscope}%
\pgfsetbuttcap%
\pgfsetroundjoin%
\definecolor{currentfill}{rgb}{0.000000,0.000000,0.000000}%
\pgfsetfillcolor{currentfill}%
\pgfsetlinewidth{0.803000pt}%
\definecolor{currentstroke}{rgb}{0.000000,0.000000,0.000000}%
\pgfsetstrokecolor{currentstroke}%
\pgfsetdash{}{0pt}%
\pgfsys@defobject{currentmarker}{\pgfqpoint{0.000000in}{0.000000in}}{\pgfqpoint{0.048611in}{0.000000in}}{%
\pgfpathmoveto{\pgfqpoint{0.000000in}{0.000000in}}%
\pgfpathlineto{\pgfqpoint{0.048611in}{0.000000in}}%
\pgfusepath{stroke,fill}%
}%
\begin{pgfscope}%
\pgfsys@transformshift{2.442496in}{0.941230in}%
\pgfsys@useobject{currentmarker}{}%
\end{pgfscope}%
\end{pgfscope}%
\begin{pgfscope}%
\definecolor{textcolor}{rgb}{0.000000,0.000000,0.000000}%
\pgfsetstrokecolor{textcolor}%
\pgfsetfillcolor{textcolor}%
\pgftext[x=2.539718in, y=0.893005in, left, base]{\color{textcolor}{\rmfamily\fontsize{10.000000}{12.000000}\selectfont\catcode`\^=\active\def^{\ifmmode\sp\else\^{}\fi}\catcode`\%=\active\def%{\%}40}}%
\end{pgfscope}%
\begin{pgfscope}%
\pgfsetbuttcap%
\pgfsetroundjoin%
\definecolor{currentfill}{rgb}{0.000000,0.000000,0.000000}%
\pgfsetfillcolor{currentfill}%
\pgfsetlinewidth{0.803000pt}%
\definecolor{currentstroke}{rgb}{0.000000,0.000000,0.000000}%
\pgfsetstrokecolor{currentstroke}%
\pgfsetdash{}{0pt}%
\pgfsys@defobject{currentmarker}{\pgfqpoint{0.000000in}{0.000000in}}{\pgfqpoint{0.048611in}{0.000000in}}{%
\pgfpathmoveto{\pgfqpoint{0.000000in}{0.000000in}}%
\pgfpathlineto{\pgfqpoint{0.048611in}{0.000000in}}%
\pgfusepath{stroke,fill}%
}%
\begin{pgfscope}%
\pgfsys@transformshift{2.442496in}{1.183449in}%
\pgfsys@useobject{currentmarker}{}%
\end{pgfscope}%
\end{pgfscope}%
\begin{pgfscope}%
\definecolor{textcolor}{rgb}{0.000000,0.000000,0.000000}%
\pgfsetstrokecolor{textcolor}%
\pgfsetfillcolor{textcolor}%
\pgftext[x=2.539718in, y=1.135224in, left, base]{\color{textcolor}{\rmfamily\fontsize{10.000000}{12.000000}\selectfont\catcode`\^=\active\def^{\ifmmode\sp\else\^{}\fi}\catcode`\%=\active\def%{\%}60}}%
\end{pgfscope}%
\begin{pgfscope}%
\pgfsetbuttcap%
\pgfsetroundjoin%
\definecolor{currentfill}{rgb}{0.000000,0.000000,0.000000}%
\pgfsetfillcolor{currentfill}%
\pgfsetlinewidth{0.803000pt}%
\definecolor{currentstroke}{rgb}{0.000000,0.000000,0.000000}%
\pgfsetstrokecolor{currentstroke}%
\pgfsetdash{}{0pt}%
\pgfsys@defobject{currentmarker}{\pgfqpoint{0.000000in}{0.000000in}}{\pgfqpoint{0.048611in}{0.000000in}}{%
\pgfpathmoveto{\pgfqpoint{0.000000in}{0.000000in}}%
\pgfpathlineto{\pgfqpoint{0.048611in}{0.000000in}}%
\pgfusepath{stroke,fill}%
}%
\begin{pgfscope}%
\pgfsys@transformshift{2.442496in}{1.425668in}%
\pgfsys@useobject{currentmarker}{}%
\end{pgfscope}%
\end{pgfscope}%
\begin{pgfscope}%
\definecolor{textcolor}{rgb}{0.000000,0.000000,0.000000}%
\pgfsetstrokecolor{textcolor}%
\pgfsetfillcolor{textcolor}%
\pgftext[x=2.539718in, y=1.377442in, left, base]{\color{textcolor}{\rmfamily\fontsize{10.000000}{12.000000}\selectfont\catcode`\^=\active\def^{\ifmmode\sp\else\^{}\fi}\catcode`\%=\active\def%{\%}80}}%
\end{pgfscope}%
\begin{pgfscope}%
\pgfsetbuttcap%
\pgfsetroundjoin%
\definecolor{currentfill}{rgb}{0.000000,0.000000,0.000000}%
\pgfsetfillcolor{currentfill}%
\pgfsetlinewidth{0.803000pt}%
\definecolor{currentstroke}{rgb}{0.000000,0.000000,0.000000}%
\pgfsetstrokecolor{currentstroke}%
\pgfsetdash{}{0pt}%
\pgfsys@defobject{currentmarker}{\pgfqpoint{0.000000in}{0.000000in}}{\pgfqpoint{0.048611in}{0.000000in}}{%
\pgfpathmoveto{\pgfqpoint{0.000000in}{0.000000in}}%
\pgfpathlineto{\pgfqpoint{0.048611in}{0.000000in}}%
\pgfusepath{stroke,fill}%
}%
\begin{pgfscope}%
\pgfsys@transformshift{2.442496in}{1.667886in}%
\pgfsys@useobject{currentmarker}{}%
\end{pgfscope}%
\end{pgfscope}%
\begin{pgfscope}%
\definecolor{textcolor}{rgb}{0.000000,0.000000,0.000000}%
\pgfsetstrokecolor{textcolor}%
\pgfsetfillcolor{textcolor}%
\pgftext[x=2.539718in, y=1.619661in, left, base]{\color{textcolor}{\rmfamily\fontsize{10.000000}{12.000000}\selectfont\catcode`\^=\active\def^{\ifmmode\sp\else\^{}\fi}\catcode`\%=\active\def%{\%}100}}%
\end{pgfscope}%
\begin{pgfscope}%
\pgfsetbuttcap%
\pgfsetroundjoin%
\definecolor{currentfill}{rgb}{0.000000,0.000000,0.000000}%
\pgfsetfillcolor{currentfill}%
\pgfsetlinewidth{0.803000pt}%
\definecolor{currentstroke}{rgb}{0.000000,0.000000,0.000000}%
\pgfsetstrokecolor{currentstroke}%
\pgfsetdash{}{0pt}%
\pgfsys@defobject{currentmarker}{\pgfqpoint{0.000000in}{0.000000in}}{\pgfqpoint{0.048611in}{0.000000in}}{%
\pgfpathmoveto{\pgfqpoint{0.000000in}{0.000000in}}%
\pgfpathlineto{\pgfqpoint{0.048611in}{0.000000in}}%
\pgfusepath{stroke,fill}%
}%
\begin{pgfscope}%
\pgfsys@transformshift{2.442496in}{1.910105in}%
\pgfsys@useobject{currentmarker}{}%
\end{pgfscope}%
\end{pgfscope}%
\begin{pgfscope}%
\definecolor{textcolor}{rgb}{0.000000,0.000000,0.000000}%
\pgfsetstrokecolor{textcolor}%
\pgfsetfillcolor{textcolor}%
\pgftext[x=2.539718in, y=1.861879in, left, base]{\color{textcolor}{\rmfamily\fontsize{10.000000}{12.000000}\selectfont\catcode`\^=\active\def^{\ifmmode\sp\else\^{}\fi}\catcode`\%=\active\def%{\%}120}}%
\end{pgfscope}%
\begin{pgfscope}%
\definecolor{textcolor}{rgb}{0.000000,0.000000,0.000000}%
\pgfsetstrokecolor{textcolor}%
\pgfsetfillcolor{textcolor}%
\pgftext[x=2.803607in,y=1.183449in,,top,rotate=90.000000]{\color{textcolor}{\rmfamily\fontsize{10.000000}{12.000000}\selectfont\catcode`\^=\active\def^{\ifmmode\sp\else\^{}\fi}\catcode`\%=\active\def%{\%}$v$ / \unit{\kilo\metre\per\hour}}}%
\end{pgfscope}%
\begin{pgfscope}%
\pgfpathrectangle{\pgfqpoint{0.464911in}{0.456793in}}{\pgfqpoint{1.977584in}{1.453311in}}%
\pgfusepath{clip}%
\pgfsetrectcap%
\pgfsetroundjoin%
\pgfsetlinewidth{0.803000pt}%
\definecolor{currentstroke}{rgb}{0.172549,0.627451,0.172549}%
\pgfsetstrokecolor{currentstroke}%
\pgfsetdash{}{0pt}%
\pgfpathmoveto{\pgfqpoint{0.464911in}{0.456793in}}%
\pgfpathlineto{\pgfqpoint{0.502698in}{0.457864in}}%
\pgfpathlineto{\pgfqpoint{0.509027in}{0.467720in}}%
\pgfpathlineto{\pgfqpoint{0.515914in}{0.479876in}}%
\pgfpathlineto{\pgfqpoint{0.518148in}{0.491258in}}%
\pgfpathlineto{\pgfqpoint{0.520382in}{0.517062in}}%
\pgfpathlineto{\pgfqpoint{0.528572in}{0.643977in}}%
\pgfpathlineto{\pgfqpoint{0.529317in}{0.641699in}}%
\pgfpathlineto{\pgfqpoint{0.529875in}{0.642642in}}%
\pgfpathlineto{\pgfqpoint{0.532667in}{0.658165in}}%
\pgfpathlineto{\pgfqpoint{0.538624in}{0.706168in}}%
\pgfpathlineto{\pgfqpoint{0.560217in}{0.904647in}}%
\pgfpathlineto{\pgfqpoint{0.562450in}{0.933379in}}%
\pgfpathlineto{\pgfqpoint{0.563753in}{0.940518in}}%
\pgfpathlineto{\pgfqpoint{0.571571in}{1.005799in}}%
\pgfpathlineto{\pgfqpoint{0.597631in}{1.247558in}}%
\pgfpathlineto{\pgfqpoint{0.603774in}{1.298051in}}%
\pgfpathlineto{\pgfqpoint{0.606194in}{1.325320in}}%
\pgfpathlineto{\pgfqpoint{0.608428in}{1.346946in}}%
\pgfpathlineto{\pgfqpoint{0.613826in}{1.390216in}}%
\pgfpathlineto{\pgfqpoint{0.632068in}{1.531344in}}%
\pgfpathlineto{\pgfqpoint{0.640816in}{1.578138in}}%
\pgfpathlineto{\pgfqpoint{0.650310in}{1.613361in}}%
\pgfpathlineto{\pgfqpoint{0.654405in}{1.625774in}}%
\pgfpathlineto{\pgfqpoint{0.656266in}{1.630953in}}%
\pgfpathlineto{\pgfqpoint{0.658128in}{1.633649in}}%
\pgfpathlineto{\pgfqpoint{0.666690in}{1.647035in}}%
\pgfpathlineto{\pgfqpoint{0.675811in}{1.656829in}}%
\pgfpathlineto{\pgfqpoint{0.681023in}{1.661368in}}%
\pgfpathlineto{\pgfqpoint{0.682140in}{1.662972in}}%
\pgfpathlineto{\pgfqpoint{0.682513in}{1.662819in}}%
\pgfpathlineto{\pgfqpoint{0.684746in}{1.662962in}}%
\pgfpathlineto{\pgfqpoint{0.698707in}{1.666599in}}%
\pgfpathlineto{\pgfqpoint{0.715274in}{1.668363in}}%
\pgfpathlineto{\pgfqpoint{0.722906in}{1.667537in}}%
\pgfpathlineto{\pgfqpoint{0.727931in}{1.664926in}}%
\pgfpathlineto{\pgfqpoint{0.730537in}{1.664116in}}%
\pgfpathlineto{\pgfqpoint{0.734633in}{1.660210in}}%
\pgfpathlineto{\pgfqpoint{0.738914in}{1.653796in}}%
\pgfpathlineto{\pgfqpoint{0.741892in}{1.641013in}}%
\pgfpathlineto{\pgfqpoint{0.746546in}{1.601435in}}%
\pgfpathlineto{\pgfqpoint{0.765532in}{1.433343in}}%
\pgfpathlineto{\pgfqpoint{0.775026in}{1.372515in}}%
\pgfpathlineto{\pgfqpoint{0.781168in}{1.345396in}}%
\pgfpathlineto{\pgfqpoint{0.786939in}{1.329632in}}%
\pgfpathlineto{\pgfqpoint{0.792709in}{1.319944in}}%
\pgfpathlineto{\pgfqpoint{0.802016in}{1.309745in}}%
\pgfpathlineto{\pgfqpoint{0.810579in}{1.302083in}}%
\pgfpathlineto{\pgfqpoint{0.825657in}{1.293186in}}%
\pgfpathlineto{\pgfqpoint{0.829007in}{1.293453in}}%
\pgfpathlineto{\pgfqpoint{0.830496in}{1.296109in}}%
\pgfpathlineto{\pgfqpoint{0.831055in}{1.295641in}}%
\pgfpathlineto{\pgfqpoint{0.833102in}{1.295110in}}%
\pgfpathlineto{\pgfqpoint{0.847808in}{1.296891in}}%
\pgfpathlineto{\pgfqpoint{0.898625in}{1.304727in}}%
\pgfpathlineto{\pgfqpoint{0.947022in}{1.306885in}}%
\pgfpathlineto{\pgfqpoint{0.965078in}{1.305923in}}%
\pgfpathlineto{\pgfqpoint{0.999328in}{1.303340in}}%
\pgfpathlineto{\pgfqpoint{1.035440in}{1.303891in}}%
\pgfpathlineto{\pgfqpoint{1.091097in}{1.304327in}}%
\pgfpathlineto{\pgfqpoint{1.097426in}{1.302331in}}%
\pgfpathlineto{\pgfqpoint{1.100962in}{1.298806in}}%
\pgfpathlineto{\pgfqpoint{1.110269in}{1.286352in}}%
\pgfpathlineto{\pgfqpoint{1.112503in}{1.273384in}}%
\pgfpathlineto{\pgfqpoint{1.117715in}{1.224700in}}%
\pgfpathlineto{\pgfqpoint{1.152896in}{0.875894in}}%
\pgfpathlineto{\pgfqpoint{1.169463in}{0.719317in}}%
\pgfpathlineto{\pgfqpoint{1.171697in}{0.692956in}}%
\pgfpathlineto{\pgfqpoint{1.172813in}{0.689686in}}%
\pgfpathlineto{\pgfqpoint{1.178398in}{0.661209in}}%
\pgfpathlineto{\pgfqpoint{1.189008in}{0.611261in}}%
\pgfpathlineto{\pgfqpoint{1.195895in}{0.596552in}}%
\pgfpathlineto{\pgfqpoint{1.196267in}{0.596782in}}%
\pgfpathlineto{\pgfqpoint{1.198129in}{0.599659in}}%
\pgfpathlineto{\pgfqpoint{1.201107in}{0.609244in}}%
\pgfpathlineto{\pgfqpoint{1.205761in}{0.633196in}}%
\pgfpathlineto{\pgfqpoint{1.213020in}{0.684053in}}%
\pgfpathlineto{\pgfqpoint{1.233868in}{0.850759in}}%
\pgfpathlineto{\pgfqpoint{1.235358in}{0.858783in}}%
\pgfpathlineto{\pgfqpoint{1.236288in}{0.861618in}}%
\pgfpathlineto{\pgfqpoint{1.242059in}{0.895107in}}%
\pgfpathlineto{\pgfqpoint{1.263279in}{1.020010in}}%
\pgfpathlineto{\pgfqpoint{1.266257in}{1.038219in}}%
\pgfpathlineto{\pgfqpoint{1.268119in}{1.052012in}}%
\pgfpathlineto{\pgfqpoint{1.268863in}{1.053091in}}%
\pgfpathlineto{\pgfqpoint{1.272586in}{1.067542in}}%
\pgfpathlineto{\pgfqpoint{1.297157in}{1.164126in}}%
\pgfpathlineto{\pgfqpoint{1.308884in}{1.199672in}}%
\pgfpathlineto{\pgfqpoint{1.319122in}{1.224251in}}%
\pgfpathlineto{\pgfqpoint{1.320983in}{1.230341in}}%
\pgfpathlineto{\pgfqpoint{1.322845in}{1.236410in}}%
\pgfpathlineto{\pgfqpoint{1.324334in}{1.238085in}}%
\pgfpathlineto{\pgfqpoint{1.345927in}{1.272429in}}%
\pgfpathlineto{\pgfqpoint{1.358957in}{1.288145in}}%
\pgfpathlineto{\pgfqpoint{1.384086in}{1.316736in}}%
\pgfpathlineto{\pgfqpoint{1.390973in}{1.329829in}}%
\pgfpathlineto{\pgfqpoint{1.397674in}{1.347616in}}%
\pgfpathlineto{\pgfqpoint{1.404189in}{1.371491in}}%
\pgfpathlineto{\pgfqpoint{1.410704in}{1.403133in}}%
\pgfpathlineto{\pgfqpoint{1.411821in}{1.400248in}}%
\pgfpathlineto{\pgfqpoint{1.412193in}{1.400056in}}%
\pgfpathlineto{\pgfqpoint{1.412566in}{1.400803in}}%
\pgfpathlineto{\pgfqpoint{1.414427in}{1.411177in}}%
\pgfpathlineto{\pgfqpoint{1.428016in}{1.494498in}}%
\pgfpathlineto{\pgfqpoint{1.431925in}{1.502764in}}%
\pgfpathlineto{\pgfqpoint{1.435834in}{1.509605in}}%
\pgfpathlineto{\pgfqpoint{1.437881in}{1.508416in}}%
\pgfpathlineto{\pgfqpoint{1.441976in}{1.502439in}}%
\pgfpathlineto{\pgfqpoint{1.443093in}{1.501441in}}%
\pgfpathlineto{\pgfqpoint{1.443466in}{1.501607in}}%
\pgfpathlineto{\pgfqpoint{1.445513in}{1.501136in}}%
\pgfpathlineto{\pgfqpoint{1.448491in}{1.498017in}}%
\pgfpathlineto{\pgfqpoint{1.450353in}{1.491944in}}%
\pgfpathlineto{\pgfqpoint{1.453890in}{1.469965in}}%
\pgfpathlineto{\pgfqpoint{1.474179in}{1.352892in}}%
\pgfpathlineto{\pgfqpoint{1.482556in}{1.320184in}}%
\pgfpathlineto{\pgfqpoint{1.491118in}{1.295497in}}%
\pgfpathlineto{\pgfqpoint{1.503031in}{1.268559in}}%
\pgfpathlineto{\pgfqpoint{1.524438in}{1.220625in}}%
\pgfpathlineto{\pgfqpoint{1.536537in}{1.187254in}}%
\pgfpathlineto{\pgfqpoint{1.547892in}{1.149334in}}%
\pgfpathlineto{\pgfqpoint{1.565762in}{1.079272in}}%
\pgfpathlineto{\pgfqpoint{1.581956in}{1.001565in}}%
\pgfpathlineto{\pgfqpoint{1.588657in}{0.973772in}}%
\pgfpathlineto{\pgfqpoint{1.594800in}{0.958237in}}%
\pgfpathlineto{\pgfqpoint{1.600198in}{0.944293in}}%
\pgfpathlineto{\pgfqpoint{1.607644in}{0.912591in}}%
\pgfpathlineto{\pgfqpoint{1.622535in}{0.838631in}}%
\pgfpathlineto{\pgfqpoint{1.629795in}{0.823068in}}%
\pgfpathlineto{\pgfqpoint{1.633890in}{0.814821in}}%
\pgfpathlineto{\pgfqpoint{1.635565in}{0.812104in}}%
\pgfpathlineto{\pgfqpoint{1.640219in}{0.807918in}}%
\pgfpathlineto{\pgfqpoint{1.646920in}{0.799240in}}%
\pgfpathlineto{\pgfqpoint{1.653249in}{0.789309in}}%
\pgfpathlineto{\pgfqpoint{1.655669in}{0.776038in}}%
\pgfpathlineto{\pgfqpoint{1.666837in}{0.705857in}}%
\pgfpathlineto{\pgfqpoint{1.675772in}{0.686507in}}%
\pgfpathlineto{\pgfqpoint{1.678750in}{0.685026in}}%
\pgfpathlineto{\pgfqpoint{1.682846in}{0.685519in}}%
\pgfpathlineto{\pgfqpoint{1.688802in}{0.688861in}}%
\pgfpathlineto{\pgfqpoint{1.694573in}{0.691299in}}%
\pgfpathlineto{\pgfqpoint{1.696992in}{0.689994in}}%
\pgfpathlineto{\pgfqpoint{1.700901in}{0.683920in}}%
\pgfpathlineto{\pgfqpoint{1.709278in}{0.668787in}}%
\pgfpathlineto{\pgfqpoint{1.711884in}{0.653083in}}%
\pgfpathlineto{\pgfqpoint{1.721563in}{0.590118in}}%
\pgfpathlineto{\pgfqpoint{1.726217in}{0.581592in}}%
\pgfpathlineto{\pgfqpoint{1.738875in}{0.565907in}}%
\pgfpathlineto{\pgfqpoint{1.742411in}{0.565368in}}%
\pgfpathlineto{\pgfqpoint{1.746507in}{0.567191in}}%
\pgfpathlineto{\pgfqpoint{1.755814in}{0.574501in}}%
\pgfpathlineto{\pgfqpoint{1.761770in}{0.577694in}}%
\pgfpathlineto{\pgfqpoint{1.769030in}{0.578717in}}%
\pgfpathlineto{\pgfqpoint{1.789506in}{0.578384in}}%
\pgfpathlineto{\pgfqpoint{1.795090in}{0.576060in}}%
\pgfpathlineto{\pgfqpoint{1.799743in}{0.571524in}}%
\pgfpathlineto{\pgfqpoint{1.809609in}{0.557861in}}%
\pgfpathlineto{\pgfqpoint{1.814076in}{0.544174in}}%
\pgfpathlineto{\pgfqpoint{1.820591in}{0.526376in}}%
\pgfpathlineto{\pgfqpoint{1.824128in}{0.522652in}}%
\pgfpathlineto{\pgfqpoint{1.827851in}{0.522517in}}%
\pgfpathlineto{\pgfqpoint{1.838647in}{0.524131in}}%
\pgfpathlineto{\pgfqpoint{1.841067in}{0.527477in}}%
\pgfpathlineto{\pgfqpoint{1.843487in}{0.535043in}}%
\pgfpathlineto{\pgfqpoint{1.846279in}{0.552338in}}%
\pgfpathlineto{\pgfqpoint{1.852980in}{0.612969in}}%
\pgfpathlineto{\pgfqpoint{1.854656in}{0.642157in}}%
\pgfpathlineto{\pgfqpoint{1.855400in}{0.639668in}}%
\pgfpathlineto{\pgfqpoint{1.855773in}{0.638965in}}%
\pgfpathlineto{\pgfqpoint{1.856331in}{0.639931in}}%
\pgfpathlineto{\pgfqpoint{1.860426in}{0.654966in}}%
\pgfpathlineto{\pgfqpoint{1.865824in}{0.685737in}}%
\pgfpathlineto{\pgfqpoint{1.874945in}{0.753494in}}%
\pgfpathlineto{\pgfqpoint{1.887417in}{0.857968in}}%
\pgfpathlineto{\pgfqpoint{1.888906in}{0.866867in}}%
\pgfpathlineto{\pgfqpoint{1.889651in}{0.868756in}}%
\pgfpathlineto{\pgfqpoint{1.896724in}{0.904974in}}%
\pgfpathlineto{\pgfqpoint{1.918689in}{1.032438in}}%
\pgfpathlineto{\pgfqpoint{1.919992in}{1.043275in}}%
\pgfpathlineto{\pgfqpoint{1.920364in}{1.043120in}}%
\pgfpathlineto{\pgfqpoint{1.921109in}{1.044133in}}%
\pgfpathlineto{\pgfqpoint{1.927438in}{1.066461in}}%
\pgfpathlineto{\pgfqpoint{1.957779in}{1.177127in}}%
\pgfpathlineto{\pgfqpoint{1.969878in}{1.210256in}}%
\pgfpathlineto{\pgfqpoint{1.976766in}{1.227773in}}%
\pgfpathlineto{\pgfqpoint{1.978441in}{1.231597in}}%
\pgfpathlineto{\pgfqpoint{1.979930in}{1.233345in}}%
\pgfpathlineto{\pgfqpoint{1.999103in}{1.262617in}}%
\pgfpathlineto{\pgfqpoint{2.011388in}{1.276668in}}%
\pgfpathlineto{\pgfqpoint{2.031492in}{1.298632in}}%
\pgfpathlineto{\pgfqpoint{2.038565in}{1.310975in}}%
\pgfpathlineto{\pgfqpoint{2.045266in}{1.327462in}}%
\pgfpathlineto{\pgfqpoint{2.051781in}{1.349665in}}%
\pgfpathlineto{\pgfqpoint{2.058668in}{1.381547in}}%
\pgfpathlineto{\pgfqpoint{2.059227in}{1.384115in}}%
\pgfpathlineto{\pgfqpoint{2.059785in}{1.382958in}}%
\pgfpathlineto{\pgfqpoint{2.060530in}{1.381050in}}%
\pgfpathlineto{\pgfqpoint{2.061088in}{1.382064in}}%
\pgfpathlineto{\pgfqpoint{2.063136in}{1.393794in}}%
\pgfpathlineto{\pgfqpoint{2.070395in}{1.449759in}}%
\pgfpathlineto{\pgfqpoint{2.073001in}{1.464485in}}%
\pgfpathlineto{\pgfqpoint{2.079330in}{1.493156in}}%
\pgfpathlineto{\pgfqpoint{2.082681in}{1.498197in}}%
\pgfpathlineto{\pgfqpoint{2.084356in}{1.500711in}}%
\pgfpathlineto{\pgfqpoint{2.086590in}{1.504040in}}%
\pgfpathlineto{\pgfqpoint{2.088451in}{1.502869in}}%
\pgfpathlineto{\pgfqpoint{2.093291in}{1.495688in}}%
\pgfpathlineto{\pgfqpoint{2.095153in}{1.494481in}}%
\pgfpathlineto{\pgfqpoint{2.096828in}{1.493793in}}%
\pgfpathlineto{\pgfqpoint{2.098317in}{1.489876in}}%
\pgfpathlineto{\pgfqpoint{2.101481in}{1.472639in}}%
\pgfpathlineto{\pgfqpoint{2.122702in}{1.347425in}}%
\pgfpathlineto{\pgfqpoint{2.131264in}{1.312868in}}%
\pgfpathlineto{\pgfqpoint{2.139455in}{1.288684in}}%
\pgfpathlineto{\pgfqpoint{2.149692in}{1.265676in}}%
\pgfpathlineto{\pgfqpoint{2.173705in}{1.213991in}}%
\pgfpathlineto{\pgfqpoint{2.185618in}{1.181363in}}%
\pgfpathlineto{\pgfqpoint{2.197531in}{1.142075in}}%
\pgfpathlineto{\pgfqpoint{2.210747in}{1.090284in}}%
\pgfpathlineto{\pgfqpoint{2.228803in}{1.001931in}}%
\pgfpathlineto{\pgfqpoint{2.235504in}{0.974480in}}%
\pgfpathlineto{\pgfqpoint{2.241833in}{0.958432in}}%
\pgfpathlineto{\pgfqpoint{2.247231in}{0.944183in}}%
\pgfpathlineto{\pgfqpoint{2.255049in}{0.910825in}}%
\pgfpathlineto{\pgfqpoint{2.267893in}{0.844572in}}%
\pgfpathlineto{\pgfqpoint{2.274036in}{0.828373in}}%
\pgfpathlineto{\pgfqpoint{2.285763in}{0.809001in}}%
\pgfpathlineto{\pgfqpoint{2.290975in}{0.800644in}}%
\pgfpathlineto{\pgfqpoint{2.293581in}{0.791139in}}%
\pgfpathlineto{\pgfqpoint{2.297862in}{0.764423in}}%
\pgfpathlineto{\pgfqpoint{2.307170in}{0.707279in}}%
\pgfpathlineto{\pgfqpoint{2.311637in}{0.697659in}}%
\pgfpathlineto{\pgfqpoint{2.313498in}{0.692314in}}%
\pgfpathlineto{\pgfqpoint{2.315360in}{0.687869in}}%
\pgfpathlineto{\pgfqpoint{2.321689in}{0.684325in}}%
\pgfpathlineto{\pgfqpoint{2.326715in}{0.685471in}}%
\pgfpathlineto{\pgfqpoint{2.340861in}{0.690161in}}%
\pgfpathlineto{\pgfqpoint{2.343095in}{0.687212in}}%
\pgfpathlineto{\pgfqpoint{2.353891in}{0.667433in}}%
\pgfpathlineto{\pgfqpoint{2.357242in}{0.647957in}}%
\pgfpathlineto{\pgfqpoint{2.360965in}{0.615583in}}%
\pgfpathlineto{\pgfqpoint{2.364315in}{0.598364in}}%
\pgfpathlineto{\pgfqpoint{2.369155in}{0.586281in}}%
\pgfpathlineto{\pgfqpoint{2.375298in}{0.577152in}}%
\pgfpathlineto{\pgfqpoint{2.384419in}{0.566153in}}%
\pgfpathlineto{\pgfqpoint{2.387956in}{0.565457in}}%
\pgfpathlineto{\pgfqpoint{2.391865in}{0.567047in}}%
\pgfpathlineto{\pgfqpoint{2.400055in}{0.573396in}}%
\pgfpathlineto{\pgfqpoint{2.406942in}{0.577635in}}%
\pgfpathlineto{\pgfqpoint{2.413830in}{0.578749in}}%
\pgfpathlineto{\pgfqpoint{2.430582in}{0.578260in}}%
\pgfpathlineto{\pgfqpoint{2.442496in}{0.578037in}}%
\pgfpathlineto{\pgfqpoint{2.442496in}{0.578037in}}%
\pgfusepath{stroke}%
\end{pgfscope}%
\begin{pgfscope}%
\pgfsetrectcap%
\pgfsetmiterjoin%
\pgfsetlinewidth{0.803000pt}%
\definecolor{currentstroke}{rgb}{0.000000,0.000000,0.000000}%
\pgfsetstrokecolor{currentstroke}%
\pgfsetdash{}{0pt}%
\pgfpathmoveto{\pgfqpoint{0.464911in}{0.456793in}}%
\pgfpathlineto{\pgfqpoint{0.464911in}{1.910105in}}%
\pgfusepath{stroke}%
\end{pgfscope}%
\begin{pgfscope}%
\pgfsetrectcap%
\pgfsetmiterjoin%
\pgfsetlinewidth{0.803000pt}%
\definecolor{currentstroke}{rgb}{0.000000,0.000000,0.000000}%
\pgfsetstrokecolor{currentstroke}%
\pgfsetdash{}{0pt}%
\pgfpathmoveto{\pgfqpoint{2.442496in}{0.456793in}}%
\pgfpathlineto{\pgfqpoint{2.442496in}{1.910105in}}%
\pgfusepath{stroke}%
\end{pgfscope}%
\begin{pgfscope}%
\pgfsetrectcap%
\pgfsetmiterjoin%
\pgfsetlinewidth{0.803000pt}%
\definecolor{currentstroke}{rgb}{0.000000,0.000000,0.000000}%
\pgfsetstrokecolor{currentstroke}%
\pgfsetdash{}{0pt}%
\pgfpathmoveto{\pgfqpoint{0.464911in}{0.456793in}}%
\pgfpathlineto{\pgfqpoint{2.442496in}{0.456793in}}%
\pgfusepath{stroke}%
\end{pgfscope}%
\end{pgfpicture}%
\makeatother%
\endgroup%

    \end{minipage}

    %% ── Zeile 2: Captions (beide bündig auf gleicher Höhe) ───────────
    \begin{minipage}[t]{0.49\textwidth}
      \subcaption{Räumliche Trajektorie mit farbkodierten Manöverabschnitten.
        \label{fig:sim_strecke_tikz}}
    \end{minipage}
    \hfill
    \begin{minipage}[t]{0.49\textwidth}
      \subcaption{Signalverläufe mit klassifizierten Manöverabschnitten.
        \label{fig:sim_strecke_carmaker}}
    \end{minipage}

    \smallskip

    %% ── Gemeinsame Legende ───────────────────────────────────────────
    \begin{tikzpicture}[baseline=(current bounding box.center)]

  %% ── Spalte 1: ax / ay ────────────────────────────────────────────
  \draw[color=clrAx, line width=0.8pt]
    (0cm, 0mm) -- ++(14pt, 0);
  \node[font=\tiny, anchor=west]
    at ({16pt}, 0mm) {$a_\mathrm{x}$};

  \draw[color=clrAy, line width=0.8pt]
    (0cm, -4mm) -- ++(14pt, 0);
  \node[font=\tiny, anchor=west]
    at ({16pt}, -4mm) {$a_\mathrm{y}$};

  %% ── Spalte 2: v / Gerade ─────────────────────────────────────────
  \draw[color=clrV, line width=0.8pt]
    (1.5cm, 0mm) -- ++(14pt, 0);
  \node[font=\tiny, anchor=west]
    at ({1.5cm + 16pt}, 0mm) {$v$};

  \fill[clrGerade, fill opacity=0.50]
    (1.5cm, -5.5mm) rectangle ++(14pt, 3mm);
  \node[font=\tiny, anchor=west]
    at ({1.5cm + 16pt}, -4mm) {konst. Gerade};

  %% ── Spalte 3: Beschl. / Bremsung ─────────────────────────────
  \fill[clrAccel, fill opacity=0.50]
    (4.25cm, -1.5mm) rectangle ++(14pt, 3mm);
  \node[font=\tiny, anchor=west]
    at ({4.25cm + 16pt}, 0mm) {Beschl.};

  \fill[clrBrake, fill opacity=0.50]
    (4.25cm, -5.5mm) rectangle ++(14pt, 3mm);
  \node[font=\tiny, anchor=west]
    at ({4.25cm + 16pt}, -4mm) {Bremsung};

  %% ── Spalte 3: Kurve li. / Kurve re. ─────────────────────────────
  \fill[clrCurveL, fill opacity=0.50]
    (6.5cm, -1.5mm) rectangle ++(14pt, 3mm);
  \node[font=\tiny, anchor=west]
    at ({6.5cm + 16pt}, 0mm) {Kurvenfahrt li.};

  \fill[clrCurveR, fill opacity=0.50]
    (6.5cm, -5.5mm) rectangle ++(14pt, 3mm);
  \node[font=\tiny, anchor=west]
    at ({6.5cm + 16pt}, -4mm) {Kurvenfahrt re.};

  %% ── Spalte 4: Gebr. Kurve li. / Gebr. Kurve re. ─────────────────
  \fill[clrCurveBrakeL, fill opacity=0.50]
    (9.5cm, -1.5mm) rectangle ++(14pt, 3mm);
  \node[font=\tiny, anchor=west]
    at ({9.5cm + 16pt}, 0mm) {gebr. Kurve li.};

  \fill[clrCurveBrakeR, fill opacity=0.50]
    (9.5cm, -5.5mm) rectangle ++(14pt, 3mm);
  \node[font=\tiny, anchor=west]
    at ({9.5cm + 16pt}, -4mm) {gebr. Kurve re.};

  %% ── Spalte 5: Beschl. Kurve li. / Beschl. Kurve re. ─────────────
  \fill[clrCurveAccelL, fill opacity=0.50]
    (12.5cm, -1.5mm) rectangle ++(14pt, 3mm);
  \node[font=\tiny, anchor=west]
    at ({12.5cm + 16pt}, 0mm) {beschl. Kurve li.};

  \fill[clrCurveAccelR, fill opacity=0.50]
    (12.5cm, -5.5mm) rectangle ++(14pt, 3mm);
  \node[font=\tiny, anchor=west]
    at ({12.5cm + 16pt}, -4mm) {beschl. Kurve re.};

\end{tikzpicture}

    \caption{Synthetische Simulationsstrecke als \textit{Ground Truth} zur Kalibrierung der Beschleunigungsschwellwerte.}
    \label{fig:sim_strecke}
  \end{figure}


% End of File
% !TeX root = ../main.tex
% 04-validation.tex
\chapter{Validierung der Reifenkräfte}
\label{chap:validation}

\lipsum[1]~\cite{hka_profil}.



% End of File
% !TeX root = ../main.tex
% 05-preparation.tex
\chapter{Entwicklung von \dots}
\label{chap:preparation}

\lipsum[1]



% End of File
% !TeX root = ../main.tex
% 06-evaluation.tex
\chapter{Prüfstandsversuche und Auswertung}
\label{chap:evaluation}

\lipsum[1]

\begin{figure}[htbp]
    \centering
    %% Creator: Matplotlib, PGF backend
%%
%% To include the figure in your LaTeX document, write
%%   \input{<filename>.pgf}
%%
%% Make sure the required packages are loaded in your preamble
%%   \usepackage{pgf}
%%
%% Also ensure that all the required font packages are loaded; for instance,
%% the lmodern package is sometimes necessary when using math font.
%%   \usepackage{lmodern}
%%
%% Figures using additional raster images can only be included by \input if
%% they are in the same directory as the main LaTeX file. For loading figures
%% from other directories you can use the `import` package
%%   \usepackage{import}
%%
%% and then include the figures with
%%   \import{<path to file>}{<filename>.pgf}
%%
%% Matplotlib used the following preamble
%%   \def\mathdefault#1{#1}
%%   \everymath=\expandafter{\the\everymath\displaystyle}
%%   \IfFileExists{scrextend.sty}{
%%     \usepackage[fontsize=10.000000pt]{scrextend}
%%   }{
%%     \renewcommand{\normalsize}{\fontsize{10.000000}{12.000000}\selectfont}
%%     \normalsize
%%   }
%%   \usepackage{amsmath}
%%   \usepackage{amssymb}
%%   \usepackage{newunicodechar}
%%   \newunicodechar{✓}{\checkmark}
%%   \usepackage{siunitx}
%%   \makeatletter\@ifpackageloaded{underscore}{}{\usepackage[strings]{underscore}}\makeatother
%%
\begingroup%
\makeatletter%
\begin{pgfpicture}%
\pgfpathrectangle{\pgfpointorigin}{\pgfqpoint{4.700000in}{2.227519in}}%
\pgfusepath{use as bounding box, clip}%
\begin{pgfscope}%
\pgfsetbuttcap%
\pgfsetmiterjoin%
\definecolor{currentfill}{rgb}{1.000000,1.000000,1.000000}%
\pgfsetfillcolor{currentfill}%
\pgfsetlinewidth{0.000000pt}%
\definecolor{currentstroke}{rgb}{1.000000,1.000000,1.000000}%
\pgfsetstrokecolor{currentstroke}%
\pgfsetdash{}{0pt}%
\pgfpathmoveto{\pgfqpoint{0.000000in}{0.000000in}}%
\pgfpathlineto{\pgfqpoint{4.700000in}{0.000000in}}%
\pgfpathlineto{\pgfqpoint{4.700000in}{2.227519in}}%
\pgfpathlineto{\pgfqpoint{0.000000in}{2.227519in}}%
\pgfpathlineto{\pgfqpoint{0.000000in}{0.000000in}}%
\pgfpathclose%
\pgfusepath{fill}%
\end{pgfscope}%
\begin{pgfscope}%
\pgfsetbuttcap%
\pgfsetmiterjoin%
\definecolor{currentfill}{rgb}{1.000000,1.000000,1.000000}%
\pgfsetfillcolor{currentfill}%
\pgfsetlinewidth{0.000000pt}%
\definecolor{currentstroke}{rgb}{0.000000,0.000000,0.000000}%
\pgfsetstrokecolor{currentstroke}%
\pgfsetstrokeopacity{0.000000}%
\pgfsetdash{}{0pt}%
\pgfpathmoveto{\pgfqpoint{0.367483in}{0.513891in}}%
\pgfpathlineto{\pgfqpoint{4.595833in}{0.513891in}}%
\pgfpathlineto{\pgfqpoint{4.595833in}{2.147238in}}%
\pgfpathlineto{\pgfqpoint{0.367483in}{2.147238in}}%
\pgfpathlineto{\pgfqpoint{0.367483in}{0.513891in}}%
\pgfpathclose%
\pgfusepath{fill}%
\end{pgfscope}%
\begin{pgfscope}%
\pgfpathrectangle{\pgfqpoint{0.367483in}{0.513891in}}{\pgfqpoint{4.228350in}{1.633348in}}%
\pgfusepath{clip}%
\pgfsetbuttcap%
\pgfsetmiterjoin%
\pgfsetlinewidth{0.803000pt}%
\definecolor{currentstroke}{rgb}{0.400000,0.400000,0.400000}%
\pgfsetstrokecolor{currentstroke}%
\pgfsetdash{{0.800000pt}{1.320000pt}}{0.000000pt}%
\pgfpathmoveto{\pgfqpoint{1.847405in}{0.615975in}}%
\pgfpathlineto{\pgfqpoint{2.270240in}{0.615975in}}%
\pgfpathlineto{\pgfqpoint{2.270240in}{0.922228in}}%
\pgfpathlineto{\pgfqpoint{1.847405in}{0.922228in}}%
\pgfpathlineto{\pgfqpoint{1.847405in}{0.615975in}}%
\pgfpathclose%
\pgfusepath{stroke}%
\end{pgfscope}%
\begin{pgfscope}%
\pgfpathrectangle{\pgfqpoint{0.367483in}{0.513891in}}{\pgfqpoint{4.228350in}{1.633348in}}%
\pgfusepath{clip}%
\pgfsetrectcap%
\pgfsetroundjoin%
\pgfsetlinewidth{0.803000pt}%
\definecolor{currentstroke}{rgb}{0.690196,0.690196,0.690196}%
\pgfsetstrokecolor{currentstroke}%
\pgfsetstrokeopacity{0.300000}%
\pgfsetdash{}{0pt}%
\pgfpathmoveto{\pgfqpoint{0.367483in}{0.513891in}}%
\pgfpathlineto{\pgfqpoint{0.367483in}{2.147238in}}%
\pgfusepath{stroke}%
\end{pgfscope}%
\begin{pgfscope}%
\pgfsetbuttcap%
\pgfsetroundjoin%
\definecolor{currentfill}{rgb}{0.000000,0.000000,0.000000}%
\pgfsetfillcolor{currentfill}%
\pgfsetlinewidth{0.803000pt}%
\definecolor{currentstroke}{rgb}{0.000000,0.000000,0.000000}%
\pgfsetstrokecolor{currentstroke}%
\pgfsetdash{}{0pt}%
\pgfsys@defobject{currentmarker}{\pgfqpoint{0.000000in}{-0.048611in}}{\pgfqpoint{0.000000in}{0.000000in}}{%
\pgfpathmoveto{\pgfqpoint{0.000000in}{0.000000in}}%
\pgfpathlineto{\pgfqpoint{0.000000in}{-0.048611in}}%
\pgfusepath{stroke,fill}%
}%
\begin{pgfscope}%
\pgfsys@transformshift{0.367483in}{0.513891in}%
\pgfsys@useobject{currentmarker}{}%
\end{pgfscope}%
\end{pgfscope}%
\begin{pgfscope}%
\definecolor{textcolor}{rgb}{0.000000,0.000000,0.000000}%
\pgfsetstrokecolor{textcolor}%
\pgfsetfillcolor{textcolor}%
\pgftext[x=0.367483in,y=0.416668in,,top]{\color{textcolor}{\rmfamily\fontsize{10.000000}{12.000000}\selectfont\catcode`\^=\active\def^{\ifmmode\sp\else\^{}\fi}\catcode`\%=\active\def%{\%}0}}%
\end{pgfscope}%
\begin{pgfscope}%
\pgfpathrectangle{\pgfqpoint{0.367483in}{0.513891in}}{\pgfqpoint{4.228350in}{1.633348in}}%
\pgfusepath{clip}%
\pgfsetrectcap%
\pgfsetroundjoin%
\pgfsetlinewidth{0.803000pt}%
\definecolor{currentstroke}{rgb}{0.690196,0.690196,0.690196}%
\pgfsetstrokecolor{currentstroke}%
\pgfsetstrokeopacity{0.300000}%
\pgfsetdash{}{0pt}%
\pgfpathmoveto{\pgfqpoint{0.896027in}{0.513891in}}%
\pgfpathlineto{\pgfqpoint{0.896027in}{2.147238in}}%
\pgfusepath{stroke}%
\end{pgfscope}%
\begin{pgfscope}%
\pgfsetbuttcap%
\pgfsetroundjoin%
\definecolor{currentfill}{rgb}{0.000000,0.000000,0.000000}%
\pgfsetfillcolor{currentfill}%
\pgfsetlinewidth{0.803000pt}%
\definecolor{currentstroke}{rgb}{0.000000,0.000000,0.000000}%
\pgfsetstrokecolor{currentstroke}%
\pgfsetdash{}{0pt}%
\pgfsys@defobject{currentmarker}{\pgfqpoint{0.000000in}{-0.048611in}}{\pgfqpoint{0.000000in}{0.000000in}}{%
\pgfpathmoveto{\pgfqpoint{0.000000in}{0.000000in}}%
\pgfpathlineto{\pgfqpoint{0.000000in}{-0.048611in}}%
\pgfusepath{stroke,fill}%
}%
\begin{pgfscope}%
\pgfsys@transformshift{0.896027in}{0.513891in}%
\pgfsys@useobject{currentmarker}{}%
\end{pgfscope}%
\end{pgfscope}%
\begin{pgfscope}%
\definecolor{textcolor}{rgb}{0.000000,0.000000,0.000000}%
\pgfsetstrokecolor{textcolor}%
\pgfsetfillcolor{textcolor}%
\pgftext[x=0.896027in,y=0.416668in,,top]{\color{textcolor}{\rmfamily\fontsize{10.000000}{12.000000}\selectfont\catcode`\^=\active\def^{\ifmmode\sp\else\^{}\fi}\catcode`\%=\active\def%{\%}50}}%
\end{pgfscope}%
\begin{pgfscope}%
\pgfpathrectangle{\pgfqpoint{0.367483in}{0.513891in}}{\pgfqpoint{4.228350in}{1.633348in}}%
\pgfusepath{clip}%
\pgfsetrectcap%
\pgfsetroundjoin%
\pgfsetlinewidth{0.803000pt}%
\definecolor{currentstroke}{rgb}{0.690196,0.690196,0.690196}%
\pgfsetstrokecolor{currentstroke}%
\pgfsetstrokeopacity{0.300000}%
\pgfsetdash{}{0pt}%
\pgfpathmoveto{\pgfqpoint{1.424570in}{0.513891in}}%
\pgfpathlineto{\pgfqpoint{1.424570in}{2.147238in}}%
\pgfusepath{stroke}%
\end{pgfscope}%
\begin{pgfscope}%
\pgfsetbuttcap%
\pgfsetroundjoin%
\definecolor{currentfill}{rgb}{0.000000,0.000000,0.000000}%
\pgfsetfillcolor{currentfill}%
\pgfsetlinewidth{0.803000pt}%
\definecolor{currentstroke}{rgb}{0.000000,0.000000,0.000000}%
\pgfsetstrokecolor{currentstroke}%
\pgfsetdash{}{0pt}%
\pgfsys@defobject{currentmarker}{\pgfqpoint{0.000000in}{-0.048611in}}{\pgfqpoint{0.000000in}{0.000000in}}{%
\pgfpathmoveto{\pgfqpoint{0.000000in}{0.000000in}}%
\pgfpathlineto{\pgfqpoint{0.000000in}{-0.048611in}}%
\pgfusepath{stroke,fill}%
}%
\begin{pgfscope}%
\pgfsys@transformshift{1.424570in}{0.513891in}%
\pgfsys@useobject{currentmarker}{}%
\end{pgfscope}%
\end{pgfscope}%
\begin{pgfscope}%
\definecolor{textcolor}{rgb}{0.000000,0.000000,0.000000}%
\pgfsetstrokecolor{textcolor}%
\pgfsetfillcolor{textcolor}%
\pgftext[x=1.424570in,y=0.416668in,,top]{\color{textcolor}{\rmfamily\fontsize{10.000000}{12.000000}\selectfont\catcode`\^=\active\def^{\ifmmode\sp\else\^{}\fi}\catcode`\%=\active\def%{\%}100}}%
\end{pgfscope}%
\begin{pgfscope}%
\pgfpathrectangle{\pgfqpoint{0.367483in}{0.513891in}}{\pgfqpoint{4.228350in}{1.633348in}}%
\pgfusepath{clip}%
\pgfsetrectcap%
\pgfsetroundjoin%
\pgfsetlinewidth{0.803000pt}%
\definecolor{currentstroke}{rgb}{0.690196,0.690196,0.690196}%
\pgfsetstrokecolor{currentstroke}%
\pgfsetstrokeopacity{0.300000}%
\pgfsetdash{}{0pt}%
\pgfpathmoveto{\pgfqpoint{1.953114in}{0.513891in}}%
\pgfpathlineto{\pgfqpoint{1.953114in}{2.147238in}}%
\pgfusepath{stroke}%
\end{pgfscope}%
\begin{pgfscope}%
\pgfsetbuttcap%
\pgfsetroundjoin%
\definecolor{currentfill}{rgb}{0.000000,0.000000,0.000000}%
\pgfsetfillcolor{currentfill}%
\pgfsetlinewidth{0.803000pt}%
\definecolor{currentstroke}{rgb}{0.000000,0.000000,0.000000}%
\pgfsetstrokecolor{currentstroke}%
\pgfsetdash{}{0pt}%
\pgfsys@defobject{currentmarker}{\pgfqpoint{0.000000in}{-0.048611in}}{\pgfqpoint{0.000000in}{0.000000in}}{%
\pgfpathmoveto{\pgfqpoint{0.000000in}{0.000000in}}%
\pgfpathlineto{\pgfqpoint{0.000000in}{-0.048611in}}%
\pgfusepath{stroke,fill}%
}%
\begin{pgfscope}%
\pgfsys@transformshift{1.953114in}{0.513891in}%
\pgfsys@useobject{currentmarker}{}%
\end{pgfscope}%
\end{pgfscope}%
\begin{pgfscope}%
\definecolor{textcolor}{rgb}{0.000000,0.000000,0.000000}%
\pgfsetstrokecolor{textcolor}%
\pgfsetfillcolor{textcolor}%
\pgftext[x=1.953114in,y=0.416668in,,top]{\color{textcolor}{\rmfamily\fontsize{10.000000}{12.000000}\selectfont\catcode`\^=\active\def^{\ifmmode\sp\else\^{}\fi}\catcode`\%=\active\def%{\%}150}}%
\end{pgfscope}%
\begin{pgfscope}%
\pgfpathrectangle{\pgfqpoint{0.367483in}{0.513891in}}{\pgfqpoint{4.228350in}{1.633348in}}%
\pgfusepath{clip}%
\pgfsetrectcap%
\pgfsetroundjoin%
\pgfsetlinewidth{0.803000pt}%
\definecolor{currentstroke}{rgb}{0.690196,0.690196,0.690196}%
\pgfsetstrokecolor{currentstroke}%
\pgfsetstrokeopacity{0.300000}%
\pgfsetdash{}{0pt}%
\pgfpathmoveto{\pgfqpoint{2.481658in}{0.513891in}}%
\pgfpathlineto{\pgfqpoint{2.481658in}{2.147238in}}%
\pgfusepath{stroke}%
\end{pgfscope}%
\begin{pgfscope}%
\pgfsetbuttcap%
\pgfsetroundjoin%
\definecolor{currentfill}{rgb}{0.000000,0.000000,0.000000}%
\pgfsetfillcolor{currentfill}%
\pgfsetlinewidth{0.803000pt}%
\definecolor{currentstroke}{rgb}{0.000000,0.000000,0.000000}%
\pgfsetstrokecolor{currentstroke}%
\pgfsetdash{}{0pt}%
\pgfsys@defobject{currentmarker}{\pgfqpoint{0.000000in}{-0.048611in}}{\pgfqpoint{0.000000in}{0.000000in}}{%
\pgfpathmoveto{\pgfqpoint{0.000000in}{0.000000in}}%
\pgfpathlineto{\pgfqpoint{0.000000in}{-0.048611in}}%
\pgfusepath{stroke,fill}%
}%
\begin{pgfscope}%
\pgfsys@transformshift{2.481658in}{0.513891in}%
\pgfsys@useobject{currentmarker}{}%
\end{pgfscope}%
\end{pgfscope}%
\begin{pgfscope}%
\definecolor{textcolor}{rgb}{0.000000,0.000000,0.000000}%
\pgfsetstrokecolor{textcolor}%
\pgfsetfillcolor{textcolor}%
\pgftext[x=2.481658in,y=0.416668in,,top]{\color{textcolor}{\rmfamily\fontsize{10.000000}{12.000000}\selectfont\catcode`\^=\active\def^{\ifmmode\sp\else\^{}\fi}\catcode`\%=\active\def%{\%}200}}%
\end{pgfscope}%
\begin{pgfscope}%
\pgfpathrectangle{\pgfqpoint{0.367483in}{0.513891in}}{\pgfqpoint{4.228350in}{1.633348in}}%
\pgfusepath{clip}%
\pgfsetrectcap%
\pgfsetroundjoin%
\pgfsetlinewidth{0.803000pt}%
\definecolor{currentstroke}{rgb}{0.690196,0.690196,0.690196}%
\pgfsetstrokecolor{currentstroke}%
\pgfsetstrokeopacity{0.300000}%
\pgfsetdash{}{0pt}%
\pgfpathmoveto{\pgfqpoint{3.010202in}{0.513891in}}%
\pgfpathlineto{\pgfqpoint{3.010202in}{2.147238in}}%
\pgfusepath{stroke}%
\end{pgfscope}%
\begin{pgfscope}%
\pgfsetbuttcap%
\pgfsetroundjoin%
\definecolor{currentfill}{rgb}{0.000000,0.000000,0.000000}%
\pgfsetfillcolor{currentfill}%
\pgfsetlinewidth{0.803000pt}%
\definecolor{currentstroke}{rgb}{0.000000,0.000000,0.000000}%
\pgfsetstrokecolor{currentstroke}%
\pgfsetdash{}{0pt}%
\pgfsys@defobject{currentmarker}{\pgfqpoint{0.000000in}{-0.048611in}}{\pgfqpoint{0.000000in}{0.000000in}}{%
\pgfpathmoveto{\pgfqpoint{0.000000in}{0.000000in}}%
\pgfpathlineto{\pgfqpoint{0.000000in}{-0.048611in}}%
\pgfusepath{stroke,fill}%
}%
\begin{pgfscope}%
\pgfsys@transformshift{3.010202in}{0.513891in}%
\pgfsys@useobject{currentmarker}{}%
\end{pgfscope}%
\end{pgfscope}%
\begin{pgfscope}%
\definecolor{textcolor}{rgb}{0.000000,0.000000,0.000000}%
\pgfsetstrokecolor{textcolor}%
\pgfsetfillcolor{textcolor}%
\pgftext[x=3.010202in,y=0.416668in,,top]{\color{textcolor}{\rmfamily\fontsize{10.000000}{12.000000}\selectfont\catcode`\^=\active\def^{\ifmmode\sp\else\^{}\fi}\catcode`\%=\active\def%{\%}250}}%
\end{pgfscope}%
\begin{pgfscope}%
\pgfpathrectangle{\pgfqpoint{0.367483in}{0.513891in}}{\pgfqpoint{4.228350in}{1.633348in}}%
\pgfusepath{clip}%
\pgfsetrectcap%
\pgfsetroundjoin%
\pgfsetlinewidth{0.803000pt}%
\definecolor{currentstroke}{rgb}{0.690196,0.690196,0.690196}%
\pgfsetstrokecolor{currentstroke}%
\pgfsetstrokeopacity{0.300000}%
\pgfsetdash{}{0pt}%
\pgfpathmoveto{\pgfqpoint{3.538745in}{0.513891in}}%
\pgfpathlineto{\pgfqpoint{3.538745in}{2.147238in}}%
\pgfusepath{stroke}%
\end{pgfscope}%
\begin{pgfscope}%
\pgfsetbuttcap%
\pgfsetroundjoin%
\definecolor{currentfill}{rgb}{0.000000,0.000000,0.000000}%
\pgfsetfillcolor{currentfill}%
\pgfsetlinewidth{0.803000pt}%
\definecolor{currentstroke}{rgb}{0.000000,0.000000,0.000000}%
\pgfsetstrokecolor{currentstroke}%
\pgfsetdash{}{0pt}%
\pgfsys@defobject{currentmarker}{\pgfqpoint{0.000000in}{-0.048611in}}{\pgfqpoint{0.000000in}{0.000000in}}{%
\pgfpathmoveto{\pgfqpoint{0.000000in}{0.000000in}}%
\pgfpathlineto{\pgfqpoint{0.000000in}{-0.048611in}}%
\pgfusepath{stroke,fill}%
}%
\begin{pgfscope}%
\pgfsys@transformshift{3.538745in}{0.513891in}%
\pgfsys@useobject{currentmarker}{}%
\end{pgfscope}%
\end{pgfscope}%
\begin{pgfscope}%
\definecolor{textcolor}{rgb}{0.000000,0.000000,0.000000}%
\pgfsetstrokecolor{textcolor}%
\pgfsetfillcolor{textcolor}%
\pgftext[x=3.538745in,y=0.416668in,,top]{\color{textcolor}{\rmfamily\fontsize{10.000000}{12.000000}\selectfont\catcode`\^=\active\def^{\ifmmode\sp\else\^{}\fi}\catcode`\%=\active\def%{\%}300}}%
\end{pgfscope}%
\begin{pgfscope}%
\pgfpathrectangle{\pgfqpoint{0.367483in}{0.513891in}}{\pgfqpoint{4.228350in}{1.633348in}}%
\pgfusepath{clip}%
\pgfsetrectcap%
\pgfsetroundjoin%
\pgfsetlinewidth{0.803000pt}%
\definecolor{currentstroke}{rgb}{0.690196,0.690196,0.690196}%
\pgfsetstrokecolor{currentstroke}%
\pgfsetstrokeopacity{0.300000}%
\pgfsetdash{}{0pt}%
\pgfpathmoveto{\pgfqpoint{4.067289in}{0.513891in}}%
\pgfpathlineto{\pgfqpoint{4.067289in}{2.147238in}}%
\pgfusepath{stroke}%
\end{pgfscope}%
\begin{pgfscope}%
\pgfsetbuttcap%
\pgfsetroundjoin%
\definecolor{currentfill}{rgb}{0.000000,0.000000,0.000000}%
\pgfsetfillcolor{currentfill}%
\pgfsetlinewidth{0.803000pt}%
\definecolor{currentstroke}{rgb}{0.000000,0.000000,0.000000}%
\pgfsetstrokecolor{currentstroke}%
\pgfsetdash{}{0pt}%
\pgfsys@defobject{currentmarker}{\pgfqpoint{0.000000in}{-0.048611in}}{\pgfqpoint{0.000000in}{0.000000in}}{%
\pgfpathmoveto{\pgfqpoint{0.000000in}{0.000000in}}%
\pgfpathlineto{\pgfqpoint{0.000000in}{-0.048611in}}%
\pgfusepath{stroke,fill}%
}%
\begin{pgfscope}%
\pgfsys@transformshift{4.067289in}{0.513891in}%
\pgfsys@useobject{currentmarker}{}%
\end{pgfscope}%
\end{pgfscope}%
\begin{pgfscope}%
\definecolor{textcolor}{rgb}{0.000000,0.000000,0.000000}%
\pgfsetstrokecolor{textcolor}%
\pgfsetfillcolor{textcolor}%
\pgftext[x=4.067289in,y=0.416668in,,top]{\color{textcolor}{\rmfamily\fontsize{10.000000}{12.000000}\selectfont\catcode`\^=\active\def^{\ifmmode\sp\else\^{}\fi}\catcode`\%=\active\def%{\%}350}}%
\end{pgfscope}%
\begin{pgfscope}%
\pgfpathrectangle{\pgfqpoint{0.367483in}{0.513891in}}{\pgfqpoint{4.228350in}{1.633348in}}%
\pgfusepath{clip}%
\pgfsetrectcap%
\pgfsetroundjoin%
\pgfsetlinewidth{0.803000pt}%
\definecolor{currentstroke}{rgb}{0.690196,0.690196,0.690196}%
\pgfsetstrokecolor{currentstroke}%
\pgfsetstrokeopacity{0.300000}%
\pgfsetdash{}{0pt}%
\pgfpathmoveto{\pgfqpoint{4.595833in}{0.513891in}}%
\pgfpathlineto{\pgfqpoint{4.595833in}{2.147238in}}%
\pgfusepath{stroke}%
\end{pgfscope}%
\begin{pgfscope}%
\pgfsetbuttcap%
\pgfsetroundjoin%
\definecolor{currentfill}{rgb}{0.000000,0.000000,0.000000}%
\pgfsetfillcolor{currentfill}%
\pgfsetlinewidth{0.803000pt}%
\definecolor{currentstroke}{rgb}{0.000000,0.000000,0.000000}%
\pgfsetstrokecolor{currentstroke}%
\pgfsetdash{}{0pt}%
\pgfsys@defobject{currentmarker}{\pgfqpoint{0.000000in}{-0.048611in}}{\pgfqpoint{0.000000in}{0.000000in}}{%
\pgfpathmoveto{\pgfqpoint{0.000000in}{0.000000in}}%
\pgfpathlineto{\pgfqpoint{0.000000in}{-0.048611in}}%
\pgfusepath{stroke,fill}%
}%
\begin{pgfscope}%
\pgfsys@transformshift{4.595833in}{0.513891in}%
\pgfsys@useobject{currentmarker}{}%
\end{pgfscope}%
\end{pgfscope}%
\begin{pgfscope}%
\definecolor{textcolor}{rgb}{0.000000,0.000000,0.000000}%
\pgfsetstrokecolor{textcolor}%
\pgfsetfillcolor{textcolor}%
\pgftext[x=4.595833in,y=0.416668in,,top]{\color{textcolor}{\rmfamily\fontsize{10.000000}{12.000000}\selectfont\catcode`\^=\active\def^{\ifmmode\sp\else\^{}\fi}\catcode`\%=\active\def%{\%}400}}%
\end{pgfscope}%
\begin{pgfscope}%
\definecolor{textcolor}{rgb}{0.000000,0.000000,0.000000}%
\pgfsetstrokecolor{textcolor}%
\pgfsetfillcolor{textcolor}%
\pgftext[x=2.481658in,y=0.222224in,,top]{\color{textcolor}{\rmfamily\fontsize{10.000000}{12.000000}\selectfont\catcode`\^=\active\def^{\ifmmode\sp\else\^{}\fi}\catcode`\%=\active\def%{\%}$\sum W_\mathrm{wear}$ / \unit{\mega\joule}}}%
\end{pgfscope}%
\begin{pgfscope}%
\pgfpathrectangle{\pgfqpoint{0.367483in}{0.513891in}}{\pgfqpoint{4.228350in}{1.633348in}}%
\pgfusepath{clip}%
\pgfsetrectcap%
\pgfsetroundjoin%
\pgfsetlinewidth{0.803000pt}%
\definecolor{currentstroke}{rgb}{0.690196,0.690196,0.690196}%
\pgfsetstrokecolor{currentstroke}%
\pgfsetstrokeopacity{0.300000}%
\pgfsetdash{}{0pt}%
\pgfpathmoveto{\pgfqpoint{0.367483in}{0.513891in}}%
\pgfpathlineto{\pgfqpoint{4.595833in}{0.513891in}}%
\pgfusepath{stroke}%
\end{pgfscope}%
\begin{pgfscope}%
\pgfsetbuttcap%
\pgfsetroundjoin%
\definecolor{currentfill}{rgb}{0.000000,0.000000,0.000000}%
\pgfsetfillcolor{currentfill}%
\pgfsetlinewidth{0.803000pt}%
\definecolor{currentstroke}{rgb}{0.000000,0.000000,0.000000}%
\pgfsetstrokecolor{currentstroke}%
\pgfsetdash{}{0pt}%
\pgfsys@defobject{currentmarker}{\pgfqpoint{-0.048611in}{0.000000in}}{\pgfqpoint{-0.000000in}{0.000000in}}{%
\pgfpathmoveto{\pgfqpoint{-0.000000in}{0.000000in}}%
\pgfpathlineto{\pgfqpoint{-0.048611in}{0.000000in}}%
\pgfusepath{stroke,fill}%
}%
\begin{pgfscope}%
\pgfsys@transformshift{0.367483in}{0.513891in}%
\pgfsys@useobject{currentmarker}{}%
\end{pgfscope}%
\end{pgfscope}%
\begin{pgfscope}%
\definecolor{textcolor}{rgb}{0.000000,0.000000,0.000000}%
\pgfsetstrokecolor{textcolor}%
\pgfsetfillcolor{textcolor}%
\pgftext[x=0.200816in, y=0.461129in, left, base]{\color{textcolor}{\rmfamily\fontsize{10.000000}{12.000000}\selectfont\catcode`\^=\active\def^{\ifmmode\sp\else\^{}\fi}\catcode`\%=\active\def%{\%}0}}%
\end{pgfscope}%
\begin{pgfscope}%
\pgfpathrectangle{\pgfqpoint{0.367483in}{0.513891in}}{\pgfqpoint{4.228350in}{1.633348in}}%
\pgfusepath{clip}%
\pgfsetrectcap%
\pgfsetroundjoin%
\pgfsetlinewidth{0.803000pt}%
\definecolor{currentstroke}{rgb}{0.690196,0.690196,0.690196}%
\pgfsetstrokecolor{currentstroke}%
\pgfsetstrokeopacity{0.300000}%
\pgfsetdash{}{0pt}%
\pgfpathmoveto{\pgfqpoint{0.367483in}{0.922228in}}%
\pgfpathlineto{\pgfqpoint{4.595833in}{0.922228in}}%
\pgfusepath{stroke}%
\end{pgfscope}%
\begin{pgfscope}%
\pgfsetbuttcap%
\pgfsetroundjoin%
\definecolor{currentfill}{rgb}{0.000000,0.000000,0.000000}%
\pgfsetfillcolor{currentfill}%
\pgfsetlinewidth{0.803000pt}%
\definecolor{currentstroke}{rgb}{0.000000,0.000000,0.000000}%
\pgfsetstrokecolor{currentstroke}%
\pgfsetdash{}{0pt}%
\pgfsys@defobject{currentmarker}{\pgfqpoint{-0.048611in}{0.000000in}}{\pgfqpoint{-0.000000in}{0.000000in}}{%
\pgfpathmoveto{\pgfqpoint{-0.000000in}{0.000000in}}%
\pgfpathlineto{\pgfqpoint{-0.048611in}{0.000000in}}%
\pgfusepath{stroke,fill}%
}%
\begin{pgfscope}%
\pgfsys@transformshift{0.367483in}{0.922228in}%
\pgfsys@useobject{currentmarker}{}%
\end{pgfscope}%
\end{pgfscope}%
\begin{pgfscope}%
\definecolor{textcolor}{rgb}{0.000000,0.000000,0.000000}%
\pgfsetstrokecolor{textcolor}%
\pgfsetfillcolor{textcolor}%
\pgftext[x=0.200816in, y=0.869466in, left, base]{\color{textcolor}{\rmfamily\fontsize{10.000000}{12.000000}\selectfont\catcode`\^=\active\def^{\ifmmode\sp\else\^{}\fi}\catcode`\%=\active\def%{\%}2}}%
\end{pgfscope}%
\begin{pgfscope}%
\pgfpathrectangle{\pgfqpoint{0.367483in}{0.513891in}}{\pgfqpoint{4.228350in}{1.633348in}}%
\pgfusepath{clip}%
\pgfsetrectcap%
\pgfsetroundjoin%
\pgfsetlinewidth{0.803000pt}%
\definecolor{currentstroke}{rgb}{0.690196,0.690196,0.690196}%
\pgfsetstrokecolor{currentstroke}%
\pgfsetstrokeopacity{0.300000}%
\pgfsetdash{}{0pt}%
\pgfpathmoveto{\pgfqpoint{0.367483in}{1.330565in}}%
\pgfpathlineto{\pgfqpoint{4.595833in}{1.330565in}}%
\pgfusepath{stroke}%
\end{pgfscope}%
\begin{pgfscope}%
\pgfsetbuttcap%
\pgfsetroundjoin%
\definecolor{currentfill}{rgb}{0.000000,0.000000,0.000000}%
\pgfsetfillcolor{currentfill}%
\pgfsetlinewidth{0.803000pt}%
\definecolor{currentstroke}{rgb}{0.000000,0.000000,0.000000}%
\pgfsetstrokecolor{currentstroke}%
\pgfsetdash{}{0pt}%
\pgfsys@defobject{currentmarker}{\pgfqpoint{-0.048611in}{0.000000in}}{\pgfqpoint{-0.000000in}{0.000000in}}{%
\pgfpathmoveto{\pgfqpoint{-0.000000in}{0.000000in}}%
\pgfpathlineto{\pgfqpoint{-0.048611in}{0.000000in}}%
\pgfusepath{stroke,fill}%
}%
\begin{pgfscope}%
\pgfsys@transformshift{0.367483in}{1.330565in}%
\pgfsys@useobject{currentmarker}{}%
\end{pgfscope}%
\end{pgfscope}%
\begin{pgfscope}%
\definecolor{textcolor}{rgb}{0.000000,0.000000,0.000000}%
\pgfsetstrokecolor{textcolor}%
\pgfsetfillcolor{textcolor}%
\pgftext[x=0.200816in, y=1.277803in, left, base]{\color{textcolor}{\rmfamily\fontsize{10.000000}{12.000000}\selectfont\catcode`\^=\active\def^{\ifmmode\sp\else\^{}\fi}\catcode`\%=\active\def%{\%}4}}%
\end{pgfscope}%
\begin{pgfscope}%
\pgfpathrectangle{\pgfqpoint{0.367483in}{0.513891in}}{\pgfqpoint{4.228350in}{1.633348in}}%
\pgfusepath{clip}%
\pgfsetrectcap%
\pgfsetroundjoin%
\pgfsetlinewidth{0.803000pt}%
\definecolor{currentstroke}{rgb}{0.690196,0.690196,0.690196}%
\pgfsetstrokecolor{currentstroke}%
\pgfsetstrokeopacity{0.300000}%
\pgfsetdash{}{0pt}%
\pgfpathmoveto{\pgfqpoint{0.367483in}{1.738902in}}%
\pgfpathlineto{\pgfqpoint{4.595833in}{1.738902in}}%
\pgfusepath{stroke}%
\end{pgfscope}%
\begin{pgfscope}%
\pgfsetbuttcap%
\pgfsetroundjoin%
\definecolor{currentfill}{rgb}{0.000000,0.000000,0.000000}%
\pgfsetfillcolor{currentfill}%
\pgfsetlinewidth{0.803000pt}%
\definecolor{currentstroke}{rgb}{0.000000,0.000000,0.000000}%
\pgfsetstrokecolor{currentstroke}%
\pgfsetdash{}{0pt}%
\pgfsys@defobject{currentmarker}{\pgfqpoint{-0.048611in}{0.000000in}}{\pgfqpoint{-0.000000in}{0.000000in}}{%
\pgfpathmoveto{\pgfqpoint{-0.000000in}{0.000000in}}%
\pgfpathlineto{\pgfqpoint{-0.048611in}{0.000000in}}%
\pgfusepath{stroke,fill}%
}%
\begin{pgfscope}%
\pgfsys@transformshift{0.367483in}{1.738902in}%
\pgfsys@useobject{currentmarker}{}%
\end{pgfscope}%
\end{pgfscope}%
\begin{pgfscope}%
\definecolor{textcolor}{rgb}{0.000000,0.000000,0.000000}%
\pgfsetstrokecolor{textcolor}%
\pgfsetfillcolor{textcolor}%
\pgftext[x=0.200816in, y=1.686140in, left, base]{\color{textcolor}{\rmfamily\fontsize{10.000000}{12.000000}\selectfont\catcode`\^=\active\def^{\ifmmode\sp\else\^{}\fi}\catcode`\%=\active\def%{\%}6}}%
\end{pgfscope}%
\begin{pgfscope}%
\pgfpathrectangle{\pgfqpoint{0.367483in}{0.513891in}}{\pgfqpoint{4.228350in}{1.633348in}}%
\pgfusepath{clip}%
\pgfsetrectcap%
\pgfsetroundjoin%
\pgfsetlinewidth{0.803000pt}%
\definecolor{currentstroke}{rgb}{0.690196,0.690196,0.690196}%
\pgfsetstrokecolor{currentstroke}%
\pgfsetstrokeopacity{0.300000}%
\pgfsetdash{}{0pt}%
\pgfpathmoveto{\pgfqpoint{0.367483in}{2.147238in}}%
\pgfpathlineto{\pgfqpoint{4.595833in}{2.147238in}}%
\pgfusepath{stroke}%
\end{pgfscope}%
\begin{pgfscope}%
\pgfsetbuttcap%
\pgfsetroundjoin%
\definecolor{currentfill}{rgb}{0.000000,0.000000,0.000000}%
\pgfsetfillcolor{currentfill}%
\pgfsetlinewidth{0.803000pt}%
\definecolor{currentstroke}{rgb}{0.000000,0.000000,0.000000}%
\pgfsetstrokecolor{currentstroke}%
\pgfsetdash{}{0pt}%
\pgfsys@defobject{currentmarker}{\pgfqpoint{-0.048611in}{0.000000in}}{\pgfqpoint{-0.000000in}{0.000000in}}{%
\pgfpathmoveto{\pgfqpoint{-0.000000in}{0.000000in}}%
\pgfpathlineto{\pgfqpoint{-0.048611in}{0.000000in}}%
\pgfusepath{stroke,fill}%
}%
\begin{pgfscope}%
\pgfsys@transformshift{0.367483in}{2.147238in}%
\pgfsys@useobject{currentmarker}{}%
\end{pgfscope}%
\end{pgfscope}%
\begin{pgfscope}%
\definecolor{textcolor}{rgb}{0.000000,0.000000,0.000000}%
\pgfsetstrokecolor{textcolor}%
\pgfsetfillcolor{textcolor}%
\pgftext[x=0.200816in, y=2.094477in, left, base]{\color{textcolor}{\rmfamily\fontsize{10.000000}{12.000000}\selectfont\catcode`\^=\active\def^{\ifmmode\sp\else\^{}\fi}\catcode`\%=\active\def%{\%}8}}%
\end{pgfscope}%
\begin{pgfscope}%
\definecolor{textcolor}{rgb}{0.000000,0.000000,0.000000}%
\pgfsetstrokecolor{textcolor}%
\pgfsetfillcolor{textcolor}%
\pgftext[x=0.145260in,y=1.330565in,,bottom,rotate=90.000000]{\color{textcolor}{\rmfamily\fontsize{10.000000}{12.000000}\selectfont\catcode`\^=\active\def^{\ifmmode\sp\else\^{}\fi}\catcode`\%=\active\def%{\%}$EF_\mathrm{m,spez,PM10}$ / \unit[per-mode = symbol]{\milli\gram\per\kilo\metre\per\mega\joule}}}%
\end{pgfscope}%
\begin{pgfscope}%
\pgfpathrectangle{\pgfqpoint{0.367483in}{0.513891in}}{\pgfqpoint{4.228350in}{1.633348in}}%
\pgfusepath{clip}%
\pgfsetbuttcap%
\pgfsetroundjoin%
\pgfsetlinewidth{1.003750pt}%
\definecolor{currentstroke}{rgb}{0.121569,0.466667,0.705882}%
\pgfsetstrokecolor{currentstroke}%
\pgfsetdash{{3.700000pt}{1.600000pt}}{0.000000pt}%
\pgfpathmoveto{\pgfqpoint{0.601875in}{1.805799in}}%
\pgfpathlineto{\pgfqpoint{0.615233in}{1.742620in}}%
\pgfpathlineto{\pgfqpoint{0.628591in}{1.685613in}}%
\pgfpathlineto{\pgfqpoint{0.641949in}{1.633903in}}%
\pgfpathlineto{\pgfqpoint{0.655306in}{1.586775in}}%
\pgfpathlineto{\pgfqpoint{0.668664in}{1.543636in}}%
\pgfpathlineto{\pgfqpoint{0.682022in}{1.503993in}}%
\pgfpathlineto{\pgfqpoint{0.695379in}{1.467433in}}%
\pgfpathlineto{\pgfqpoint{0.708737in}{1.433603in}}%
\pgfpathlineto{\pgfqpoint{0.722095in}{1.402205in}}%
\pgfpathlineto{\pgfqpoint{0.735453in}{1.372982in}}%
\pgfpathlineto{\pgfqpoint{0.748810in}{1.345712in}}%
\pgfpathlineto{\pgfqpoint{0.762168in}{1.320202in}}%
\pgfpathlineto{\pgfqpoint{0.775526in}{1.296286in}}%
\pgfpathlineto{\pgfqpoint{0.788883in}{1.273816in}}%
\pgfpathlineto{\pgfqpoint{0.802241in}{1.252662in}}%
\pgfpathlineto{\pgfqpoint{0.815599in}{1.232712in}}%
\pgfpathlineto{\pgfqpoint{0.828957in}{1.213863in}}%
\pgfpathlineto{\pgfqpoint{0.842314in}{1.196025in}}%
\pgfpathlineto{\pgfqpoint{0.869030in}{1.163069in}}%
\pgfpathlineto{\pgfqpoint{0.895745in}{1.133296in}}%
\pgfpathlineto{\pgfqpoint{0.922461in}{1.106260in}}%
\pgfpathlineto{\pgfqpoint{0.949176in}{1.081593in}}%
\pgfpathlineto{\pgfqpoint{0.975891in}{1.058994in}}%
\pgfpathlineto{\pgfqpoint{1.002607in}{1.038209in}}%
\pgfpathlineto{\pgfqpoint{1.029322in}{1.019025in}}%
\pgfpathlineto{\pgfqpoint{1.056038in}{1.001262in}}%
\pgfpathlineto{\pgfqpoint{1.082753in}{0.984765in}}%
\pgfpathlineto{\pgfqpoint{1.109469in}{0.969400in}}%
\pgfpathlineto{\pgfqpoint{1.136184in}{0.955055in}}%
\pgfpathlineto{\pgfqpoint{1.162900in}{0.941629in}}%
\pgfpathlineto{\pgfqpoint{1.202973in}{0.923026in}}%
\pgfpathlineto{\pgfqpoint{1.243046in}{0.906048in}}%
\pgfpathlineto{\pgfqpoint{1.283119in}{0.890489in}}%
\pgfpathlineto{\pgfqpoint{1.323192in}{0.876175in}}%
\pgfpathlineto{\pgfqpoint{1.363265in}{0.862960in}}%
\pgfpathlineto{\pgfqpoint{1.403338in}{0.850721in}}%
\pgfpathlineto{\pgfqpoint{1.456769in}{0.835738in}}%
\pgfpathlineto{\pgfqpoint{1.510200in}{0.822094in}}%
\pgfpathlineto{\pgfqpoint{1.563631in}{0.809612in}}%
\pgfpathlineto{\pgfqpoint{1.617062in}{0.798149in}}%
\pgfpathlineto{\pgfqpoint{1.683851in}{0.785070in}}%
\pgfpathlineto{\pgfqpoint{1.750639in}{0.773196in}}%
\pgfpathlineto{\pgfqpoint{1.817428in}{0.762367in}}%
\pgfpathlineto{\pgfqpoint{1.897574in}{0.750562in}}%
\pgfpathlineto{\pgfqpoint{1.977720in}{0.739880in}}%
\pgfpathlineto{\pgfqpoint{2.071224in}{0.728630in}}%
\pgfpathlineto{\pgfqpoint{2.164728in}{0.718498in}}%
\pgfpathlineto{\pgfqpoint{2.271590in}{0.708081in}}%
\pgfpathlineto{\pgfqpoint{2.391810in}{0.697618in}}%
\pgfpathlineto{\pgfqpoint{2.512029in}{0.688275in}}%
\pgfpathlineto{\pgfqpoint{2.645606in}{0.678999in}}%
\pgfpathlineto{\pgfqpoint{2.792541in}{0.669922in}}%
\pgfpathlineto{\pgfqpoint{2.952834in}{0.661144in}}%
\pgfpathlineto{\pgfqpoint{3.126484in}{0.652734in}}%
\pgfpathlineto{\pgfqpoint{3.326850in}{0.644202in}}%
\pgfpathlineto{\pgfqpoint{3.540573in}{0.636235in}}%
\pgfpathlineto{\pgfqpoint{3.781012in}{0.628412in}}%
\pgfpathlineto{\pgfqpoint{4.048167in}{0.620866in}}%
\pgfpathlineto{\pgfqpoint{4.342036in}{0.613685in}}%
\pgfpathlineto{\pgfqpoint{4.595833in}{0.608250in}}%
\pgfpathlineto{\pgfqpoint{4.595833in}{0.608250in}}%
\pgfusepath{stroke}%
\end{pgfscope}%
\begin{pgfscope}%
\pgfpathrectangle{\pgfqpoint{0.367483in}{0.513891in}}{\pgfqpoint{4.228350in}{1.633348in}}%
\pgfusepath{clip}%
\pgfsetbuttcap%
\pgfsetroundjoin%
\pgfsetlinewidth{1.003750pt}%
\definecolor{currentstroke}{rgb}{1.000000,0.498039,0.054902}%
\pgfsetstrokecolor{currentstroke}%
\pgfsetdash{{3.700000pt}{1.600000pt}}{0.000000pt}%
\pgfpathmoveto{\pgfqpoint{0.630371in}{2.050420in}}%
\pgfpathlineto{\pgfqpoint{0.643634in}{1.982618in}}%
\pgfpathlineto{\pgfqpoint{0.656896in}{1.920783in}}%
\pgfpathlineto{\pgfqpoint{0.670158in}{1.864152in}}%
\pgfpathlineto{\pgfqpoint{0.683421in}{1.812087in}}%
\pgfpathlineto{\pgfqpoint{0.696683in}{1.764049in}}%
\pgfpathlineto{\pgfqpoint{0.709946in}{1.719583in}}%
\pgfpathlineto{\pgfqpoint{0.723208in}{1.678300in}}%
\pgfpathlineto{\pgfqpoint{0.736470in}{1.639866in}}%
\pgfpathlineto{\pgfqpoint{0.749733in}{1.603991in}}%
\pgfpathlineto{\pgfqpoint{0.762995in}{1.570426in}}%
\pgfpathlineto{\pgfqpoint{0.776258in}{1.538952in}}%
\pgfpathlineto{\pgfqpoint{0.789520in}{1.509376in}}%
\pgfpathlineto{\pgfqpoint{0.802783in}{1.481529in}}%
\pgfpathlineto{\pgfqpoint{0.816045in}{1.455263in}}%
\pgfpathlineto{\pgfqpoint{0.829307in}{1.430445in}}%
\pgfpathlineto{\pgfqpoint{0.842570in}{1.406956in}}%
\pgfpathlineto{\pgfqpoint{0.855832in}{1.384691in}}%
\pgfpathlineto{\pgfqpoint{0.882357in}{1.343466in}}%
\pgfpathlineto{\pgfqpoint{0.908882in}{1.306121in}}%
\pgfpathlineto{\pgfqpoint{0.935407in}{1.272126in}}%
\pgfpathlineto{\pgfqpoint{0.961931in}{1.241044in}}%
\pgfpathlineto{\pgfqpoint{0.988456in}{1.212512in}}%
\pgfpathlineto{\pgfqpoint{1.014981in}{1.186225in}}%
\pgfpathlineto{\pgfqpoint{1.041506in}{1.161923in}}%
\pgfpathlineto{\pgfqpoint{1.068031in}{1.139389in}}%
\pgfpathlineto{\pgfqpoint{1.094556in}{1.118433in}}%
\pgfpathlineto{\pgfqpoint{1.121080in}{1.098893in}}%
\pgfpathlineto{\pgfqpoint{1.147605in}{1.080628in}}%
\pgfpathlineto{\pgfqpoint{1.174130in}{1.063516in}}%
\pgfpathlineto{\pgfqpoint{1.200655in}{1.047450in}}%
\pgfpathlineto{\pgfqpoint{1.227180in}{1.032335in}}%
\pgfpathlineto{\pgfqpoint{1.266967in}{1.011268in}}%
\pgfpathlineto{\pgfqpoint{1.306754in}{0.991915in}}%
\pgfpathlineto{\pgfqpoint{1.346542in}{0.974071in}}%
\pgfpathlineto{\pgfqpoint{1.386329in}{0.957566in}}%
\pgfpathlineto{\pgfqpoint{1.426116in}{0.942252in}}%
\pgfpathlineto{\pgfqpoint{1.465903in}{0.928003in}}%
\pgfpathlineto{\pgfqpoint{1.505690in}{0.914710in}}%
\pgfpathlineto{\pgfqpoint{1.558740in}{0.898312in}}%
\pgfpathlineto{\pgfqpoint{1.611790in}{0.883256in}}%
\pgfpathlineto{\pgfqpoint{1.664839in}{0.869383in}}%
\pgfpathlineto{\pgfqpoint{1.717889in}{0.856557in}}%
\pgfpathlineto{\pgfqpoint{1.784201in}{0.841821in}}%
\pgfpathlineto{\pgfqpoint{1.850513in}{0.828350in}}%
\pgfpathlineto{\pgfqpoint{1.916825in}{0.815987in}}%
\pgfpathlineto{\pgfqpoint{1.996400in}{0.802427in}}%
\pgfpathlineto{\pgfqpoint{2.075974in}{0.790080in}}%
\pgfpathlineto{\pgfqpoint{2.155549in}{0.778788in}}%
\pgfpathlineto{\pgfqpoint{2.248386in}{0.766774in}}%
\pgfpathlineto{\pgfqpoint{2.341223in}{0.755846in}}%
\pgfpathlineto{\pgfqpoint{2.447322in}{0.744503in}}%
\pgfpathlineto{\pgfqpoint{2.553421in}{0.734219in}}%
\pgfpathlineto{\pgfqpoint{2.672783in}{0.723735in}}%
\pgfpathlineto{\pgfqpoint{2.805407in}{0.713243in}}%
\pgfpathlineto{\pgfqpoint{2.951294in}{0.702898in}}%
\pgfpathlineto{\pgfqpoint{3.110443in}{0.692817in}}%
\pgfpathlineto{\pgfqpoint{3.282854in}{0.683090in}}%
\pgfpathlineto{\pgfqpoint{3.468528in}{0.673777in}}%
\pgfpathlineto{\pgfqpoint{3.667464in}{0.664917in}}%
\pgfpathlineto{\pgfqpoint{3.892925in}{0.656036in}}%
\pgfpathlineto{\pgfqpoint{4.131648in}{0.647748in}}%
\pgfpathlineto{\pgfqpoint{4.396897in}{0.639645in}}%
\pgfpathlineto{\pgfqpoint{4.595833in}{0.634209in}}%
\pgfpathlineto{\pgfqpoint{4.595833in}{0.634209in}}%
\pgfusepath{stroke}%
\end{pgfscope}%
\begin{pgfscope}%
\pgfpathrectangle{\pgfqpoint{0.367483in}{0.513891in}}{\pgfqpoint{4.228350in}{1.633348in}}%
\pgfusepath{clip}%
\pgfsetbuttcap%
\pgfsetroundjoin%
\pgfsetlinewidth{1.003750pt}%
\definecolor{currentstroke}{rgb}{0.172549,0.627451,0.172549}%
\pgfsetstrokecolor{currentstroke}%
\pgfsetdash{{3.700000pt}{1.600000pt}}{0.000000pt}%
\pgfpathmoveto{\pgfqpoint{0.582603in}{1.818830in}}%
\pgfpathlineto{\pgfqpoint{0.596025in}{1.760706in}}%
\pgfpathlineto{\pgfqpoint{0.609447in}{1.708274in}}%
\pgfpathlineto{\pgfqpoint{0.622869in}{1.660704in}}%
\pgfpathlineto{\pgfqpoint{0.636291in}{1.617324in}}%
\pgfpathlineto{\pgfqpoint{0.649713in}{1.577582in}}%
\pgfpathlineto{\pgfqpoint{0.663136in}{1.541021in}}%
\pgfpathlineto{\pgfqpoint{0.676558in}{1.507258in}}%
\pgfpathlineto{\pgfqpoint{0.689980in}{1.475971in}}%
\pgfpathlineto{\pgfqpoint{0.703402in}{1.446887in}}%
\pgfpathlineto{\pgfqpoint{0.716824in}{1.419770in}}%
\pgfpathlineto{\pgfqpoint{0.730247in}{1.394421in}}%
\pgfpathlineto{\pgfqpoint{0.743669in}{1.370664in}}%
\pgfpathlineto{\pgfqpoint{0.757091in}{1.348348in}}%
\pgfpathlineto{\pgfqpoint{0.770513in}{1.327340in}}%
\pgfpathlineto{\pgfqpoint{0.783935in}{1.307523in}}%
\pgfpathlineto{\pgfqpoint{0.797357in}{1.288796in}}%
\pgfpathlineto{\pgfqpoint{0.824202in}{1.254252in}}%
\pgfpathlineto{\pgfqpoint{0.851046in}{1.223094in}}%
\pgfpathlineto{\pgfqpoint{0.877890in}{1.194828in}}%
\pgfpathlineto{\pgfqpoint{0.904735in}{1.169053in}}%
\pgfpathlineto{\pgfqpoint{0.931579in}{1.145441in}}%
\pgfpathlineto{\pgfqpoint{0.958423in}{1.123720in}}%
\pgfpathlineto{\pgfqpoint{0.985268in}{1.103662in}}%
\pgfpathlineto{\pgfqpoint{1.012112in}{1.085076in}}%
\pgfpathlineto{\pgfqpoint{1.038957in}{1.067799in}}%
\pgfpathlineto{\pgfqpoint{1.065801in}{1.051692in}}%
\pgfpathlineto{\pgfqpoint{1.092645in}{1.036635in}}%
\pgfpathlineto{\pgfqpoint{1.119490in}{1.022524in}}%
\pgfpathlineto{\pgfqpoint{1.159756in}{1.002939in}}%
\pgfpathlineto{\pgfqpoint{1.200023in}{0.985023in}}%
\pgfpathlineto{\pgfqpoint{1.240289in}{0.968566in}}%
\pgfpathlineto{\pgfqpoint{1.280556in}{0.953387in}}%
\pgfpathlineto{\pgfqpoint{1.320822in}{0.939338in}}%
\pgfpathlineto{\pgfqpoint{1.361089in}{0.926292in}}%
\pgfpathlineto{\pgfqpoint{1.414777in}{0.910274in}}%
\pgfpathlineto{\pgfqpoint{1.468466in}{0.895633in}}%
\pgfpathlineto{\pgfqpoint{1.522155in}{0.882192in}}%
\pgfpathlineto{\pgfqpoint{1.575844in}{0.869804in}}%
\pgfpathlineto{\pgfqpoint{1.642954in}{0.855613in}}%
\pgfpathlineto{\pgfqpoint{1.710065in}{0.842672in}}%
\pgfpathlineto{\pgfqpoint{1.777176in}{0.830818in}}%
\pgfpathlineto{\pgfqpoint{1.857709in}{0.817837in}}%
\pgfpathlineto{\pgfqpoint{1.938242in}{0.806029in}}%
\pgfpathlineto{\pgfqpoint{2.032198in}{0.793528in}}%
\pgfpathlineto{\pgfqpoint{2.126153in}{0.782206in}}%
\pgfpathlineto{\pgfqpoint{2.233530in}{0.770498in}}%
\pgfpathlineto{\pgfqpoint{2.340908in}{0.759915in}}%
\pgfpathlineto{\pgfqpoint{2.461707in}{0.749153in}}%
\pgfpathlineto{\pgfqpoint{2.595929in}{0.738404in}}%
\pgfpathlineto{\pgfqpoint{2.743573in}{0.727819in}}%
\pgfpathlineto{\pgfqpoint{2.891217in}{0.718328in}}%
\pgfpathlineto{\pgfqpoint{3.052283in}{0.709024in}}%
\pgfpathlineto{\pgfqpoint{3.240193in}{0.699336in}}%
\pgfpathlineto{\pgfqpoint{3.441526in}{0.690116in}}%
\pgfpathlineto{\pgfqpoint{3.656281in}{0.681381in}}%
\pgfpathlineto{\pgfqpoint{3.897880in}{0.672677in}}%
\pgfpathlineto{\pgfqpoint{4.166323in}{0.664154in}}%
\pgfpathlineto{\pgfqpoint{4.461611in}{0.655921in}}%
\pgfpathlineto{\pgfqpoint{4.595833in}{0.652513in}}%
\pgfpathlineto{\pgfqpoint{4.595833in}{0.652513in}}%
\pgfusepath{stroke}%
\end{pgfscope}%
\begin{pgfscope}%
\pgfsetrectcap%
\pgfsetmiterjoin%
\pgfsetlinewidth{0.803000pt}%
\definecolor{currentstroke}{rgb}{0.000000,0.000000,0.000000}%
\pgfsetstrokecolor{currentstroke}%
\pgfsetdash{}{0pt}%
\pgfpathmoveto{\pgfqpoint{0.367483in}{0.513891in}}%
\pgfpathlineto{\pgfqpoint{0.367483in}{2.147238in}}%
\pgfusepath{stroke}%
\end{pgfscope}%
\begin{pgfscope}%
\pgfsetrectcap%
\pgfsetmiterjoin%
\pgfsetlinewidth{0.803000pt}%
\definecolor{currentstroke}{rgb}{0.000000,0.000000,0.000000}%
\pgfsetstrokecolor{currentstroke}%
\pgfsetdash{}{0pt}%
\pgfpathmoveto{\pgfqpoint{4.595833in}{0.513891in}}%
\pgfpathlineto{\pgfqpoint{4.595833in}{2.147238in}}%
\pgfusepath{stroke}%
\end{pgfscope}%
\begin{pgfscope}%
\pgfsetrectcap%
\pgfsetmiterjoin%
\pgfsetlinewidth{0.803000pt}%
\definecolor{currentstroke}{rgb}{0.000000,0.000000,0.000000}%
\pgfsetstrokecolor{currentstroke}%
\pgfsetdash{}{0pt}%
\pgfpathmoveto{\pgfqpoint{0.367483in}{0.513891in}}%
\pgfpathlineto{\pgfqpoint{4.595833in}{0.513891in}}%
\pgfusepath{stroke}%
\end{pgfscope}%
\begin{pgfscope}%
\pgfsetrectcap%
\pgfsetmiterjoin%
\pgfsetlinewidth{0.803000pt}%
\definecolor{currentstroke}{rgb}{0.000000,0.000000,0.000000}%
\pgfsetstrokecolor{currentstroke}%
\pgfsetdash{}{0pt}%
\pgfpathmoveto{\pgfqpoint{0.367483in}{2.147238in}}%
\pgfpathlineto{\pgfqpoint{4.595833in}{2.147238in}}%
\pgfusepath{stroke}%
\end{pgfscope}%
\begin{pgfscope}%
\pgfpathrectangle{\pgfqpoint{0.367483in}{0.513891in}}{\pgfqpoint{4.228350in}{1.633348in}}%
\pgfusepath{clip}%
\pgfsetbuttcap%
\pgfsetroundjoin%
\definecolor{currentfill}{rgb}{0.121569,0.466667,0.705882}%
\pgfsetfillcolor{currentfill}%
\pgfsetlinewidth{1.003750pt}%
\definecolor{currentstroke}{rgb}{0.121569,0.466667,0.705882}%
\pgfsetstrokecolor{currentstroke}%
\pgfsetdash{}{0pt}%
\pgfsys@defobject{currentmarker}{\pgfqpoint{-0.021960in}{-0.021960in}}{\pgfqpoint{0.021960in}{0.021960in}}{%
\pgfpathmoveto{\pgfqpoint{0.000000in}{-0.021960in}}%
\pgfpathcurveto{\pgfqpoint{0.005824in}{-0.021960in}}{\pgfqpoint{0.011410in}{-0.019646in}}{\pgfqpoint{0.015528in}{-0.015528in}}%
\pgfpathcurveto{\pgfqpoint{0.019646in}{-0.011410in}}{\pgfqpoint{0.021960in}{-0.005824in}}{\pgfqpoint{0.021960in}{0.000000in}}%
\pgfpathcurveto{\pgfqpoint{0.021960in}{0.005824in}}{\pgfqpoint{0.019646in}{0.011410in}}{\pgfqpoint{0.015528in}{0.015528in}}%
\pgfpathcurveto{\pgfqpoint{0.011410in}{0.019646in}}{\pgfqpoint{0.005824in}{0.021960in}}{\pgfqpoint{0.000000in}{0.021960in}}%
\pgfpathcurveto{\pgfqpoint{-0.005824in}{0.021960in}}{\pgfqpoint{-0.011410in}{0.019646in}}{\pgfqpoint{-0.015528in}{0.015528in}}%
\pgfpathcurveto{\pgfqpoint{-0.019646in}{0.011410in}}{\pgfqpoint{-0.021960in}{0.005824in}}{\pgfqpoint{-0.021960in}{0.000000in}}%
\pgfpathcurveto{\pgfqpoint{-0.021960in}{-0.005824in}}{\pgfqpoint{-0.019646in}{-0.011410in}}{\pgfqpoint{-0.015528in}{-0.015528in}}%
\pgfpathcurveto{\pgfqpoint{-0.011410in}{-0.019646in}}{\pgfqpoint{-0.005824in}{-0.021960in}}{\pgfqpoint{0.000000in}{-0.021960in}}%
\pgfpathlineto{\pgfqpoint{0.000000in}{-0.021960in}}%
\pgfpathclose%
\pgfusepath{stroke,fill}%
}%
\begin{pgfscope}%
\pgfsys@transformshift{0.627919in}{1.684257in}%
\pgfsys@useobject{currentmarker}{}%
\end{pgfscope}%
\begin{pgfscope}%
\pgfsys@transformshift{1.558326in}{0.871112in}%
\pgfsys@useobject{currentmarker}{}%
\end{pgfscope}%
\begin{pgfscope}%
\pgfsys@transformshift{1.949330in}{0.772531in}%
\pgfsys@useobject{currentmarker}{}%
\end{pgfscope}%
\begin{pgfscope}%
\pgfsys@transformshift{2.126166in}{0.661511in}%
\pgfsys@useobject{currentmarker}{}%
\end{pgfscope}%
\begin{pgfscope}%
\pgfsys@transformshift{2.394588in}{0.662111in}%
\pgfsys@useobject{currentmarker}{}%
\end{pgfscope}%
\begin{pgfscope}%
\pgfsys@transformshift{2.693169in}{0.663779in}%
\pgfsys@useobject{currentmarker}{}%
\end{pgfscope}%
\begin{pgfscope}%
\pgfsys@transformshift{2.975994in}{0.691669in}%
\pgfsys@useobject{currentmarker}{}%
\end{pgfscope}%
\begin{pgfscope}%
\pgfsys@transformshift{3.403733in}{0.616227in}%
\pgfsys@useobject{currentmarker}{}%
\end{pgfscope}%
\begin{pgfscope}%
\pgfsys@transformshift{3.591874in}{0.643181in}%
\pgfsys@useobject{currentmarker}{}%
\end{pgfscope}%
\begin{pgfscope}%
\pgfsys@transformshift{3.876160in}{0.615953in}%
\pgfsys@useobject{currentmarker}{}%
\end{pgfscope}%
\end{pgfscope}%
\begin{pgfscope}%
\pgfpathrectangle{\pgfqpoint{0.367483in}{0.513891in}}{\pgfqpoint{4.228350in}{1.633348in}}%
\pgfusepath{clip}%
\pgfsetbuttcap%
\pgfsetroundjoin%
\definecolor{currentfill}{rgb}{1.000000,0.498039,0.054902}%
\pgfsetfillcolor{currentfill}%
\pgfsetlinewidth{1.003750pt}%
\definecolor{currentstroke}{rgb}{1.000000,0.498039,0.054902}%
\pgfsetstrokecolor{currentstroke}%
\pgfsetdash{}{0pt}%
\pgfsys@defobject{currentmarker}{\pgfqpoint{-0.021960in}{-0.021960in}}{\pgfqpoint{0.021960in}{0.021960in}}{%
\pgfpathmoveto{\pgfqpoint{-0.021960in}{-0.021960in}}%
\pgfpathlineto{\pgfqpoint{0.021960in}{-0.021960in}}%
\pgfpathlineto{\pgfqpoint{0.021960in}{0.021960in}}%
\pgfpathlineto{\pgfqpoint{-0.021960in}{0.021960in}}%
\pgfpathlineto{\pgfqpoint{-0.021960in}{-0.021960in}}%
\pgfpathclose%
\pgfusepath{stroke,fill}%
}%
\begin{pgfscope}%
\pgfsys@transformshift{0.659581in}{1.908491in}%
\pgfsys@useobject{currentmarker}{}%
\end{pgfscope}%
\begin{pgfscope}%
\pgfsys@transformshift{1.778424in}{0.866622in}%
\pgfsys@useobject{currentmarker}{}%
\end{pgfscope}%
\begin{pgfscope}%
\pgfsys@transformshift{2.149519in}{0.738439in}%
\pgfsys@useobject{currentmarker}{}%
\end{pgfscope}%
\begin{pgfscope}%
\pgfsys@transformshift{2.528488in}{0.753525in}%
\pgfsys@useobject{currentmarker}{}%
\end{pgfscope}%
\end{pgfscope}%
\begin{pgfscope}%
\pgfpathrectangle{\pgfqpoint{0.367483in}{0.513891in}}{\pgfqpoint{4.228350in}{1.633348in}}%
\pgfusepath{clip}%
\pgfsetbuttcap%
\pgfsetroundjoin%
\definecolor{currentfill}{rgb}{0.172549,0.627451,0.172549}%
\pgfsetfillcolor{currentfill}%
\pgfsetlinewidth{1.003750pt}%
\definecolor{currentstroke}{rgb}{0.172549,0.627451,0.172549}%
\pgfsetstrokecolor{currentstroke}%
\pgfsetdash{}{0pt}%
\pgfsys@defobject{currentmarker}{\pgfqpoint{-0.021960in}{-0.021960in}}{\pgfqpoint{0.021960in}{0.021960in}}{%
\pgfpathmoveto{\pgfqpoint{0.000000in}{0.021960in}}%
\pgfpathlineto{\pgfqpoint{-0.021960in}{-0.021960in}}%
\pgfpathlineto{\pgfqpoint{0.021960in}{-0.021960in}}%
\pgfpathlineto{\pgfqpoint{0.000000in}{0.021960in}}%
\pgfpathclose%
\pgfusepath{stroke,fill}%
}%
\begin{pgfscope}%
\pgfsys@transformshift{0.606505in}{1.715953in}%
\pgfsys@useobject{currentmarker}{}%
\end{pgfscope}%
\begin{pgfscope}%
\pgfsys@transformshift{1.424489in}{0.975164in}%
\pgfsys@useobject{currentmarker}{}%
\end{pgfscope}%
\begin{pgfscope}%
\pgfsys@transformshift{1.715586in}{0.800875in}%
\pgfsys@useobject{currentmarker}{}%
\end{pgfscope}%
\begin{pgfscope}%
\pgfsys@transformshift{2.015482in}{0.764677in}%
\pgfsys@useobject{currentmarker}{}%
\end{pgfscope}%
\begin{pgfscope}%
\pgfsys@transformshift{2.310481in}{0.652380in}%
\pgfsys@useobject{currentmarker}{}%
\end{pgfscope}%
\begin{pgfscope}%
\pgfsys@transformshift{2.525084in}{0.896388in}%
\pgfsys@useobject{currentmarker}{}%
\end{pgfscope}%
\begin{pgfscope}%
\pgfsys@transformshift{2.839834in}{0.716063in}%
\pgfsys@useobject{currentmarker}{}%
\end{pgfscope}%
\begin{pgfscope}%
\pgfsys@transformshift{3.130930in}{0.648272in}%
\pgfsys@useobject{currentmarker}{}%
\end{pgfscope}%
\begin{pgfscope}%
\pgfsys@transformshift{3.430827in}{0.712579in}%
\pgfsys@useobject{currentmarker}{}%
\end{pgfscope}%
\end{pgfscope}%
\begin{pgfscope}%
\pgfsetbuttcap%
\pgfsetmiterjoin%
\definecolor{currentfill}{rgb}{1.000000,1.000000,1.000000}%
\pgfsetfillcolor{currentfill}%
\pgfsetlinewidth{0.000000pt}%
\definecolor{currentstroke}{rgb}{0.000000,0.000000,0.000000}%
\pgfsetstrokecolor{currentstroke}%
\pgfsetstrokeopacity{0.000000}%
\pgfsetdash{}{0pt}%
\pgfpathmoveto{\pgfqpoint{2.375949in}{1.085562in}}%
\pgfpathlineto{\pgfqpoint{4.490124in}{1.085562in}}%
\pgfpathlineto{\pgfqpoint{4.490124in}{1.575567in}}%
\pgfpathlineto{\pgfqpoint{2.375949in}{1.575567in}}%
\pgfpathlineto{\pgfqpoint{2.375949in}{1.085562in}}%
\pgfpathclose%
\pgfusepath{fill}%
\end{pgfscope}%
\begin{pgfscope}%
\pgfsetbuttcap%
\pgfsetmiterjoin%
\pgfsetlinewidth{0.803000pt}%
\definecolor{currentstroke}{rgb}{0.400000,0.400000,0.400000}%
\pgfsetstrokecolor{currentstroke}%
\pgfsetdash{{0.800000pt}{1.320000pt}}{0.000000pt}%
\pgfpathmoveto{\pgfqpoint{2.375949in}{1.575567in}}%
\pgfpathlineto{\pgfqpoint{1.847405in}{0.922228in}}%
\pgfusepath{stroke}%
\end{pgfscope}%
\begin{pgfscope}%
\pgfsetbuttcap%
\pgfsetmiterjoin%
\pgfsetlinewidth{0.803000pt}%
\definecolor{currentstroke}{rgb}{0.400000,0.400000,0.400000}%
\pgfsetstrokecolor{currentstroke}%
\pgfsetdash{{0.800000pt}{1.320000pt}}{0.000000pt}%
\pgfpathmoveto{\pgfqpoint{4.490124in}{1.085562in}}%
\pgfpathlineto{\pgfqpoint{2.270240in}{0.615975in}}%
\pgfusepath{stroke}%
\end{pgfscope}%
\begin{pgfscope}%
\pgfpathrectangle{\pgfqpoint{2.375949in}{1.085562in}}{\pgfqpoint{2.114175in}{0.490004in}}%
\pgfusepath{clip}%
\pgfsetbuttcap%
\pgfsetroundjoin%
\pgfsetlinewidth{0.501875pt}%
\definecolor{currentstroke}{rgb}{0.690196,0.690196,0.690196}%
\pgfsetstrokecolor{currentstroke}%
\pgfsetdash{{0.500000pt}{0.825000pt}}{0.000000pt}%
\pgfpathmoveto{\pgfqpoint{2.375949in}{1.085562in}}%
\pgfpathlineto{\pgfqpoint{2.375949in}{1.575567in}}%
\pgfusepath{stroke}%
\end{pgfscope}%
\begin{pgfscope}%
\pgfsetbuttcap%
\pgfsetroundjoin%
\definecolor{currentfill}{rgb}{0.000000,0.000000,0.000000}%
\pgfsetfillcolor{currentfill}%
\pgfsetlinewidth{0.803000pt}%
\definecolor{currentstroke}{rgb}{0.000000,0.000000,0.000000}%
\pgfsetstrokecolor{currentstroke}%
\pgfsetdash{}{0pt}%
\pgfsys@defobject{currentmarker}{\pgfqpoint{0.000000in}{-0.048611in}}{\pgfqpoint{0.000000in}{0.000000in}}{%
\pgfpathmoveto{\pgfqpoint{0.000000in}{0.000000in}}%
\pgfpathlineto{\pgfqpoint{0.000000in}{-0.048611in}}%
\pgfusepath{stroke,fill}%
}%
\begin{pgfscope}%
\pgfsys@transformshift{2.375949in}{1.085562in}%
\pgfsys@useobject{currentmarker}{}%
\end{pgfscope}%
\end{pgfscope}%
\begin{pgfscope}%
\pgfsetbuttcap%
\pgfsetmiterjoin%
\definecolor{currentfill}{rgb}{1.000000,1.000000,1.000000}%
\pgfsetfillcolor{currentfill}%
\pgfsetfillopacity{0.700000}%
\pgfsetlinewidth{0.000000pt}%
\definecolor{currentstroke}{rgb}{0.000000,0.000000,0.000000}%
\pgfsetstrokecolor{currentstroke}%
\pgfsetstrokeopacity{0.700000}%
\pgfsetdash{}{0pt}%
\pgfpathmoveto{\pgfqpoint{2.291228in}{0.896673in}}%
\pgfpathlineto{\pgfqpoint{2.460670in}{0.896673in}}%
\pgfpathlineto{\pgfqpoint{2.460670in}{0.996673in}}%
\pgfpathlineto{\pgfqpoint{2.291228in}{0.996673in}}%
\pgfpathlineto{\pgfqpoint{2.291228in}{0.896673in}}%
\pgfpathclose%
\pgfusepath{fill}%
\end{pgfscope}%
\begin{pgfscope}%
\definecolor{textcolor}{rgb}{0.000000,0.000000,0.000000}%
\pgfsetstrokecolor{textcolor}%
\pgfsetfillcolor{textcolor}%
\pgftext[x=2.375949in,y=0.988340in,,top]{\color{textcolor}{\rmfamily\fontsize{6.000000}{7.200000}\selectfont\catcode`\^=\active\def^{\ifmmode\sp\else\^{}\fi}\catcode`\%=\active\def%{\%}140}}%
\end{pgfscope}%
\begin{pgfscope}%
\pgfpathrectangle{\pgfqpoint{2.375949in}{1.085562in}}{\pgfqpoint{2.114175in}{0.490004in}}%
\pgfusepath{clip}%
\pgfsetbuttcap%
\pgfsetroundjoin%
\pgfsetlinewidth{0.501875pt}%
\definecolor{currentstroke}{rgb}{0.690196,0.690196,0.690196}%
\pgfsetstrokecolor{currentstroke}%
\pgfsetdash{{0.500000pt}{0.825000pt}}{0.000000pt}%
\pgfpathmoveto{\pgfqpoint{2.904493in}{1.085562in}}%
\pgfpathlineto{\pgfqpoint{2.904493in}{1.575567in}}%
\pgfusepath{stroke}%
\end{pgfscope}%
\begin{pgfscope}%
\pgfsetbuttcap%
\pgfsetroundjoin%
\definecolor{currentfill}{rgb}{0.000000,0.000000,0.000000}%
\pgfsetfillcolor{currentfill}%
\pgfsetlinewidth{0.803000pt}%
\definecolor{currentstroke}{rgb}{0.000000,0.000000,0.000000}%
\pgfsetstrokecolor{currentstroke}%
\pgfsetdash{}{0pt}%
\pgfsys@defobject{currentmarker}{\pgfqpoint{0.000000in}{-0.048611in}}{\pgfqpoint{0.000000in}{0.000000in}}{%
\pgfpathmoveto{\pgfqpoint{0.000000in}{0.000000in}}%
\pgfpathlineto{\pgfqpoint{0.000000in}{-0.048611in}}%
\pgfusepath{stroke,fill}%
}%
\begin{pgfscope}%
\pgfsys@transformshift{2.904493in}{1.085562in}%
\pgfsys@useobject{currentmarker}{}%
\end{pgfscope}%
\end{pgfscope}%
\begin{pgfscope}%
\pgfsetbuttcap%
\pgfsetmiterjoin%
\definecolor{currentfill}{rgb}{1.000000,1.000000,1.000000}%
\pgfsetfillcolor{currentfill}%
\pgfsetfillopacity{0.700000}%
\pgfsetlinewidth{0.000000pt}%
\definecolor{currentstroke}{rgb}{0.000000,0.000000,0.000000}%
\pgfsetstrokecolor{currentstroke}%
\pgfsetstrokeopacity{0.700000}%
\pgfsetdash{}{0pt}%
\pgfpathmoveto{\pgfqpoint{2.819772in}{0.896673in}}%
\pgfpathlineto{\pgfqpoint{2.989214in}{0.896673in}}%
\pgfpathlineto{\pgfqpoint{2.989214in}{0.996673in}}%
\pgfpathlineto{\pgfqpoint{2.819772in}{0.996673in}}%
\pgfpathlineto{\pgfqpoint{2.819772in}{0.896673in}}%
\pgfpathclose%
\pgfusepath{fill}%
\end{pgfscope}%
\begin{pgfscope}%
\definecolor{textcolor}{rgb}{0.000000,0.000000,0.000000}%
\pgfsetstrokecolor{textcolor}%
\pgfsetfillcolor{textcolor}%
\pgftext[x=2.904493in,y=0.988340in,,top]{\color{textcolor}{\rmfamily\fontsize{6.000000}{7.200000}\selectfont\catcode`\^=\active\def^{\ifmmode\sp\else\^{}\fi}\catcode`\%=\active\def%{\%}150}}%
\end{pgfscope}%
\begin{pgfscope}%
\pgfpathrectangle{\pgfqpoint{2.375949in}{1.085562in}}{\pgfqpoint{2.114175in}{0.490004in}}%
\pgfusepath{clip}%
\pgfsetbuttcap%
\pgfsetroundjoin%
\pgfsetlinewidth{0.501875pt}%
\definecolor{currentstroke}{rgb}{0.690196,0.690196,0.690196}%
\pgfsetstrokecolor{currentstroke}%
\pgfsetdash{{0.500000pt}{0.825000pt}}{0.000000pt}%
\pgfpathmoveto{\pgfqpoint{3.433037in}{1.085562in}}%
\pgfpathlineto{\pgfqpoint{3.433037in}{1.575567in}}%
\pgfusepath{stroke}%
\end{pgfscope}%
\begin{pgfscope}%
\pgfsetbuttcap%
\pgfsetroundjoin%
\definecolor{currentfill}{rgb}{0.000000,0.000000,0.000000}%
\pgfsetfillcolor{currentfill}%
\pgfsetlinewidth{0.803000pt}%
\definecolor{currentstroke}{rgb}{0.000000,0.000000,0.000000}%
\pgfsetstrokecolor{currentstroke}%
\pgfsetdash{}{0pt}%
\pgfsys@defobject{currentmarker}{\pgfqpoint{0.000000in}{-0.048611in}}{\pgfqpoint{0.000000in}{0.000000in}}{%
\pgfpathmoveto{\pgfqpoint{0.000000in}{0.000000in}}%
\pgfpathlineto{\pgfqpoint{0.000000in}{-0.048611in}}%
\pgfusepath{stroke,fill}%
}%
\begin{pgfscope}%
\pgfsys@transformshift{3.433037in}{1.085562in}%
\pgfsys@useobject{currentmarker}{}%
\end{pgfscope}%
\end{pgfscope}%
\begin{pgfscope}%
\pgfsetbuttcap%
\pgfsetmiterjoin%
\definecolor{currentfill}{rgb}{1.000000,1.000000,1.000000}%
\pgfsetfillcolor{currentfill}%
\pgfsetfillopacity{0.700000}%
\pgfsetlinewidth{0.000000pt}%
\definecolor{currentstroke}{rgb}{0.000000,0.000000,0.000000}%
\pgfsetstrokecolor{currentstroke}%
\pgfsetstrokeopacity{0.700000}%
\pgfsetdash{}{0pt}%
\pgfpathmoveto{\pgfqpoint{3.348316in}{0.896673in}}%
\pgfpathlineto{\pgfqpoint{3.517758in}{0.896673in}}%
\pgfpathlineto{\pgfqpoint{3.517758in}{0.996673in}}%
\pgfpathlineto{\pgfqpoint{3.348316in}{0.996673in}}%
\pgfpathlineto{\pgfqpoint{3.348316in}{0.896673in}}%
\pgfpathclose%
\pgfusepath{fill}%
\end{pgfscope}%
\begin{pgfscope}%
\definecolor{textcolor}{rgb}{0.000000,0.000000,0.000000}%
\pgfsetstrokecolor{textcolor}%
\pgfsetfillcolor{textcolor}%
\pgftext[x=3.433037in,y=0.988340in,,top]{\color{textcolor}{\rmfamily\fontsize{6.000000}{7.200000}\selectfont\catcode`\^=\active\def^{\ifmmode\sp\else\^{}\fi}\catcode`\%=\active\def%{\%}160}}%
\end{pgfscope}%
\begin{pgfscope}%
\pgfpathrectangle{\pgfqpoint{2.375949in}{1.085562in}}{\pgfqpoint{2.114175in}{0.490004in}}%
\pgfusepath{clip}%
\pgfsetbuttcap%
\pgfsetroundjoin%
\pgfsetlinewidth{0.501875pt}%
\definecolor{currentstroke}{rgb}{0.690196,0.690196,0.690196}%
\pgfsetstrokecolor{currentstroke}%
\pgfsetdash{{0.500000pt}{0.825000pt}}{0.000000pt}%
\pgfpathmoveto{\pgfqpoint{3.961580in}{1.085562in}}%
\pgfpathlineto{\pgfqpoint{3.961580in}{1.575567in}}%
\pgfusepath{stroke}%
\end{pgfscope}%
\begin{pgfscope}%
\pgfsetbuttcap%
\pgfsetroundjoin%
\definecolor{currentfill}{rgb}{0.000000,0.000000,0.000000}%
\pgfsetfillcolor{currentfill}%
\pgfsetlinewidth{0.803000pt}%
\definecolor{currentstroke}{rgb}{0.000000,0.000000,0.000000}%
\pgfsetstrokecolor{currentstroke}%
\pgfsetdash{}{0pt}%
\pgfsys@defobject{currentmarker}{\pgfqpoint{0.000000in}{-0.048611in}}{\pgfqpoint{0.000000in}{0.000000in}}{%
\pgfpathmoveto{\pgfqpoint{0.000000in}{0.000000in}}%
\pgfpathlineto{\pgfqpoint{0.000000in}{-0.048611in}}%
\pgfusepath{stroke,fill}%
}%
\begin{pgfscope}%
\pgfsys@transformshift{3.961580in}{1.085562in}%
\pgfsys@useobject{currentmarker}{}%
\end{pgfscope}%
\end{pgfscope}%
\begin{pgfscope}%
\pgfsetbuttcap%
\pgfsetmiterjoin%
\definecolor{currentfill}{rgb}{1.000000,1.000000,1.000000}%
\pgfsetfillcolor{currentfill}%
\pgfsetfillopacity{0.700000}%
\pgfsetlinewidth{0.000000pt}%
\definecolor{currentstroke}{rgb}{0.000000,0.000000,0.000000}%
\pgfsetstrokecolor{currentstroke}%
\pgfsetstrokeopacity{0.700000}%
\pgfsetdash{}{0pt}%
\pgfpathmoveto{\pgfqpoint{3.876859in}{0.896673in}}%
\pgfpathlineto{\pgfqpoint{4.046301in}{0.896673in}}%
\pgfpathlineto{\pgfqpoint{4.046301in}{0.996673in}}%
\pgfpathlineto{\pgfqpoint{3.876859in}{0.996673in}}%
\pgfpathlineto{\pgfqpoint{3.876859in}{0.896673in}}%
\pgfpathclose%
\pgfusepath{fill}%
\end{pgfscope}%
\begin{pgfscope}%
\definecolor{textcolor}{rgb}{0.000000,0.000000,0.000000}%
\pgfsetstrokecolor{textcolor}%
\pgfsetfillcolor{textcolor}%
\pgftext[x=3.961580in,y=0.988340in,,top]{\color{textcolor}{\rmfamily\fontsize{6.000000}{7.200000}\selectfont\catcode`\^=\active\def^{\ifmmode\sp\else\^{}\fi}\catcode`\%=\active\def%{\%}170}}%
\end{pgfscope}%
\begin{pgfscope}%
\pgfpathrectangle{\pgfqpoint{2.375949in}{1.085562in}}{\pgfqpoint{2.114175in}{0.490004in}}%
\pgfusepath{clip}%
\pgfsetbuttcap%
\pgfsetroundjoin%
\pgfsetlinewidth{0.501875pt}%
\definecolor{currentstroke}{rgb}{0.690196,0.690196,0.690196}%
\pgfsetstrokecolor{currentstroke}%
\pgfsetdash{{0.500000pt}{0.825000pt}}{0.000000pt}%
\pgfpathmoveto{\pgfqpoint{4.490124in}{1.085562in}}%
\pgfpathlineto{\pgfqpoint{4.490124in}{1.575567in}}%
\pgfusepath{stroke}%
\end{pgfscope}%
\begin{pgfscope}%
\pgfsetbuttcap%
\pgfsetroundjoin%
\definecolor{currentfill}{rgb}{0.000000,0.000000,0.000000}%
\pgfsetfillcolor{currentfill}%
\pgfsetlinewidth{0.803000pt}%
\definecolor{currentstroke}{rgb}{0.000000,0.000000,0.000000}%
\pgfsetstrokecolor{currentstroke}%
\pgfsetdash{}{0pt}%
\pgfsys@defobject{currentmarker}{\pgfqpoint{0.000000in}{-0.048611in}}{\pgfqpoint{0.000000in}{0.000000in}}{%
\pgfpathmoveto{\pgfqpoint{0.000000in}{0.000000in}}%
\pgfpathlineto{\pgfqpoint{0.000000in}{-0.048611in}}%
\pgfusepath{stroke,fill}%
}%
\begin{pgfscope}%
\pgfsys@transformshift{4.490124in}{1.085562in}%
\pgfsys@useobject{currentmarker}{}%
\end{pgfscope}%
\end{pgfscope}%
\begin{pgfscope}%
\pgfsetbuttcap%
\pgfsetmiterjoin%
\definecolor{currentfill}{rgb}{1.000000,1.000000,1.000000}%
\pgfsetfillcolor{currentfill}%
\pgfsetfillopacity{0.700000}%
\pgfsetlinewidth{0.000000pt}%
\definecolor{currentstroke}{rgb}{0.000000,0.000000,0.000000}%
\pgfsetstrokecolor{currentstroke}%
\pgfsetstrokeopacity{0.700000}%
\pgfsetdash{}{0pt}%
\pgfpathmoveto{\pgfqpoint{4.405403in}{0.896673in}}%
\pgfpathlineto{\pgfqpoint{4.574845in}{0.896673in}}%
\pgfpathlineto{\pgfqpoint{4.574845in}{0.996673in}}%
\pgfpathlineto{\pgfqpoint{4.405403in}{0.996673in}}%
\pgfpathlineto{\pgfqpoint{4.405403in}{0.896673in}}%
\pgfpathclose%
\pgfusepath{fill}%
\end{pgfscope}%
\begin{pgfscope}%
\definecolor{textcolor}{rgb}{0.000000,0.000000,0.000000}%
\pgfsetstrokecolor{textcolor}%
\pgfsetfillcolor{textcolor}%
\pgftext[x=4.490124in,y=0.988340in,,top]{\color{textcolor}{\rmfamily\fontsize{6.000000}{7.200000}\selectfont\catcode`\^=\active\def^{\ifmmode\sp\else\^{}\fi}\catcode`\%=\active\def%{\%}180}}%
\end{pgfscope}%
\begin{pgfscope}%
\pgfpathrectangle{\pgfqpoint{2.375949in}{1.085562in}}{\pgfqpoint{2.114175in}{0.490004in}}%
\pgfusepath{clip}%
\pgfsetbuttcap%
\pgfsetroundjoin%
\pgfsetlinewidth{1.405250pt}%
\definecolor{currentstroke}{rgb}{0.121569,0.466667,0.705882}%
\pgfsetstrokecolor{currentstroke}%
\pgfsetdash{{5.180000pt}{2.240000pt}}{0.000000pt}%
\pgfpathmoveto{\pgfqpoint{2.375949in}{1.312497in}}%
\pgfpathlineto{\pgfqpoint{2.721229in}{1.296723in}}%
\pgfpathlineto{\pgfqpoint{3.066509in}{1.282234in}}%
\pgfpathlineto{\pgfqpoint{3.411789in}{1.268877in}}%
\pgfpathlineto{\pgfqpoint{3.757069in}{1.256522in}}%
\pgfpathlineto{\pgfqpoint{4.171404in}{1.242865in}}%
\pgfpathlineto{\pgfqpoint{4.495124in}{1.232987in}}%
\pgfpathlineto{\pgfqpoint{4.495124in}{1.232987in}}%
\pgfusepath{stroke}%
\end{pgfscope}%
\begin{pgfscope}%
\pgfpathrectangle{\pgfqpoint{2.375949in}{1.085562in}}{\pgfqpoint{2.114175in}{0.490004in}}%
\pgfusepath{clip}%
\pgfsetbuttcap%
\pgfsetroundjoin%
\pgfsetlinewidth{1.405250pt}%
\definecolor{currentstroke}{rgb}{1.000000,0.498039,0.054902}%
\pgfsetstrokecolor{currentstroke}%
\pgfsetdash{{5.180000pt}{2.240000pt}}{0.000000pt}%
\pgfpathmoveto{\pgfqpoint{2.375949in}{1.426331in}}%
\pgfpathlineto{\pgfqpoint{2.652173in}{1.409672in}}%
\pgfpathlineto{\pgfqpoint{2.928397in}{1.394124in}}%
\pgfpathlineto{\pgfqpoint{3.204621in}{1.379577in}}%
\pgfpathlineto{\pgfqpoint{3.549901in}{1.362660in}}%
\pgfpathlineto{\pgfqpoint{3.895180in}{1.346999in}}%
\pgfpathlineto{\pgfqpoint{4.240460in}{1.332460in}}%
\pgfpathlineto{\pgfqpoint{4.495124in}{1.322389in}}%
\pgfpathlineto{\pgfqpoint{4.495124in}{1.322389in}}%
\pgfusepath{stroke}%
\end{pgfscope}%
\begin{pgfscope}%
\pgfpathrectangle{\pgfqpoint{2.375949in}{1.085562in}}{\pgfqpoint{2.114175in}{0.490004in}}%
\pgfusepath{clip}%
\pgfsetbuttcap%
\pgfsetroundjoin%
\pgfsetlinewidth{1.405250pt}%
\definecolor{currentstroke}{rgb}{0.172549,0.627451,0.172549}%
\pgfsetstrokecolor{currentstroke}%
\pgfsetdash{{5.180000pt}{2.240000pt}}{0.000000pt}%
\pgfpathmoveto{\pgfqpoint{2.375949in}{1.411088in}}%
\pgfpathlineto{\pgfqpoint{2.721229in}{1.394589in}}%
\pgfpathlineto{\pgfqpoint{3.066509in}{1.379331in}}%
\pgfpathlineto{\pgfqpoint{3.411789in}{1.365173in}}%
\pgfpathlineto{\pgfqpoint{3.757069in}{1.351996in}}%
\pgfpathlineto{\pgfqpoint{4.171404in}{1.337335in}}%
\pgfpathlineto{\pgfqpoint{4.495124in}{1.326665in}}%
\pgfpathlineto{\pgfqpoint{4.495124in}{1.326665in}}%
\pgfusepath{stroke}%
\end{pgfscope}%
\begin{pgfscope}%
\pgfpathrectangle{\pgfqpoint{2.375949in}{1.085562in}}{\pgfqpoint{2.114175in}{0.490004in}}%
\pgfusepath{clip}%
\pgfsetrectcap%
\pgfsetroundjoin%
\pgfsetlinewidth{0.501875pt}%
\definecolor{currentstroke}{rgb}{0.000000,0.000000,0.000000}%
\pgfsetstrokecolor{currentstroke}%
\pgfsetdash{}{0pt}%
\pgfpathmoveto{\pgfqpoint{3.267835in}{1.085562in}}%
\pgfpathlineto{\pgfqpoint{3.267835in}{1.575567in}}%
\pgfusepath{stroke}%
\end{pgfscope}%
\begin{pgfscope}%
\pgfsetrectcap%
\pgfsetmiterjoin%
\pgfsetlinewidth{0.803000pt}%
\definecolor{currentstroke}{rgb}{0.000000,0.000000,0.000000}%
\pgfsetstrokecolor{currentstroke}%
\pgfsetdash{}{0pt}%
\pgfpathmoveto{\pgfqpoint{2.375949in}{1.085562in}}%
\pgfpathlineto{\pgfqpoint{2.375949in}{1.575567in}}%
\pgfusepath{stroke}%
\end{pgfscope}%
\begin{pgfscope}%
\pgfsetrectcap%
\pgfsetmiterjoin%
\pgfsetlinewidth{0.803000pt}%
\definecolor{currentstroke}{rgb}{0.000000,0.000000,0.000000}%
\pgfsetstrokecolor{currentstroke}%
\pgfsetdash{}{0pt}%
\pgfpathmoveto{\pgfqpoint{4.490124in}{1.085562in}}%
\pgfpathlineto{\pgfqpoint{4.490124in}{1.575567in}}%
\pgfusepath{stroke}%
\end{pgfscope}%
\begin{pgfscope}%
\pgfsetrectcap%
\pgfsetmiterjoin%
\pgfsetlinewidth{0.803000pt}%
\definecolor{currentstroke}{rgb}{0.000000,0.000000,0.000000}%
\pgfsetstrokecolor{currentstroke}%
\pgfsetdash{}{0pt}%
\pgfpathmoveto{\pgfqpoint{2.375949in}{1.085562in}}%
\pgfpathlineto{\pgfqpoint{4.490124in}{1.085562in}}%
\pgfusepath{stroke}%
\end{pgfscope}%
\begin{pgfscope}%
\pgfsetrectcap%
\pgfsetmiterjoin%
\pgfsetlinewidth{0.803000pt}%
\definecolor{currentstroke}{rgb}{0.000000,0.000000,0.000000}%
\pgfsetstrokecolor{currentstroke}%
\pgfsetdash{}{0pt}%
\pgfpathmoveto{\pgfqpoint{2.375949in}{1.575567in}}%
\pgfpathlineto{\pgfqpoint{4.490124in}{1.575567in}}%
\pgfusepath{stroke}%
\end{pgfscope}%
\begin{pgfscope}%
\pgfpathrectangle{\pgfqpoint{2.375949in}{1.085562in}}{\pgfqpoint{2.114175in}{0.490004in}}%
\pgfusepath{clip}%
\pgfsetbuttcap%
\pgfsetroundjoin%
\definecolor{currentfill}{rgb}{0.121569,0.466667,0.705882}%
\pgfsetfillcolor{currentfill}%
\pgfsetlinewidth{1.003750pt}%
\definecolor{currentstroke}{rgb}{0.121569,0.466667,0.705882}%
\pgfsetstrokecolor{currentstroke}%
\pgfsetdash{}{0pt}%
\pgfsys@defobject{currentmarker}{\pgfqpoint{-0.021960in}{-0.021960in}}{\pgfqpoint{0.021960in}{0.021960in}}{%
\pgfpathmoveto{\pgfqpoint{0.000000in}{-0.021960in}}%
\pgfpathcurveto{\pgfqpoint{0.005824in}{-0.021960in}}{\pgfqpoint{0.011410in}{-0.019646in}}{\pgfqpoint{0.015528in}{-0.015528in}}%
\pgfpathcurveto{\pgfqpoint{0.019646in}{-0.011410in}}{\pgfqpoint{0.021960in}{-0.005824in}}{\pgfqpoint{0.021960in}{0.000000in}}%
\pgfpathcurveto{\pgfqpoint{0.021960in}{0.005824in}}{\pgfqpoint{0.019646in}{0.011410in}}{\pgfqpoint{0.015528in}{0.015528in}}%
\pgfpathcurveto{\pgfqpoint{0.011410in}{0.019646in}}{\pgfqpoint{0.005824in}{0.021960in}}{\pgfqpoint{0.000000in}{0.021960in}}%
\pgfpathcurveto{\pgfqpoint{-0.005824in}{0.021960in}}{\pgfqpoint{-0.011410in}{0.019646in}}{\pgfqpoint{-0.015528in}{0.015528in}}%
\pgfpathcurveto{\pgfqpoint{-0.019646in}{0.011410in}}{\pgfqpoint{-0.021960in}{0.005824in}}{\pgfqpoint{-0.021960in}{0.000000in}}%
\pgfpathcurveto{\pgfqpoint{-0.021960in}{-0.005824in}}{\pgfqpoint{-0.019646in}{-0.011410in}}{\pgfqpoint{-0.015528in}{-0.015528in}}%
\pgfpathcurveto{\pgfqpoint{-0.011410in}{-0.019646in}}{\pgfqpoint{-0.005824in}{-0.021960in}}{\pgfqpoint{0.000000in}{-0.021960in}}%
\pgfpathlineto{\pgfqpoint{0.000000in}{-0.021960in}}%
\pgfpathclose%
\pgfusepath{stroke,fill}%
}%
\begin{pgfscope}%
\pgfsys@transformshift{-3.721483in}{2.794813in}%
\pgfsys@useobject{currentmarker}{}%
\end{pgfscope}%
\begin{pgfscope}%
\pgfsys@transformshift{0.930553in}{1.493782in}%
\pgfsys@useobject{currentmarker}{}%
\end{pgfscope}%
\begin{pgfscope}%
\pgfsys@transformshift{2.885570in}{1.336053in}%
\pgfsys@useobject{currentmarker}{}%
\end{pgfscope}%
\begin{pgfscope}%
\pgfsys@transformshift{3.769753in}{1.158421in}%
\pgfsys@useobject{currentmarker}{}%
\end{pgfscope}%
\begin{pgfscope}%
\pgfsys@transformshift{5.111860in}{1.159381in}%
\pgfsys@useobject{currentmarker}{}%
\end{pgfscope}%
\begin{pgfscope}%
\pgfsys@transformshift{6.604767in}{1.162049in}%
\pgfsys@useobject{currentmarker}{}%
\end{pgfscope}%
\begin{pgfscope}%
\pgfsys@transformshift{8.018892in}{1.206673in}%
\pgfsys@useobject{currentmarker}{}%
\end{pgfscope}%
\begin{pgfscope}%
\pgfsys@transformshift{10.157589in}{1.085966in}%
\pgfsys@useobject{currentmarker}{}%
\end{pgfscope}%
\begin{pgfscope}%
\pgfsys@transformshift{11.098294in}{1.129092in}%
\pgfsys@useobject{currentmarker}{}%
\end{pgfscope}%
\begin{pgfscope}%
\pgfsys@transformshift{12.519721in}{1.085527in}%
\pgfsys@useobject{currentmarker}{}%
\end{pgfscope}%
\end{pgfscope}%
\begin{pgfscope}%
\pgfpathrectangle{\pgfqpoint{2.375949in}{1.085562in}}{\pgfqpoint{2.114175in}{0.490004in}}%
\pgfusepath{clip}%
\pgfsetbuttcap%
\pgfsetroundjoin%
\definecolor{currentfill}{rgb}{1.000000,0.498039,0.054902}%
\pgfsetfillcolor{currentfill}%
\pgfsetlinewidth{1.003750pt}%
\definecolor{currentstroke}{rgb}{1.000000,0.498039,0.054902}%
\pgfsetstrokecolor{currentstroke}%
\pgfsetdash{}{0pt}%
\pgfsys@defobject{currentmarker}{\pgfqpoint{-0.021960in}{-0.021960in}}{\pgfqpoint{0.021960in}{0.021960in}}{%
\pgfpathmoveto{\pgfqpoint{-0.021960in}{-0.021960in}}%
\pgfpathlineto{\pgfqpoint{0.021960in}{-0.021960in}}%
\pgfpathlineto{\pgfqpoint{0.021960in}{0.021960in}}%
\pgfpathlineto{\pgfqpoint{-0.021960in}{0.021960in}}%
\pgfpathlineto{\pgfqpoint{-0.021960in}{-0.021960in}}%
\pgfpathclose%
\pgfusepath{stroke,fill}%
}%
\begin{pgfscope}%
\pgfsys@transformshift{-3.563173in}{3.153589in}%
\pgfsys@useobject{currentmarker}{}%
\end{pgfscope}%
\begin{pgfscope}%
\pgfsys@transformshift{2.031040in}{1.486598in}%
\pgfsys@useobject{currentmarker}{}%
\end{pgfscope}%
\begin{pgfscope}%
\pgfsys@transformshift{3.886517in}{1.281505in}%
\pgfsys@useobject{currentmarker}{}%
\end{pgfscope}%
\begin{pgfscope}%
\pgfsys@transformshift{5.781364in}{1.305643in}%
\pgfsys@useobject{currentmarker}{}%
\end{pgfscope}%
\end{pgfscope}%
\begin{pgfscope}%
\pgfpathrectangle{\pgfqpoint{2.375949in}{1.085562in}}{\pgfqpoint{2.114175in}{0.490004in}}%
\pgfusepath{clip}%
\pgfsetbuttcap%
\pgfsetroundjoin%
\definecolor{currentfill}{rgb}{0.172549,0.627451,0.172549}%
\pgfsetfillcolor{currentfill}%
\pgfsetlinewidth{1.003750pt}%
\definecolor{currentstroke}{rgb}{0.172549,0.627451,0.172549}%
\pgfsetstrokecolor{currentstroke}%
\pgfsetdash{}{0pt}%
\pgfsys@defobject{currentmarker}{\pgfqpoint{-0.021960in}{-0.021960in}}{\pgfqpoint{0.021960in}{0.021960in}}{%
\pgfpathmoveto{\pgfqpoint{0.000000in}{0.021960in}}%
\pgfpathlineto{\pgfqpoint{-0.021960in}{-0.021960in}}%
\pgfpathlineto{\pgfqpoint{0.021960in}{-0.021960in}}%
\pgfpathlineto{\pgfqpoint{0.000000in}{0.021960in}}%
\pgfpathclose%
\pgfusepath{stroke,fill}%
}%
\begin{pgfscope}%
\pgfsys@transformshift{-3.828554in}{2.845527in}%
\pgfsys@useobject{currentmarker}{}%
\end{pgfscope}%
\begin{pgfscope}%
\pgfsys@transformshift{0.261368in}{1.660264in}%
\pgfsys@useobject{currentmarker}{}%
\end{pgfscope}%
\begin{pgfscope}%
\pgfsys@transformshift{1.716852in}{1.381402in}%
\pgfsys@useobject{currentmarker}{}%
\end{pgfscope}%
\begin{pgfscope}%
\pgfsys@transformshift{3.216334in}{1.323485in}%
\pgfsys@useobject{currentmarker}{}%
\end{pgfscope}%
\begin{pgfscope}%
\pgfsys@transformshift{4.691329in}{1.143810in}%
\pgfsys@useobject{currentmarker}{}%
\end{pgfscope}%
\begin{pgfscope}%
\pgfsys@transformshift{5.764341in}{1.534224in}%
\pgfsys@useobject{currentmarker}{}%
\end{pgfscope}%
\begin{pgfscope}%
\pgfsys@transformshift{7.338091in}{1.245703in}%
\pgfsys@useobject{currentmarker}{}%
\end{pgfscope}%
\begin{pgfscope}%
\pgfsys@transformshift{8.793574in}{1.137237in}%
\pgfsys@useobject{currentmarker}{}%
\end{pgfscope}%
\begin{pgfscope}%
\pgfsys@transformshift{10.293057in}{1.240130in}%
\pgfsys@useobject{currentmarker}{}%
\end{pgfscope}%
\end{pgfscope}%
\begin{pgfscope}%
\pgfsetbuttcap%
\pgfsetmiterjoin%
\definecolor{currentfill}{rgb}{1.000000,1.000000,1.000000}%
\pgfsetfillcolor{currentfill}%
\pgfsetfillopacity{0.700000}%
\pgfsetlinewidth{0.000000pt}%
\definecolor{currentstroke}{rgb}{0.000000,0.000000,0.000000}%
\pgfsetstrokecolor{currentstroke}%
\pgfsetstrokeopacity{0.700000}%
\pgfsetdash{}{0pt}%
\pgfpathmoveto{\pgfqpoint{3.294262in}{1.466086in}}%
\pgfpathlineto{\pgfqpoint{4.154051in}{1.466086in}}%
\pgfpathquadraticcurveto{\pgfqpoint{4.154884in}{1.466086in}}{\pgfqpoint{4.154884in}{1.466919in}}%
\pgfpathlineto{\pgfqpoint{4.154884in}{1.551066in}}%
\pgfpathquadraticcurveto{\pgfqpoint{4.154884in}{1.551900in}}{\pgfqpoint{4.154051in}{1.551900in}}%
\pgfpathlineto{\pgfqpoint{3.294262in}{1.551900in}}%
\pgfpathquadraticcurveto{\pgfqpoint{3.293428in}{1.551900in}}{\pgfqpoint{3.293428in}{1.551066in}}%
\pgfpathlineto{\pgfqpoint{3.293428in}{1.466919in}}%
\pgfpathquadraticcurveto{\pgfqpoint{3.293428in}{1.466086in}}{\pgfqpoint{3.294262in}{1.466086in}}%
\pgfpathlineto{\pgfqpoint{3.294262in}{1.466086in}}%
\pgfpathclose%
\pgfusepath{fill}%
\end{pgfscope}%
\begin{pgfscope}%
\definecolor{textcolor}{rgb}{0.000000,0.000000,0.000000}%
\pgfsetstrokecolor{textcolor}%
\pgfsetfillcolor{textcolor}%
\pgftext[x=3.294262in,y=1.551066in,left,top]{\color{textcolor}{\rmfamily\fontsize{6.000000}{7.200000}\selectfont\catcode`\^=\active\def^{\ifmmode\sp\else\^{}\fi}\catcode`\%=\active\def%{\%}Ref.pkt. (SRT=53)}}%
\end{pgfscope}%
\begin{pgfscope}%
\pgfsetbuttcap%
\pgfsetmiterjoin%
\definecolor{currentfill}{rgb}{1.000000,1.000000,1.000000}%
\pgfsetfillcolor{currentfill}%
\pgfsetfillopacity{0.800000}%
\pgfsetlinewidth{1.003750pt}%
\definecolor{currentstroke}{rgb}{0.800000,0.800000,0.800000}%
\pgfsetstrokecolor{currentstroke}%
\pgfsetstrokeopacity{0.800000}%
\pgfsetdash{}{0pt}%
\pgfpathmoveto{\pgfqpoint{1.797521in}{1.633071in}}%
\pgfpathlineto{\pgfqpoint{4.537500in}{1.633071in}}%
\pgfpathquadraticcurveto{\pgfqpoint{4.554166in}{1.633071in}}{\pgfqpoint{4.554166in}{1.649738in}}%
\pgfpathlineto{\pgfqpoint{4.554166in}{2.088905in}}%
\pgfpathquadraticcurveto{\pgfqpoint{4.554166in}{2.105572in}}{\pgfqpoint{4.537500in}{2.105572in}}%
\pgfpathlineto{\pgfqpoint{1.797521in}{2.105572in}}%
\pgfpathquadraticcurveto{\pgfqpoint{1.780854in}{2.105572in}}{\pgfqpoint{1.780854in}{2.088905in}}%
\pgfpathlineto{\pgfqpoint{1.780854in}{1.649738in}}%
\pgfpathquadraticcurveto{\pgfqpoint{1.780854in}{1.633071in}}{\pgfqpoint{1.797521in}{1.633071in}}%
\pgfpathlineto{\pgfqpoint{1.797521in}{1.633071in}}%
\pgfpathclose%
\pgfusepath{stroke,fill}%
\end{pgfscope}%
\begin{pgfscope}%
\pgfsetbuttcap%
\pgfsetroundjoin%
\pgfsetlinewidth{1.003750pt}%
\definecolor{currentstroke}{rgb}{0.121569,0.466667,0.705882}%
\pgfsetstrokecolor{currentstroke}%
\pgfsetdash{{3.700000pt}{1.600000pt}}{0.000000pt}%
\pgfpathmoveto{\pgfqpoint{1.814187in}{2.014738in}}%
\pgfpathlineto{\pgfqpoint{1.897521in}{2.014738in}}%
\pgfpathlineto{\pgfqpoint{1.980854in}{2.014738in}}%
\pgfusepath{stroke}%
\end{pgfscope}%
\begin{pgfscope}%
\pgfsetbuttcap%
\pgfsetroundjoin%
\definecolor{currentfill}{rgb}{0.121569,0.466667,0.705882}%
\pgfsetfillcolor{currentfill}%
\pgfsetlinewidth{1.003750pt}%
\definecolor{currentstroke}{rgb}{0.121569,0.466667,0.705882}%
\pgfsetstrokecolor{currentstroke}%
\pgfsetdash{}{0pt}%
\pgfsys@defobject{currentmarker}{\pgfqpoint{-0.020833in}{-0.020833in}}{\pgfqpoint{0.020833in}{0.020833in}}{%
\pgfpathmoveto{\pgfqpoint{0.000000in}{-0.020833in}}%
\pgfpathcurveto{\pgfqpoint{0.005525in}{-0.020833in}}{\pgfqpoint{0.010825in}{-0.018638in}}{\pgfqpoint{0.014731in}{-0.014731in}}%
\pgfpathcurveto{\pgfqpoint{0.018638in}{-0.010825in}}{\pgfqpoint{0.020833in}{-0.005525in}}{\pgfqpoint{0.020833in}{0.000000in}}%
\pgfpathcurveto{\pgfqpoint{0.020833in}{0.005525in}}{\pgfqpoint{0.018638in}{0.010825in}}{\pgfqpoint{0.014731in}{0.014731in}}%
\pgfpathcurveto{\pgfqpoint{0.010825in}{0.018638in}}{\pgfqpoint{0.005525in}{0.020833in}}{\pgfqpoint{0.000000in}{0.020833in}}%
\pgfpathcurveto{\pgfqpoint{-0.005525in}{0.020833in}}{\pgfqpoint{-0.010825in}{0.018638in}}{\pgfqpoint{-0.014731in}{0.014731in}}%
\pgfpathcurveto{\pgfqpoint{-0.018638in}{0.010825in}}{\pgfqpoint{-0.020833in}{0.005525in}}{\pgfqpoint{-0.020833in}{0.000000in}}%
\pgfpathcurveto{\pgfqpoint{-0.020833in}{-0.005525in}}{\pgfqpoint{-0.018638in}{-0.010825in}}{\pgfqpoint{-0.014731in}{-0.014731in}}%
\pgfpathcurveto{\pgfqpoint{-0.010825in}{-0.018638in}}{\pgfqpoint{-0.005525in}{-0.020833in}}{\pgfqpoint{0.000000in}{-0.020833in}}%
\pgfpathlineto{\pgfqpoint{0.000000in}{-0.020833in}}%
\pgfpathclose%
\pgfusepath{stroke,fill}%
}%
\begin{pgfscope}%
\pgfsys@transformshift{1.897521in}{2.014738in}%
\pgfsys@useobject{currentmarker}{}%
\end{pgfscope}%
\end{pgfscope}%
\begin{pgfscope}%
\definecolor{textcolor}{rgb}{0.000000,0.000000,0.000000}%
\pgfsetstrokecolor{textcolor}%
\pgfsetfillcolor{textcolor}%
\pgftext[x=2.047521in,y=1.985572in,left,base]{\color{textcolor}{\rmfamily\fontsize{6.000000}{7.200000}\selectfont\catcode`\^=\active\def^{\ifmmode\sp\else\^{}\fi}\catcode`\%=\active\def%{\%}Zyklus \num{1}  ($a=$\,\num{104.412},\ $b=$\,\num{-0.905},\ $R^2=$\,\num{0.988},\ $n=$\,\num{10})}}%
\end{pgfscope}%
\begin{pgfscope}%
\pgfsetbuttcap%
\pgfsetroundjoin%
\pgfsetlinewidth{1.003750pt}%
\definecolor{currentstroke}{rgb}{1.000000,0.498039,0.054902}%
\pgfsetstrokecolor{currentstroke}%
\pgfsetdash{{3.700000pt}{1.600000pt}}{0.000000pt}%
\pgfpathmoveto{\pgfqpoint{1.814187in}{1.865571in}}%
\pgfpathlineto{\pgfqpoint{1.897521in}{1.865571in}}%
\pgfpathlineto{\pgfqpoint{1.980854in}{1.865571in}}%
\pgfusepath{stroke}%
\end{pgfscope}%
\begin{pgfscope}%
\pgfsetbuttcap%
\pgfsetmiterjoin%
\definecolor{currentfill}{rgb}{1.000000,0.498039,0.054902}%
\pgfsetfillcolor{currentfill}%
\pgfsetlinewidth{1.003750pt}%
\definecolor{currentstroke}{rgb}{1.000000,0.498039,0.054902}%
\pgfsetstrokecolor{currentstroke}%
\pgfsetdash{}{0pt}%
\pgfsys@defobject{currentmarker}{\pgfqpoint{-0.020833in}{-0.020833in}}{\pgfqpoint{0.020833in}{0.020833in}}{%
\pgfpathmoveto{\pgfqpoint{-0.020833in}{-0.020833in}}%
\pgfpathlineto{\pgfqpoint{0.020833in}{-0.020833in}}%
\pgfpathlineto{\pgfqpoint{0.020833in}{0.020833in}}%
\pgfpathlineto{\pgfqpoint{-0.020833in}{0.020833in}}%
\pgfpathlineto{\pgfqpoint{-0.020833in}{-0.020833in}}%
\pgfpathclose%
\pgfusepath{stroke,fill}%
}%
\begin{pgfscope}%
\pgfsys@transformshift{1.897521in}{1.865571in}%
\pgfsys@useobject{currentmarker}{}%
\end{pgfscope}%
\end{pgfscope}%
\begin{pgfscope}%
\definecolor{textcolor}{rgb}{0.000000,0.000000,0.000000}%
\pgfsetstrokecolor{textcolor}%
\pgfsetfillcolor{textcolor}%
\pgftext[x=2.047521in,y=1.836405in,left,base]{\color{textcolor}{\rmfamily\fontsize{6.000000}{7.200000}\selectfont\catcode`\^=\active\def^{\ifmmode\sp\else\^{}\fi}\catcode`\%=\active\def%{\%}Zyklus \num{2}  ($a=$\,\num{143.318},\ $b=$\,\num{-0.917},\ $R^2=$\,\num{0.997},\ $n=$\,\num{4})}}%
\end{pgfscope}%
\begin{pgfscope}%
\pgfsetbuttcap%
\pgfsetroundjoin%
\pgfsetlinewidth{1.003750pt}%
\definecolor{currentstroke}{rgb}{0.172549,0.627451,0.172549}%
\pgfsetstrokecolor{currentstroke}%
\pgfsetdash{{3.700000pt}{1.600000pt}}{0.000000pt}%
\pgfpathmoveto{\pgfqpoint{1.814187in}{1.716404in}}%
\pgfpathlineto{\pgfqpoint{1.897521in}{1.716404in}}%
\pgfpathlineto{\pgfqpoint{1.980854in}{1.716404in}}%
\pgfusepath{stroke}%
\end{pgfscope}%
\begin{pgfscope}%
\pgfsetbuttcap%
\pgfsetmiterjoin%
\definecolor{currentfill}{rgb}{0.172549,0.627451,0.172549}%
\pgfsetfillcolor{currentfill}%
\pgfsetlinewidth{1.003750pt}%
\definecolor{currentstroke}{rgb}{0.172549,0.627451,0.172549}%
\pgfsetstrokecolor{currentstroke}%
\pgfsetdash{}{0pt}%
\pgfsys@defobject{currentmarker}{\pgfqpoint{-0.020833in}{-0.020833in}}{\pgfqpoint{0.020833in}{0.020833in}}{%
\pgfpathmoveto{\pgfqpoint{0.000000in}{0.020833in}}%
\pgfpathlineto{\pgfqpoint{-0.020833in}{-0.020833in}}%
\pgfpathlineto{\pgfqpoint{0.020833in}{-0.020833in}}%
\pgfpathlineto{\pgfqpoint{0.000000in}{0.020833in}}%
\pgfpathclose%
\pgfusepath{stroke,fill}%
}%
\begin{pgfscope}%
\pgfsys@transformshift{1.897521in}{1.716404in}%
\pgfsys@useobject{currentmarker}{}%
\end{pgfscope}%
\end{pgfscope}%
\begin{pgfscope}%
\definecolor{textcolor}{rgb}{0.000000,0.000000,0.000000}%
\pgfsetstrokecolor{textcolor}%
\pgfsetfillcolor{textcolor}%
\pgftext[x=2.047521in,y=1.687238in,left,base]{\color{textcolor}{\rmfamily\fontsize{6.000000}{7.200000}\selectfont\catcode`\^=\active\def^{\ifmmode\sp\else\^{}\fi}\catcode`\%=\active\def%{\%}Zyklus \num{3}  ($a=$\,\num{61.760},\ $b=$\,\num{-0.753},\ $R^2=$\,\num{0.948},\ $n=$\,\num{9})}}%
\end{pgfscope}%
\end{pgfpicture}%
\makeatother%
\endgroup%

    \subcaption{Spezifischen Feinstaubemissionen je Aufrau-Zyklus über kumulierter Verschleißarbeit $\sum_i W_{\mathrm{wear},i}$ sowie zugehörige Power-Law-Fits~$\hat{y}_k(x)$.}
    \label{fig:drift_emission}
\end{figure}



% End of File
% !TeX root = ../main.tex
% 07-application.tex
\chapter{Anwendung}
\label{chap:application}

\lipsum[1]



% End of File
% !TeX root = ../main.tex
% 08-conclusion-outlook.tex
\chapter{Résumé und Ausblick}
\label{chap:conclusion_outlook}

\lipsum[1]



% End of File

% ========================
% BACKMATTER (z. B. Literatur, wieder einfacher Abstand)
% ========================
% \backmatter           % unnummerierte Kapitel

\printbibliography[heading=bibintoc, title={Literaturverzeichnis}]

% ========================
% ANHANG (Kapitel A, B, C, ... – meist wie Hauptteil gesetzt)
% ========================
\appendix             % Kapitelnummern werden zu A, B, C, ...

% !TeX root = ../main.tex
% app-b-additional-02-related.tex
\chapter{Ergänzende Dokumentation zum Stand der Technik}
\label{appchap:add_doc_related}

\lipsum[1]



% End of File
% !TeX root = ../main.tex
% app-c-additional-03-methodology.tex
\chapter{Ergänzende Dokumentation zur Methodik}
\label{appchap:add_doc_method}

\lipsum[1]



% End of File  % Ergänzende Dokumentation zur Methodik
% !TeX root = ../main.tex
% app-d-additional-04-validation.tex
\chapter{Ergänzende Dokumentation zur Validierung}
\label{appchap:add_doc_valid}

\lipsum[1]



% End of File  % Ergänzende Dokumentation zur Validierung
% !TeX root = ../main.tex
% app-e-additional-05-preparation.tex
\chapter{Ergänzende Dokumentation zur Entwicklung}
\label{appchap:add_doc_preparation}

\lipsum[1]



% End of File
% !TeX root = ../main.tex
% app-f-additional-06-evaluation.tex
\chapter[head={Ergänzung zu Versuchen \& Auswertung}]{Ergänzende Dokumentation zu den Prüfstandsversuchen und Auswertung}
\label{appchap:suppl_illustrations_evaluation}

\lipsum[1]



% End of File
% !TeX root = ../main.tex
% app-g-additional-07-application.tex
\chapter[head={Ergänzung zur Anwendung der Modelle}]{Ergänzende Dokumentation zur Anwendung der Modelle}
\label{appchap:add_doc_application}

\lipsum[1]



% End of File  % Ergänzende Dokumentation zur Anwendung
% !TeX root = ../main.tex
% app-h-projectplan.tex

\chapter{Vollständiger Projektplan}
\label{chap:pojectplan_full}

% \vspace(2cm)
\begin{figure}[htbp]
  \centering
  \rotatebox{90}{%
    \resizebox{0.7429\textheight}{!}{%
      \begin{ganttchart}[
        hgrid,
        vgrid,
        x unit=0.45cm, % Falls das Diagramm breiter/schmaler sein soll (z.B. 0.5cm oder 0.4cm)
        y unit title=0.7cm,
        y unit chart=0.45cm,
        title label font=\bfseries\footnotesize,
        bar label font=\footnotesize,
        group label font=\bfseries\footnotesize,
        milestone label font=\footnotesize\color{hka_blue},
        milestone/.append style={fill=hka_red},
        group/.append style={fill=hka_blue, draw=hka_blue}
      ]{1}{26}
                
        % Titelzeilen für die Wochen
        \gantttitle{Woche der Masterthesis}{26} \\
        \gantttitlelist{1,...,26}{1} \\
            
        % --- AP1: Rahmen & Literatur ---
        \ganttgroup{AP1: Literaturrecherche}{1}{5} \\
        \ganttmilestone{M1: Literaturreview abgeschlossen}{5} \\
            
        % --- AP2: CarMaker ---
        \ganttgroup{AP2: Einarbeitung Simulation}{2}{7} \\
        \ganttmilestone{M2: Schnittstellen verstanden}{7} \\
            
        % --- AP3: Realfahrdaten ---
        \ganttgroup{AP3: Validierung von \dots}{5}{9} \\
        \ganttmilestone{M3: Validierung abgeschlossen}{9} \\
            
        % --- AP4: Manövererkennung ---
        \ganttgroup{AP4: Entwicklung von \dots}{6}{13} \\
        \ganttmilestone{M4: Vorläufige \dots}{13} \\
            
        % --- AP5: Prüfstand ---
        \ganttgroup{AP5: Prüfstandsbegleitung}{8}{16} \\
        \ganttmilestone{M5: Prüfstandsversuche durchgeführt}{16} \\
            
        % --- AP6: Auswertung & Berechnung 1 ---
        \ganttgroup{AP6: Berechnungsvorschriften}{9}{16} \\
        \ganttmilestone{M6: Auswertungsmethodik \& -pipeline}{16} \\
            
        % --- AP7: Auswertung & Berechnung 2 ---
        \ganttgroup{AP7: Anwendung und Modelle}{15}{21} \\
        \ganttmilestone{M7: Anwendung}{21} \\
            
        % --- AP8: Thesis ---
        \ganttgroup{AP8: Schreiben / Feinschliff}{20}{23} \\
        \ganttmilestone{M8: Vollständiger Entwurf}{23} \\
            
        % --- AP9: Abschluss ---
        \ganttgroup{AP9: Puffer}{24}{24} \\
        \ganttmilestone{M9: Abgabe}{25}
            
      \end{ganttchart}
    }
  }
  \caption{Zeitlicher Projektplan der Masterthesis}
  \label{fig:gantt}
\end{figure}

% End of File

\end{document}